% concatenated from main.tex
\documentclass[aos]{imsart}

%% Packages
\RequirePackage{amsthm,amsmath,amsfonts,amssymb}
%\RequirePackage[numbers,sort&compress]{natbib}
\RequirePackage[authoryear]{natbib}%% uncomment this for author-year citations
\RequirePackage[colorlinks,citecolor=blue,urlcolor=blue]{hyperref}
\RequirePackage{graphicx}

\RequirePackage{cleveref}
\RequirePackage{enumitem}
\RequirePackage{fix-cm}
\RequirePackage{xcolor} % for 'teal'

\RequirePackage{algorithm}
\RequirePackage{algpseudocode}

\startlocaldefs
%%%%%%%%%%%%%%%%%%%%%%%%%%%%%%%%%%%%%%%%%%%%%%
%%                                          %%
%% Uncomment next line to change            %%
%% the type of equation numbering           %%
%%                                          %%
%%%%%%%%%%%%%%%%%%%%%%%%%%%%%%%%%%%%%%%%%%%%%%
%\numberwithin{equation}{section}
%%%%%%%%%%%%%%%%%%%%%%%%%%%%%%%%%%%%%%%%%%%%%%
%%                                          %%
%% For Axiom, Claim, Corollary, Hypothesis, %%
%% Lemma, Theorem, Proposition              %%
%% use \theoremstyle{plain}                 %%
%%                                          %%
%%%%%%%%%%%%%%%%%%%%%%%%%%%%%%%%%%%%%%%%%%%%%%
\theoremstyle{plain}
\newtheorem{axiom}{Axiom}
\newtheorem{claim}[axiom]{Claim}
\newtheorem{theorem}{Theorem}[section]
\newtheorem{lemma}[theorem]{Lemma}

\newtheorem{assumption}{Assumption}
%%%%%%%%%%%%%%%%%%%%%%%%%%%%%%%%%%%%%%%%%%%%%%
%%                                          %%
%% For Assumption, Definition, Example,     %%
%% Notation, Property, Remark, Fact         %%
%% use \theoremstyle{definition}            %%
%%                                          %%
%%%%%%%%%%%%%%%%%%%%%%%%%%%%%%%%%%%%%%%%%%%%%%
\theoremstyle{definition}
\newtheorem{definition}[theorem]{Definition}
\newtheorem*{example}{Example}
\newtheorem*{fact}{Fact}

\newtheorem{remark}{Remark}
%%%%%%%%%%%%%%%%%%%%%%%%%%%%%%%%%%%%%%%%%%%%%%
%% Please put your definitions here:        %%

\DeclareMathOperator*{\argmin}{arg\,min}
\DeclareMathOperator*{\argmax}{arg\,max}

% misc
\def\pmk{p_{\text{-}k}}
\def\rmk{r_{\text{-}k}}
\def\wmk{w_{\text{-}k}}
\def\Lambdamk{{\Lambda}_{\text{-}k}}
\def\Lambdamk{\Lambda_{\text{-}k}}
\def\wtLambdamk{\wt{\Lambda}_{\text{-}k}}
\def\whLambdamk{\wh{\Lambda}_{\text{-}k}}
\def\Aemk{A_{e,\text{-}k}}
\newcommand{\Ajmk}[1]{A_{#1,\text{-}k}}
\def\mat{\textnormal{\textsc{mat}}}
\def\map{\textnormal{\textsc{map}}}
\def\reshape{\textnormal{\textsc{reshape}}}
\def\tr{\textnormal{\text{tr}}}
\def\wt#1{\widetilde{#1}}
\def\wh#1{\widehat{#1}}

% Sets
\def\C{\mathbb{C}} 
\def\E{\mathbb{E}} 
\def\L{\mathbb{L}} 
\def\M{\mathbb{M}} 
\def\N{\mathbb{N}} 
\def\P{\mathbb{P}}
\def\R{\mathbb{R}} 
\def\RR{\mathbbm{R}}
\def\T{\mathbb{T}} 
\def\Z{\mathbb{Z}} 

% Mathcal
\def\cA{\mathcal{A}} 
\def\cB{\mathcal{B}} 
\def\cC{\mathcal{C}} 
\def\cD{\mathcal{D}} 
\def\cE{\mathcal{E}} 
\def\cF{\mathcal{F}}
\def\cG{\mathcal{G}} \def\cH{\boldsymbol{\mathcal{H}}} \def\cI{\mathcal{I}} 
\def\cII{\mathcal{II}} 
\def\cIII{\mathcal{III}} 
\def\cIV{\mathcal{IV}} 
\def\cJ{{\mathcal{J}}} 
\def\cL{\mathcal{L}}
\def\cM{\mathcal{M}} 
\def\cN{\mathcal{N}} 
\def\cO{\mathcal{O}} 
\def\co{\mathcal{o}} 
\def\cP{\mathcal{P}} 
\def\cQ{\mathcal{Q}} 
\def\cS{\mathcal{S}}
\def\cT{\mathcal{T}} 
\def\cU{\mathcal{U}} 
\def\cW{\mathcal{W}} 
\def\cX{\mathcal{X}} 
\def\cY{\mathcal{Y}} 
\def\cZ{\mathcal{Z}}

% Arrays
\def\one{\mathbf{1}} 
\def\zero{\mathbf{0}} 
\def\bi{\boldsymbol{i}} 
\def\bI{\mathbf{I}} 
\def\be{\mathbf{e}} 
\def\bw{\mathbf{w}}
\def\bh{\mathbf{h}}
\def\bu{\mathbf{u}} 
\def\ba{\mathbf{a}} 
\def\bx{\mathbf{x}} 
\def\by{\mathbf{y}}
\def\bA{\mathbf{A}}
\def\bB{\mathbf{B}}
\def\bD{\mathbf{D}}
\def\bX{\mathbf{X}} 
\def\bY{\mathbf{Y}} 
\def\bU{\mathbf{U}} 
\def\bF{\mathbf{F}} 
\def\bC{\mathbf{C}} 
\def\bG{\mathbf{G}} 
\def\bH{\mathbf{H}}
\def\bM{\mathbf{M}} 
\def\bR{\mathbf{R}} 
\def\bP{\mathbf{P}} 
\def\bQ{\mathbf{Q}} 
\def\bS{\mathbf{S}} 
\def\bV{\mathbf{V}} 
\def\bW{\boldsymbol{W}}
\def\bv{\boldsymbol{v}} 
\def\bg{\mathbf{g}} 
\def\bf{\boldsymbol{f}} 
\def\bJ{\mathbf{J}} 
\def\br{\mathbf{r}}

% Bold Greek
\def\balpha{\boldsymbol{\alpha}}
\def\bbeta{\boldsymbol{\beta}}
\def\bsigma{\boldsymbol{\sigma}}
\def\bmu{\boldsymbol{\mu}} \def\bLam{\boldsymbol{\Lambda}} \def\bGam{\boldsymbol{\Gamma}} \def\bSigma{\boldsymbol{\Sigma}} \def\bOmega{\boldsymbol{\Omega}}
\def\blam{\boldsymbol{\lambda}} \def\bxi{\boldsymbol{\xi}} \def\bPi{\boldsymbol{\Pi}} \def\bep{\boldsymbol{\epsilon}} \def\bvep{\boldsymbol{\varepsilon}}
\def\bchi{\boldsymbol{\chi}} \def\bPhi{\mathbf{\Phi}}

% Hats and Tildes
\def\hf{\hat{f}} 
\def\hbeta{\hat{\beta}} 
\def\htheta{\hat{\theta}} 
\def\halpha{\hat{\alpha}} 
\def\tbeta{\tilde{\beta}} 
\def\ty{\tilde{y}}
\def\hy{\hat{y}} 
\def\thh{\widehat{\theta}} 
\def\thj{\theta_j} 
\def\heta{\widehat{\eta}} 
\def\tA{\tilde{A}} 
\def\tB{\tilde{B}} 
\def\tD{\tilde{D}}
\def\tE{\tilde{\bE}} 
\def\tF{\tilde{\cF}} 
\def\tK{\tilde{K}} 
\def\tS{\tilde{\cS}} 
\def\tbF{\tilde{\bF}} 
\def\tbR{\tilde{\bR}} 
\def\tbC{\tilde{\bC}}
\def\hR{\widehat{\bR}} 
\def\hC{\widehat{\bC}} 
\def\hG{\widehat{\bG}} 
\def\hM{\widehat{\bM}} 
\def\hH{\widehat{\bH}}

% Misc
\def\d{\mathrm{d}} 
\def\dd{\mbox{d}} 
\def\th{\mathrm{th}} 
\def\trans{\intercal}
\def\op{\mathrm{op}} 
\def\F{{F}} 
\def\TV{\mathrm{TV}} 
\def\MSE{\mathrm{MSE}} 
\def\e{\text{e}}

% Operators
\def\cov{\mathbb{C}\mathrm{ov}} 
\def\var{\mathbb{V}\mathrm{ar}}
\def\df{\mathrm{df}} 
\def\dim{\mathrm{dim}} 
\def\sgn{\mathrm{Sgn}} 
\def\vec{\mathrm{vec}} 
\def\Vech{\mathrm{Vech}} 
\def\col{\mathrm{col}}
\def\row{\mathrm{row}} 
\def\nul{\mathrm{null}} 
\def\nuli{\mathrm{nullity}} 
\def\rank{\mathrm{rank}} 
\def\spa{\mathrm{span}}
\def\supp{\mathrm{supp}} 
\def\diag{\mathrm{diag}} 
\def\aff{\mathrm{aff}} 
\def\conv{\mathrm{conv}} 
\def\range{\mathrm{range}}

% Misc symbols
\def\go{\bg^{(1)}} 
\def\gt{\bg^{(2)}} 
\def\lc{\ell^{\circ}} 
\def\al{a_{\ell}} 
\def\bl{b_{\ell}} 
\def\alc{a_{\ell^{\circ}}} 
\def\blc{b_{\ell^{\circ}}}
\def\delc{\Delta^{\circ}} 
\def\gc{\mathbf{g}^{\circ}} 
\def\bVR{\bV_{R,r_1}} 
\def\bVC{\bV_{C,r_2}}

\newcommand{\ZC}[1]{\textcolor{teal}{ZC: #1}}
%%%%%%%%%%%%%%%%%%%%%%%%%%%%%%%%%%%%%%%%%%%%%%

\endlocaldefs

\begin{document}

\begin{frontmatter}
\title{Identification and estimation of multi-order tensor factor models}
%\title{A sample article title with some additional note\thanksref{t1}}
\runtitle{Multi-order tensor factor models}
%\thankstext{T1}{A sample additional note to the title.}

\begin{aug}
%%%%%%%%%%%%%%%%%%%%%%%%%%%%%%%%%%%%%%%%%%%%%%%
%% Only one address is permitted per author. %%
%% Only division, organization and e-mail is %%
%% included in the address.                  %%
%% Additional information can be included in %%
%% the Acknowledgments section if necessary. %%
%% ORCID can be inserted by command:         %%
%% \orcid{0000-0000-0000-0000}               %%
%%%%%%%%%%%%%%%%%%%%%%%%%%%%%%%%%%%%%%%%%%%%%%%
\author[A]{\fnms{Zetai}~\snm{Cen}
%\thanks{[\textbf{Corresponding author.}]}
\ead[label=e1]{zetai.cen@bristol.ac.uk}}\orcid{0009-0007-0215-8716} %,
%\author[B]{\fnms{Second}~\snm{Author}\ead[label=e2]{second@somewhere.com}\orcid{0000-0000-0000-0000}}
%\and
%\author[B]{\fnms{Third}~\snm{Author}\ead[label=e3]{third@somewhere.com}}
%%%%%%%%%%%%%%%%%%%%%%%%%%%%%%%%%%%%%%%%%%%%%%
%% Addresses                                %%
%%%%%%%%%%%%%%%%%%%%%%%%%%%%%%%%%%%%%%%%%%%%%%
\address[A]{School of Mathematics,
University of Bristol
%\textbf{[Additional affiliations should be put in the Acknowledgments section]}
\printead[presep={ ,\ }]{e1}}

%\address[B]{Department,
%University or Company Name \textbf{[Additional affiliations should be put in the Acknowledgments section]}\printead[presep={,\ }]{e2,e3}}
\end{aug}

\begin{abstract}
We propose a novel framework in high-dimensional factor models to simultaneously analyse multiple tensor time series, each with potentially different tensor orders and dimensionality. The connection between different tensor time series is through their global factors that are correlated to each other. A salient feature of our model is that when all tensor time series have the same order, it can be regarded as an extension of multilevel factor models from vectors to general tensors. Under very mild conditions, we separate the global and local components in the proposed model. Parameter estimation is thoroughly discussed, including a consistent factor number estimator. With strong correlation between global factors and noise allowed, we derive the rates of convergence of our estimators, which can be more superior than those of existing methods for multilevel factor models. We also develop estimators that are more computationally efficient, with rates of convergence spelt out. Extensive experiments are performed under various settings, corroborating with the pronounced theoretical results. As a real application example, we analyse a set of taxi data to study the traffic flow between Times Squares and its neighbouring areas.
\end{abstract}

\begin{keyword}[class=MSC]
\kwd[Primary ]{62H25}
\kwd{62H12}
%\kwd[; secondary ]{00X00}
\end{keyword}

\begin{keyword}
\kwd{Dimension reduction}
\kwd{global factors}
\kwd{multilevel factor models}
\kwd{tensor order}
\kwd{tensor time series}
\end{keyword}

\end{frontmatter}
%%%%%%%%%%%%%%%%%%%%%%%%%%%%%%%%%%%%%%%%%%%%%%
%% Please use \tableofcontents for articles %%
%% with 50 pages and more                   %%
%%%%%%%%%%%%%%%%%%%%%%%%%%%%%%%%%%%%%%%%%%%%%%
%\tableofcontents

\section{Introduction}
\label{sec: introduction}

With the rapid development of technology over the past few decades, high-dimensional time series are becoming more ubiquitous. This constantly demands methodologies to facilitate interpretation and forecasting. Among many other methods, factor models are found very efficient in dimension reduction and helping researchers to understand the dependence structure in data sets. Credited to these advantages of factor models, they have been applied to various fields, including but not limited to psychology \citep{McCraeJohn1992}, biology \citep{Hirzeletal2002, Hochreiteretal2006}, economics and finance \citep{ChamberlainRothschild1983, FamaFrench1993, StockWatson2002}.

More recently, researchers have open up to time series formed by matrices or, more generally, multidimensional arrays or tensors. Not only are such representations often the most natural by how the data is collected, they also provide useful insight that flattening the data might overlook. Thus, the literature has seen increasing work in matrix- or generally tensor-valued factor models. For matrix-valued time series, \cite{Wangetal2019} proposed matrix factor models by extending the framework in \cite{Lametal2011}. \cite{yu2022projected} also studied a matrix factor model and proposed a projection estimator with faster rate of convergence than the PCA-type (principal component analysis) estimator. \cite{ChenFan2023} proposed $\alpha$-PCA by generalizing the method for vector factor models by \cite{Bai2003}. Addressing the interplay between rows and columns, \cite{LamCen2024} proposed a main effect factor model, and \cite{Yuanetal2023} studied a two-way dynamic factor model. Methodologies are also ample for tensor factor models since the early work by \cite{Chenetal2022}, see \text{e.g.} \cite{barigozzi2022statistical}, \cite{Hanetal2022}, and \cite{Hanetal2024}.


Despite the vast amount of approaches in tensor factor modelling discussed above, all of them consider a single time series with tensor observation of the same order and dimensions.
Formally, for a single series of order-$K$ tensors $\{\cX_t\}$, a Tucker-type tensor factor model decomposes each $\cX_t$ as $\cX_t = \cF_t \times_{k=1}^K \bA_k + \cE_t$, where the latent core factor $\cF_t$ drives the dynamics of the observation through the factor loading matrices $\bA_k, \dots, \bA_K$.
This refrains us to conduct a comprehensive study on data sets containing multiple tensor time series, or to investigate the dependence structure among several tensor time series. In fact, this can be limited even when facing two vector time series that are closely connected in nature for panel data, yet of different dimensions. 

To tackle this problem, for vector time series, \cite{Wang2008} proposed a multilevel factor model, whose asymptotic theory is later generalized by \cite{Choietal2018}. We also note that such a framework is similar to hierarchical factor models in \cite{Moenchetal2013} for instance. In particular, their models consider $M$ series of observed vectors. For each group $m=1,\dots,M$, the vector observation $\bx_t^{(m)}\in \R^{p_m}$ is represented by $\bx_t^{(m)} = \bA_m \mathbf{g}_t + \bB_m \mathbf{f}_t^{(m)} +\mathbf{e}_t^{(m)}$. The low-rank vectors $\mathbf{g}_t$ and $\mathbf{f}_t^{(m)}$ are termed as the global and local factors respectively, with $\bA_m$ and $\bB_m$ being their corresponding loading matrices. By stacking all groups of vectors at the same timestamp, one may write
\[
\begin{pmatrix}
    \bx_t^{(1)} \\ \bx_t^{(2)} \\ \vdots \\ \bx_t^{(M)}
\end{pmatrix} = \begin{pmatrix}
    \bA_1 & \bB_1 & \zero & \ldots & \zero \\
    \bA_2 & \zero & \bB_2 & \ldots & \zero \\
    \vdots & \vdots & \vdots & \ddots & \vdots \\
    \bA_M & \zero & \zero & \ldots & \bB_M
\end{pmatrix} \begin{pmatrix}
    \mathbf{g}_t \\ \mathbf{f}_t^{(1)} \\ \vdots \\ \mathbf{f}_t^{(M)}
\end{pmatrix} + \begin{pmatrix}
    \mathbf{e}_t^{(1)} \\ \mathbf{e}_t^{(2)} \\ \vdots \\ \mathbf{e}_t^{(M)}
\end{pmatrix} ,
\]
which can be viewed as a vector factor model with specific structures. In addition to the interpretation of factor structures at each group, one benefit of the techniques developed for multilevel factor models is that, the model estimation can be more efficient by accounting for the multilevel or group structure, compared to estimating a very sparse factor loading matrix if all observed panels with unaligned dimensions are directly stacked together, as shown above. However, re-organising data sets becomes infeasible for higher-order tensor data (including matrix data which are order-2 tensors), urgently requiring appropriate methods to model the data sets.

Given $M$ groups in total, for each series of matrix observations with dimension $p_{m,1}\times p_{m,2}$, $m=1,\dots, M$, \cite{Zhangetal2025} considered multilevel factor models by decently extending the work in \cite{Wang2008} and \cite{Choietal2018} to matrix time series. However, their estimators of local factor loading matrices did not take advantage of the matrix format as \cite{yu2022projected} did, leaving the possibility to improve the estimation performance. Although their augmented iterative algorithm tried to amend the slow rates of convergence, both their initial estimator and iterative algorithm are computationally inefficient. Moreover, their Condition~6 required that all row dimensions $p_{m,1}$'s and column dimensions $p_{m,2}$'s are respectively of the same order, which might be unrealistic to many datasets empirically.

As tensor data becomes more prevalent nowadays, it is crucial to develop methodologies on multilevel factor modelling for tensor time series. Further ahead, facing multiple time series with different tensor orders gives rise to another challenge---the global factors cannot stay exactly the same over different series due to different shapes of arrays. A general solution is to consider global factors that are connected only due to their correlation. In fact, even when dealing with time series the same tensor orders, this setup is more flexible than that in multilevel factor models.

Another difficulty lies in specifying the number of factors, especially when tensor orders can be different. In the context of multilevel factor modelling, \cite{Choietal2018} proposed a consistent estimator based on several information criteria, whereas \cite{Zhangetal2025} only provided a practical approach based on eigenvalue ratio. The former method remains unclear when higher-order tensor time series are present, and for the latter, we will show in our numerical experiments that the vanilla eigenvalue-ratio estimator tend to overestimate the factor numbers, resulting in a part of idiosyncratic noise included in the estimated factors. Although this is not too problematic in practice, the effect of overestimating the factor numbers can be non-negligible \citep{BarigozziCho2020}.

To address all the above concern, we propose a multi-order tensor factor model to study a joint collection of tensor-valued time series, with at least the following contributions. As far as we are aware, our proposed model is the first work identifying and studying multiple tensor time series, each with potentially different orders and dimensions. Secondly, we adapt both techniques in \cite{Bai2003} and \cite{Lametal2011}, and establish theoretical rates of convergence that are often more superior than the results in, for example, \cite{Zhangetal2025}, even under very general dependence conditions. We also propose a consistent eigenvalue-ratio estimator for the global factor numbers, with a detailed discussion therein. The operator of tensor reshape introduced by \cite{CenLam2025_KronProd} is extended to facilitate natural representation of the global factors. Last but not least, a by-product of our model is a multilevel factor models for general tensor time series, with minor dimension requirements.



\subsection{Organisation of the paper}
After this Introduction, the rest of this paper is organised as follows. \Cref{sec: mm_tfm} introduces the data profile and formalises the multi-order tensor factor model. In detail, \Cref{subsec: model} concerns the identification issue, followed by parameter estimation in Sections~\ref{subsec: est_loading} and \ref{subsec: est_fac}. Assumptions are listed and explained in \Cref{subsec: assumption}, while theoretical guarantee of the estimated parameters is included in \Cref{subsec: theorem}. \Cref{sec: further_discussion} concerns various topics of interests. In detail, \Cref{subsec: explicit_global_fac} considers explicit forms of global factors, \Cref{subsec: est_num_fac} discusses the estimation of factor numbers, \Cref{subsec: bootstrap_implement} provides a computationally efficient implementation, and \Cref{subsec: assump_fac_satisfied} demonstrates alternative assumptions. \Cref{subsec: simulation} presents the numerical performance of the proposed methods by extensive Monte Carlo experiments, and finally, in \Cref{subsec: real_data}, a set of taxi traffic data is analysed using our model. All proof of theoretical results, together with additional details on taxi data, are relegated to the supplement.





\subsection{Notations}
%\textbf{Notations.} Some notational remarks.
Unless specified otherwise, we denote vectors, matrices and tensors by bold lower-case, bold capital, and calligraphic letters, i.e., $\bx$, $\bX$ and $\cX$, respectively. We also use $x_i, X_{ij}, \bX_{i\cdot}, \bX_{\cdot i}$ to denote, respectively, the $i$th element of a vector $\bx$, the $(i,j)$th element of $\bX$, the $i$th row vector (as a column vector) of $\bX$, and the $i$th column vector of $\bX$. We use $\otimes$ to represent the Kronecker product, and $\circ$ the Hadamard product. By convention, the total Kronecker product for an index set is computed in descending index. We use $a\asymp b$ to denote $a=O(b)$ and $b=O(a)$. Hereafter, given a positive integer $m$, define $[m]=\{1,2,\dots,m\}$. The $i$th largest eigenvalue and singular value of a matrix $\bX$ are denoted by $\lambda_i(\bX)$ and $\sigma_i(\bX)$, respectively. For an order-$K$ tensor $\cX = (X_{i_1,\ldots,i_K}) \in \mathbb{R}^{p_1\times \cdots\times p_K}$, we denote by $\mat_k(\cX) \in \R^{p_k\times \pmk}$ its \textit{mode-$k$ unfolding/matricisation}. We denote by $\cX \times_k \bA$ the \textit{mode-$k$ product} of a tensor $\cX$ with a matrix $A$, defined by $\mat_k(\cX) \times_k \bA = \bA\, \mat_k(\cX)$.


\section{Multi-order tensor factor models}\label{sec: mm_tfm}

\subsection{Model setup}\label{subsec: model}

Under some categorization $m\in[M]$ with fixed $M$, suppose that we observe a tensor time series $\{\cX_t^{(m)}\}$, $t\in[T]$, such that each $\cX_t^{(m)}$ is a mean-zero order-$K_m$ tensor of dimension $p_{m,1}\times \dots \times p_{m,K_m}$. Without loss of generality, let $K_1\leq \dots\leq K_M$. By joining $\{ \cX_t^{(1)}, \dots,\cX_t^{(M)} \}_{t\in[T]}$ as a collection (JXC hereafter), we assume JXC follows a \textit{multi-order tensor factor model} such that for each $m\in[M]$, the $m$th \textit{thread} $\{\cX_t^{(m)}\}$ admits the representation for each $t\in[T]$:
\begin{equation}
\label{eqn: model}
\cX_t^{(m)} = \cG_t^{(m)} \times_{k=1}^{K_m} \bA_{m,k} + \cF_t^{(m)} \times_{\ell=1}^{K_m} \bB_{m,\ell} + \cE_t^{(m)} ,
\end{equation}
where $\cG_t^{(m)}$ is the global core factor with dimension $r_{m,1}\times \dots\times r_{m,K_M}$, $\cF_t^{(m)}$ is the local core factor with dimension $u_{m,1}\times \dots\times u_{m,K_m}$, $\bA_{m,k}$ and $\bB_{m,\ell}$ ($k,\ell\in[K]$) are respectively the global and local factor loading matrices, and $\cE_t^{(m)}$ is the noise. For a non-trivial discussion, we assume $M>1$ throughout this paper, and otherwise \eqref{eqn: model} boils down to the classical Tucker-decomposition tensor factor model. Note that only JXC is observed and the right-hand side of \eqref{eqn: model} is latent.


Model \eqref{eqn: model} turns out to be a very general framework. A useful special case is when $K\equiv K_1=\dots =K_M$ and $\cG_t\equiv \cG_t^{(1)} =\dots=\cG_t^{(M)}$, which is commonly known as a multilevel/group factor model if $K=1$ \citep{Choietal2018, Huetal2025} or its matrix-valued extension for $K=2$ \citep{Zhangetal2025}. An immediate importance of \eqref{eqn: model} is thus a generalization of multilevel factor models to tensor time series. It might be worth to point out that even when the tensor orders are the same across $m\in[M]$, traditional factor models cannot directly characterize the factor structure in JXC, unless some forms of concatenation are carried out---which inevitably inflates the number of parameters and overlooks the structure along each mode.

Note that in most, if not all, multilevel factor models, the same global factor contributes to different groups, whereas the interplay among $\cG_t^{(m)}$ ($m\in[M]$) in our formulation \eqref{eqn: model} remains unclear. First, it should be straightforward that $\cG_t^{(m)} \neq \cG_t^{(n)}$ if $K_m\neq K_n$. More importantly, we choose not to specify, e.g., how $\cG_t^{(1)}$ affects $\cG_t^{(2)}$, and this lack of formulation is seen as an advantage in view of the generality of \eqref{eqn: model}. On the other hand, an explicit connection among the global core factors can better facilitate interpretation and forecasting on each thread $\cX_t^{(m)}$, and such a discussion is deferred to \Cref{subsec: explicit_global_fac} for interested readers.


We refer to $\cX_{G,t}^{(m)}:= \cG_t^{(m)} \times_{k=1}^{K_m} \bA_{m,k}$ and $\cX_{F,t}^{(m)}:= \cF_t^{(m)} \times_{\ell=1}^{K_m} \bB_{m,\ell}$ as the global and local components, respectively. Then with \eqref{eqn: model}, one can model multiple tensor time series with different orders at the same time, with the global component featuring the dependence of observed tensors across threads, which is new to the literature. To separate the global and local components, we impose the following identification conditions.


\begin{assumption}[Identification]\label{ass: identification}
For all $t\in[T]$, $m\in[M]$, we assume that: 
\begin{enumerate}[itemsep=0pt, label = (\alph*), left = 0pt]
    \item Any element in $\big\{ \cG_t^{(m)}, \cF_t^{(m)} \big\}$ has zero mean;
    \item Any element in $\big\{ \cG_t^{(m)}, \cF_t^{(m)} \big\}$ is uncorrelated to those in $\cF_t^{(n)}$ for all $n\neq m$;
    \item $\cG_t^{(m)}$ cannot be decomposed as $\cG_{t}^{(m)} =\cG_{t,1}^{(m)} + \cG_{t,2}^{(m)}$ with nonzero $\cG_{t,2}^{(m)}$ whose elements are uncorrelated to all elements in $\cG_t^{(n)}$ for all $n\neq m$.
\end{enumerate}
\end{assumption}



\Cref{ass: identification} is in the same spirit of Assumption~1 in \cite{Choietal2018} and Condition~6 in \cite{Zhangetal2025}. Part~(a) resolves the indeterminacy in mean. Parts~(b) and (c) together indicate that the local factor does not carry any dependence structure among the JXC, $\{\cX_t^{(1)}, \dots, \cX_t^{(M)} \}$. These assumptions are very mild in the sense that the identified global factor $\cG_t^{(m)}$ can be correlated to $\cF_t^{(m)}$. In comparison, such a correlation is prohibited in Assumption~1 in \cite{Choietal2018} for instance, due to their restricted $\cG_t^{(1)} =\dots= \cG_t^{(M)}$ which is required to be uncorrelated to the local factors. \Cref{ass: identification} allows us to conveniently identify the global and local components, and we present the result in the following theorem, with its proof relegated to the supplement.


\begin{theorem}
\label{prop: identification}
Let Assumption~\ref{ass: identification} hold. Then for any $t\in[T]$, the global and local components in \eqref{eqn: model} are identifiable in the sense that for any $m\in[M]$, two sets of parameters $(\cX_{G,t}^{(m)}, \cX_{F,t}^{(m)})$ and $(\wt\cX_{G,t}^{(m)}, \wt\cX_{F,t}^{(m)})$ satisfy $\cX_{G,t}^{(m)} + \cX_{F,t}^{(m)} = \wt\cX_{G,t}^{(m)} + \wt\cX_{F,t}^{(m)}$ if and only if $(\cX_{G,t}^{(m)}, \cX_{F,t}^{(m)}) =(\wt\cX_{G,t}^{(m)}, \wt\cX_{F,t}^{(m)})$.
\end{theorem}


\Cref{prop: identification} serves to separate the global and local components, which is sufficient for the discussion in this article, since both components admit Tucker decompositions and hence identification of parameters within has been well studied. For the same reason, the idiosyncratic noise is not involved in the proposition, as they can be separated from the global and local components by spiking eigenvalues in the unfolding covariance matrix.



\begin{remark}
Fix $m\in[M]$ and consider the dynamic of the $m$th thread only. \eqref{eqn: model} assumes a sparse tensor factor model in the sense that we can rewrite the model as
\begin{equation}
\label{eqn: Xt_rewrite_tfm}
\cX_t^{(m)} = \diag(\cG_t^{(m)}, \cF_t^{(m)}) \times_{k=1}^{K_m} (\bA_{m,k}, \bB_{m,k}) + \cE_t^{(m)} ,
\end{equation}
where $\diag(\cG_t^{(m)}, \cF_t^{(m)})$ denotes the block tensor consisting $\cG_t^{(m)}$ and $\cF_t^{(m)}$ on its diagonal. By modelling each observed tensor time series in the above way, more factors are introduced, which does not affect the theoretical results but undermines the estimation performance in practice. A more serious issue is that the dependence between threads is neglected and interpretation lost.
\end{remark}


\begin{remark}
With tensors with potentially different orders in JXC, a plausible method to model all tensor time series is to concatenate the vectorized data. In detail, note that for each $m\in[M]$, \eqref{eqn: model} implies $\vec(\cX_t^{(m)}) = (\otimes_{k=1}^{K_m} \bA_{m,k}, \otimes_{k=1}^{K_m} \bB_{m,k}) \big\{ \vec(\cG_t^{(m)})^\trans, \vec(\cF_t^{(m)})^\trans \big\}^\trans + \vec(\cE_t^{(m)})$. We thus have
{\small
\begin{align*}
    \begin{pmatrix}
    \vec(\cX_t^{(1)}) \\ \vec(\cX_t^{(2)}) \\ \vdots \\ \vec(\cX_t^{(M)})
    \end{pmatrix} 
    = \begin{pmatrix}
    \otimes_{k=1}^{K_1} \bA_{1,k} & \otimes_{k=1}^{K_1} \bB_{1,k}  & \ldots & \zero \\
    \zero & \zero &  \ldots & \zero \\
    \vdots & \vdots &  \ddots & \vdots \\
    \zero & \zero &  \ldots & \otimes_{k=1}^{K_M} \bB_{M,k}
    \end{pmatrix} \begin{pmatrix}
    \vec(\cG_t^{(1)}) \\ \vec(\cF_t^{(1)}) \\ \vdots \\ \vec(\cG_t^{(M)}) \\ \vec(\cF_t^{(M)})
    \end{pmatrix} + \begin{pmatrix}
    \vec(\cE_t^{(1)}) \\ \vec(\cE_t^{(2)}) \\ \vdots \\ \vec(\cE_t^{(M)})
    \end{pmatrix} ,
\end{align*}
}
or, if $\vec(\cG_t^{(1)}) =\dots =\vec(\cG_t^{(M)})$ under certain contexts (such as tensor multilevel factor models),
{\small
\begin{align*}
    \begin{pmatrix}
    \vec(\cX_t^{(1)}) \\ \vec(\cX_t^{(2)}) \\ \vdots \\ \vec(\cX_t^{(M)})
    \end{pmatrix} 
    = \begin{pmatrix}
    \otimes_{k=1}^{K_1} \bA_{1,k} & \otimes_{k=1}^{K_1} \bB_{1,k}  & \ldots & \zero \\
    \otimes_{k=1}^{K_2} \bA_{2,k} & \zero & \ldots & \zero \\
    \vdots & \vdots &  \ddots & \vdots \\
    \otimes_{k=1}^{K_M} \bA_{M,k} & \zero &  \ldots & \otimes_{k=1}^{K_M} \bB_{M,k}
    \end{pmatrix} \begin{pmatrix}
    \vec(\cG_t^{(1)}) \\ \vec(\cF_t^{(1)}) \\ \vdots \\ \vec(\cF_t^{(M)})
    \end{pmatrix} + \begin{pmatrix}
    \vec(\cE_t^{(1)}) \\ \vec(\cE_t^{(2)}) \\ \vdots \\ \vec(\cE_t^{(M)})
    \end{pmatrix} .
\end{align*}
}
However, both representations above significantly inflate the number of parameters in the loading matrix, and recovering the non-zero blocks within is also difficult in its own sake.
\end{remark}


\begin{remark}
\label{remark: absent_component}
In practice, we also allow for the absence of the global or local components (or both) in the representation \eqref{eqn: model} and by \Cref{ass: identification}, which is more realistic and reasonable from the perspective of data modelling. Although it can be attractive to make this mathematically rigorous in model \eqref{eqn: model} such as testing for the factor structure after estimating the components \citep[e.g.][for matrix time series]{Heetal2023}, it is beyond the scope of this paper and hence we refer to \Cref{subsec: est_num_fac} briefly discussing a practical solution and do not further pursue this line of analysis.
\end{remark}






\subsection{Estimation of the global and local loading matrices}
\label{subsec: est_loading}

In this subsection, we discuss the estimation of factor loading matrices in model \eqref{eqn: model}. First, for any $m\in[M]$, $k\in[K_m]$, define $\bA_{m,\text{-}k} := \otimes_{j\in[K_m] \setminus\{k\}} \bA_{m,j}$, and similarly $\bB_{m,\text{-}k}$. By convention, set $\bA_{m,\text{-}1}= \bB_{m,\text{-}1}= 1$ if $K_m=1$. Throughout this subsection, we treat the factor numbers $r_{m,k}$ and $u_{m,k}$ as given; the details of how to estimate them are deferred to \Cref{subsec: est_num_fac}.

To estimate the global factor loading matrices, consider the mode-$k$ unfolding of $\cX_t^{(m)}$ for each $t\in[T]$:
\begin{equation}
\label{eqn: unfold_Xt}
\mat_k(\cX_t^{(m)}) = \bA_{m,k} \mat_k(\cG_t^{(m)}) \bA_{m,\text{-}k}^\trans + \bB_{m,k} \mat_k(\cF_t^{(m)}) \bB_{m,\text{-}k}^\trans + \mat_k(\cE_t^{(m)}) .
\end{equation}
Then for any $n\neq m$, $\bi=(i_1,\ldots,i_{K_n}) \in [p_{n,1}] \times \cdots \times[p_{n,K_n}]$, we may define the filtered unfolding as
\begin{align*}
    \bOmega_{k,n,\bi}^{(m)} &:= \frac{1}{T} \sum_{t=1}^T \E\Big\{ \mat_k(\cX_t^{(m)}) \cdot (\cX_t^{(n)})_{\bi} \Big\} \\
    &= \bA_{m,k} \E\Bigg\{ \frac{1}{T} \sum_{t=1}^T \mat_k(\cG_t^{(m)}) \bA_{m,\text{-}k}^\trans \Big( \cG_t^{(n)} \times_{h=1}^{K_n} \bA_{n,h} \Big)_{\bi} \Bigg\} ,
\end{align*}
where the last equality used \Cref{ass: identification}(b) and \Cref{ass: dependence_moment}(b) (in \Cref{subsec: assumption}). Then enlightened by Equation~(4) in \cite{Lametal2011}, we introduce a non-negative definite matrix:
\begin{equation*}
\begin{split}
    \bSigma_{m,k} := \sum_{\substack{n=1\\ n\neq m}}^M \sum_{\bi \in [p_{n,1}] \times \cdots \times[p_{n,K_n}]} \bOmega_{k,n,\bi}^{(m)} \bOmega_{k,n,\bi}^{(m) \trans} ,
\end{split}
\end{equation*}
and observe that the $r_{m,k}$ leading eigenvectors of $\bSigma_{m,k}$ span the same linear space as $\bA_{m,k}$ does. Hence for any $m\in[M]$, $k\in[K_m]$, we may estimate $\bA_{m,k}$ by $\wh{\bA}_{m,k}$ which is defined as $\sqrt{p_{m,k}}$ times the $r_{m,k}$ leading eigenvectors of
\begin{equation}
\label{eqn: def_hat_Sigma_mk}
\begin{split}
    \wh\bSigma_{m,k} &:= \sum_{\substack{n=1\\ n\neq m}}^M \sum_{\bi \in [p_{n,1}] \times \cdots \times[p_{n,K_n}]} \wh\bOmega_{k,n,\bi}^{(m)} \wh\bOmega_{k,n,\bi}^{(m) \trans} ,
\end{split}
\end{equation}
where
\[
\wh\bOmega_{k,n,\bi}^{(m)} := \frac{1}{T} \sum_{t=1}^T \mat_k(\cX_t^{(m)}) \cdot (\cX_t^{(n)})_{\bi} .
\]


To estimate the local loading matrices, first consider $K_m> 1$. Let $\wh{\bA}_{m,k,\perp}$ be
%$\sqrt{p_{m,k}}$ times 
all but the $r_{m,k}$ leading eigenvectors of $\wh\bSigma_{m,k}$, \text{viz.} $\wh{\bA}_{m,k,\perp}$ is the orthogonal complement of $\wh{\bA}_{m,k}$. Then from \eqref{eqn: unfold_Xt},
\begin{align*}
    &\;\quad
    \mat_k(\cX_t^{(m)}) \wh{\bA}_{m,\text{-}k,\perp} - \bA_{m,k} \mat_k(\cG_t^{(m)}) \bA_{m,\text{-}k}^\trans \wh{\bA}_{m,\text{-}k,\perp} \\
    &= \bB_{m,k} \mat_k(\cF_t^{(m)}) \bB_{m,\text{-}k}^\trans \wh{\bA}_{m,\text{-}k,\perp} + \mat_k(\cE_t^{(m)}) \wh{\bA}_{m,\text{-}k,\perp} ,
\end{align*}
where $\wh{\bA}_{m,\text{-}k,\perp} := \otimes_{j\in[K_m] \setminus\{k\}} \wh{\bA}_{m,j,\perp}$. Intuitively, $\wh{\bA}_{m,\text{-}k,\perp}$ is close to the orthogonal complement of $\bA_{m,\text{-}k}$, and hence $\bA_{m,\text{-}k}^\trans \wh{\bA}_{m,\text{-}k,\perp}$ is close to zero. This motivates a PCA-type estimator for $\bB_{m,k}$, denoted by $\wh{\bB}_{m,k}$ which is defined as $\sqrt{p_{m,k}}$ times the $u_{m,k}$ leading eigenvectors of
\begin{equation}
\label{eqn: def_hat_Sigma_B_mk}
\begin{split}
    \wh\bSigma_{B,m,k} &:= \frac{1}{Tp_m} \sum_{t=1}^T \mat_k(\cX_t^{(m)}) \wh{\bA}_{m,\text{-}k,\perp} \wh{\bA}_{m,\text{-}k,\perp}^\trans \mat_k(\cX_t^{(m)})^\trans ,
\end{split}
\end{equation}
where $p_m:= \prod_{k=1}^{K_m} p_{m,k}$. Now suppose $K_m=1$ and the above approach to obtain $\wh{\bB}_{m,1}$ is thus infeasible. Fortunately, \eqref{eqn: unfold_Xt} can be rewritten as
\[
\mat_1(\cX_t^{(m)}) = (\bA_{m,1}, \bB_{m,1}) \big\{ \mat_1(\cG_t^{(m)})^\trans, \mat_1(\cF_t^{(m)})^\trans \big\}^\trans + \mat_k(\cE_t^{(m)}) .
\]
Hence by leveraging again the orthogonal complement of the global loading matrix, we may compute $\wh{\bB}_{m,1}$ as $\sqrt{p_{m,1}}$ times the $u_{m,k}$ leading eigenvectors of
\begin{equation}
\label{eqn: def_tilde_Sigma_B_m1}
\begin{split}
    \wt\bSigma_{B,m,1} &:= \frac{1}{Tp_{m,1}} \sum_{t=1}^T \wh{\bA}_{m,1,\perp} \wh{\bA}_{m,1,\perp}^\trans \mat_1(\cX_t^{(m)}) \mat_1(\cX_t^{(m)})^\trans \wh{\bA}_{m,1,\perp} \wh{\bA}_{m,1,\perp}^\trans .
\end{split}
\end{equation}

\iffalse
\[
\mat_k(\cX_t^{(m)}) = (\bA_{m,k}, \bB_{m,k}) \diag\big\{ \mat_k(\cG_t^{(m)}), \mat_k(\cF_t^{(m)}) \big\} (\bA_{m,\text{-}k}, \bB_{m,\text{-}k})^\trans + \mat_k(\cE_t^{(m)}).
\]
Hence let 
\begin{align*}
    &\;\quad
    \bA_{m,1,\perp} \bA_{m,1,\perp}^\trans (\bA_{m,1}, \bB_{m,1}) (\bA_{m,1}, \bB_{m,1})^\trans \bA_{m,1,\perp} \bA_{m,1,\perp}^\trans \\
    &=
    (0, \bA_{m,1,\perp} \bA_{m,1,\perp}^\trans \bB_{m,1}) (0, \bA_{m,1,\perp} \bA_{m,1,\perp}^\trans \bB_{m,1})^\trans
    = \bA_{m,1,\perp} \bA_{m,1,\perp}^\trans \bB_{m,1} \bB_{m,1}^\trans \bA_{m,1,\perp} \bA_{m,1,\perp}^\trans
\end{align*}
\fi







\subsection{Estimation of the core factors and components}
\label{subsec: est_fac}

With $\wh{\bA}_{m,k}$, $\wh{\bA}_{m,k,\perp}$, and $\wh{\bB}_{m,k}$ ($m\in[M], k\in[K_m]$) obtained in \Cref{subsec: est_loading}, we have for any $t\in[T]$, $m\in[M]$,
\begin{align*}
    &\;\quad
    (\otimes_{k=1}^K \wh{\bA}_{m,k,\perp})^\trans \vec(\cX_t^{m}) - (\otimes_{k=1}^K \wh{\bA}_{m,k,\perp})^\trans (\otimes_{k=1}^K \bA_{m,k}) \vec(\cG_t^{(m)}) \\
    &= (\otimes_{k=1}^K \wh{\bA}_{m,k,\perp})^\trans (\otimes_{k=1}^K \bB_{m,k}) \vec(\cF_t^{(m)}) + (\otimes_{k=1}^K \wh{\bA}_{m,k,\perp})^\trans \vec(\cE_t^{m}) ,
\end{align*}
where we expect $(\otimes_{k=1}^K \wh{\bA}_{m,k,\perp})^\trans (\otimes_{k=1}^K \bA_{m,k})$ to be close to zero, as discussed in \Cref{subsec: est_loading}. Then with the notation $\wh{\bC}_m := (\otimes_{k=1}^K \wh{\bA}_{m,k,\perp})^\trans (\otimes_{k=1}^K \wh{\bB}_{m,k})$, the vectorized local factor can be estimated in a least squared manner such that
\begin{equation}
\label{eqn: def_hat_Ft}
\begin{split}
    \vec(\wh{\cF}_t^{(m)}) &:= (\wh{\bC}_m^\trans \wh{\bC}_m)^{-1} \wh{\bC}_m^\trans (\otimes_{k=1}^K \wh{\bA}_{m,k,\perp})^\trans \vec(\cX_t^{m}) .
\end{split}
\end{equation}
For any $t\in[T]$, $m\in[M]$, the global factor can also be formulated as in a linear regression form:
\begin{align*}
    &\;\quad
    \vec(\cX_t^{m}) - (\otimes_{k=1}^K \bB_{m,k}) \vec(\cF_t^{(m)}) = (\otimes_{k=1}^K \bA_{m,k}) \vec(\cG_t^{(m)}) + \vec(\cE_t^{m}) .
\end{align*}
Then similar to \eqref{eqn: def_hat_Ft}, we may plug in the previous estimators and arrive at
\begin{equation}
\label{eqn: def_hat_Gt}
\begin{split}
    \vec(\wh{\cG}_t^{(m)}) &:= \Big\{ (\otimes_{k=1}^K \wh{\bA}_{m,k})^\trans (\otimes_{k=1}^K \wh{\bA}_{m,k}) \Big\}^{-1} \\
    &\;\quad
    \cdot (\otimes_{k=1}^K \wh{\bA}_{m,k})^\trans \Big\{ \vec(\cX_t^{m}) - (\otimes_{k=1}^K \wh{\bB}_{m,k}) \vec(\wh\cF_t^{(m)}) \Big\} .
\end{split}
\end{equation}

Finally, the global and local common components are respectively estimated by
\[
\wh\cX_{G,t}^{(m)}:= \wh\cG_t^{(m)} \times_{k=1}^{K_m} \wh{\bA}_{m,k}, \quad
\wh\cX_{F,t}^{(m)}:= \wh\cF_t^{(m)} \times_{\ell=1}^{K_m} \wh{\bB}_{m,\ell} .
\]



\begin{remark}
\label{remark: hat_Sigma_rewrite}
Suppose all tensor orders are the same, i.e., $K_1=\dots=K_M$. For any $k\in[K_1]$, define $p_{n,\text{-}k} =p_n/p_{n,k}$ and the matrix
\[
\wh\bOmega_{k,n,i,j}^{(m)} := \frac{1}{T} \sum_{t=1}^T \mat_k(\cX_t^{(m)})_{\cdot j} \mat_k(\cX_t^{(n)})_{\cdot i}^\trans .
\]
Then we may rewrite $\wh\bSigma_{m,k}$ defined in \eqref{eqn: def_hat_Sigma_mk} as
\begin{align*}
    \wh\bSigma_{m,k} &= \frac{1}{T^2} \sum_{\substack{n=1\\ n\neq m}}^M \sum_{t,s=1}^T \mat_k(\cX_t^{(m)}) \mat_k(\cX_s^{(m)})^\trans \sum_{\bi \in [p_{n,1}] \times \cdots \times[p_{n,K_n}]} (\cX_t^{(n)})_{\bi} (\cX_s^{(n)})_{\bi} \\
    % &= \frac{1}{T^2} \sum_{\substack{n=1\\ n\neq m}}^M \sum_{t=1}^T \sum_{s=1}^T \Big\{ \vec(\cX_t^{(n)})^\trans \otimes \mat_k(\cX_t^{(m)}) \Big\} \Big\{ \vec(\cX_s^{(n)}) \otimes \mat_k(\cX_s^{(m)})^\trans \Big\} \\
    &=
    \frac{1}{T^2} \sum_{\substack{n=1\\ n\neq m}}^M \sum_{t,s=1}^T \Big\{ \vec(\cX_s^{(n)})^\trans \vec(\cX_t^{(n)}) \Big\} \otimes \Big\{ \mat_k(\cX_t^{(m)}) \mat_k(\cX_s^{(m)})^\trans \Big\} \\
    % &= \frac{1}{T^2} \sum_{\substack{n=1\\ n\neq m}}^M \sum_{t,s=1}^T \tr\Big\{ \mat_k(\cX_s^{(n)}) \mat_k(\cX_t^{(n)})^\trans \Big\} \Big\{ \mat_k(\cX_t^{(m)}) \mat_k(\cX_s^{(m)})^\trans \Big\} \\
    &=
    \frac{1}{T^2} \sum_{\substack{n=1\\ n\neq m}}^M \sum_{t,s=1}^T \tr\Big\{ \sum_{i=1}^{p_{n,\text{-}k}} \mat_k(\cX_s^{(n)})_{\cdot i} \mat_k(\cX_t^{(n)})_{\cdot i}^\trans \Big\} \Big\{ \mat_k(\cX_t^{(m)}) \mat_k(\cX_s^{(m)})^\trans \Big\} \\
    &=
    \frac{1}{T^2} \sum_{\substack{n=1\\ n\neq m}}^M \sum_{i=1}^{p_{n,\text{-}k}} \sum_{t,s=1}^T \Big\{ \mat_k(\cX_t^{(n)})_{\cdot i}^\trans  \mat_k(\cX_s^{(n)})_{\cdot i} \Big\} \mat_k(\cX_t^{(m)}) \mat_k(\cX_s^{(m)})^\trans \\
    &=
    \frac{1}{T^2} \sum_{\substack{n=1\\ n\neq m}}^M \sum_{i,j=1}^{p_{n,\text{-}k}} \sum_{t,s=1}^T \mat_k(\cX_t^{(m)})_{\cdot j} \mat_k(\cX_t^{(n)})_{\cdot i}^\trans \mat_k(\cX_s^{(n)})_{\cdot i} \mat_k(\cX_s^{(m)})_{\cdot j}^\trans \\
    &=
    \sum_{\substack{n=1\\ n\neq m}}^M \sum_{i,j=1}^{p_{n,\text{-}k}} \wh\bOmega_{k,n,i,j}^{(m)} \wh\bOmega_{k,n,i,j}^{(m) \trans} ,
\end{align*}
which provides an equivalent way to compute $\wh\bSigma_{m,k}$. When $K_1=\dots=K_M=2$, the above $\wh\bOmega_{k,n,i,j}^{(m)}$ coincides with the notation used in \cite{Zhangetal2025} to estimate the global factor loading matrix. This indicates that under the special case of multilevel matrix factor models, our estimator $\wh\bA_{m,k}$ is the same as that by \cite{Zhangetal2025}.
\end{remark}




\section{Assumptions and Theoretical Results}\label{sec: theorem}

\subsection{Main assumptions}\label{subsec: assumption}

For each $m\in[M]$, define $p_{m,\text{-}k} =p_m/p_{m,k}$, $r_m= \prod_{k=1}^{K_m} r_{m,k}$, and $r_{m,\text{-}k} =r_m/r_{m,k}$. In the following, we present assumptions on the factor structure, noise, and dependence structure.


\begin{assumption}[Loading matrices]\label{ass: loadings}
For any $m\in[M]$, $k\in[K_m]$, there exists some constant $c$ fulfilling the follows.
%(i) there exists some constant $c$ such that $\|\bA_{m,k}\|_{\max}, \|\bB_{m,k}\|_{\max} \leq c <\infty$; and (ii) $\| p_{m,k}^{-1} \bA_{m,k}^\trans \bA_{m,k} -\bI_{r_{m,k}}\| \to 0$ and $\| p_{m,k}^{-1} \bB_{m,k}^\trans \bB_{m,k} -\bI_{r_{m,k}}\| \to 0$ as $\min\{T,p_{m,1},\dots, p_{m,K_m}\} \to \infty$.
\begin{enumerate}[itemsep=0pt, label = (\alph*), left = 0pt]
    \item $\|\bA_{m,k}\|_{\max}, \|\bB_{m,k}\|_{\max} \leq c <\infty$.
    \item $\big\| p_{m,k}^{-1} (\bA_{m,k}, \bB_{m,k})^\trans (\bA_{m,k}, \bB_{m,k}) -\bI_{r_{m,k}+u_{m,k}} \big\| \leq c p_{m,k}^{-1/2}$ for all $p_{m,k}\in \N$.
\end{enumerate}
\end{assumption}



\begin{assumption}[Core factor]\label{ass: core_factor}
For all $m\in[M]$, we assume that:
\begin{enumerate}[itemsep=0pt, label = (\alph*), left = 0pt]
    \item Any element in $\big\{ \cG_t^{(m)}, \cF_t^{(m)} \big\}$ has bounded fourth moment.
    \item For any $k\in[K_m]$, as $T\to\infty$, we have
    \[
    \frac{1}{T} \sum_{t=1}^T \mat_k(\cF_t^{(m)}) \mat_k(\cF_t^{(m)})^\trans \xrightarrow{P} \bSigma_{F,m,k} ,
    \]
    %and $T^{-1} \sum_{t=1}^T \mat_k(\cG_t^{(m)}) \mat_k(\cG_t^{(m)})^\trans \xrightarrow{P} \bSigma_{G,m,k}$, where $\bSigma_{F,m,k}$ and $\bSigma_{G,m,k}$ are positive definite matrices with finite eigenvalues.
    where $\bSigma_{F,m,k}$ is a positive definite matrix with finite eigenvalues.
    \item For any $n\in [M]$, $k\in[K_m]$, any set of matrices $\{ \bA_1, \dots, \bA_{K_n}\}$, where each $\bA_i\in \R^{p_{n,i} \times r_{n,i}}$ has full column rank and all eigenvalues bounded away from zero and infinity, there exists some positive constant $c$ such that as $T\to\infty$, for any $j\in[r_{m,k}]$:
    \[
    \sigma_{j}\bigg( \E\Big[ \mat_k(\cG_t^{(m)}) \otimes \Big\{ \vec(\cG_t^{(n)})^\trans \big( \otimes_{h=1}^{K_n} \bA_{h} \big)^\trans \Big\} \Big] \bigg) \asymp c .
    \]
\end{enumerate}
\end{assumption}



\begin{assumption}[Idiosyncratic noise]\label{ass: noise}
For all $m\in[M]$, $t\in[T]$, we assume that:
\begin{enumerate}[itemsep=0pt, label = (\alph*), left = 0pt]
    \item Any element in $\cE_t^{(m)}$ has zero mean and bounded fourth moment.
    \item For all $k\in[K_m]$, $t\in[T]$, $i\in[p_{m,k}]$, $j\in[p_{m,\text{-}k}]$, there exists some constant $c$ such that
    \[
    \sum_{s=1}^T \sum_{l=1}^{p_{m,k}} \sum_{h=1}^{p_{m,\text{-}k}} \Big| \E\Big\{ \mat_k(\cE_t^{(m)})_{ij} \mat_k(\cE_s^{(m)})_{lh} \Big\} \Big| \leq c .
    \]
    \item For all $k\in[K_m]$, $i,l_1\in[p_{m,k}]$, $j,h_1\in[p_{m,\text{-}k}]$, there exists some constant $c$ such that
    {\small
    \begin{align*}
        & \sum_{s=1}^T \sum_{l_2=1}^{p_{m,k}} \sum_{h_2=1}^{p_{m,\text{-}k}} \Big| \cov\Big\{ \mat_k(\cE_t^{(m)})_{ij} \mat_k(\cE_t^{(m)})_{l_1 j} , \mat_k(\cE_s^{(m)})_{ih_2} \mat_k(\cE_s^{(m)})_{l_2 h_2} \Big\} \Big| \leq c, \\
        & \sum_{s=1}^T \sum_{l_2=1}^{p_{m,k}} \sum_{h_2=1}^{p_{m,\text{-}k}} \Big| \cov\Big\{ \mat_k(\cE_t^{(m)})_{ij} \mat_k(\cE_t^{(m)})_{i h_1} , \mat_k(\cE_s^{(m)})_{l_2 j} \mat_k(\cE_s^{(m)})_{l_2 h_2} \Big\} \Big| \leq c, \\
        & \sum_{s=1}^T \sum_{l_2=1}^{p_{m,k}} \sum_{h_2=1}^{p_{m,\text{-}k}} \Big| \cov\Big\{ \mat_k(\cE_t^{(m)})_{ij} \mat_k(\cE_t^{(m)})_{l_1 h_1} , \mat_k(\cE_s^{(m)})_{ij} \mat_k(\cE_s^{(m)})_{l_2 h_2} \Big\} \Big| \leq c, \\
        & \sum_{s=1}^T \sum_{l_2=1}^{p_{m,k}} \sum_{h_2=1}^{p_{m,\text{-}k}} \Big| \cov\Big\{ \mat_k(\cE_t^{(m)})_{l_1 j} \mat_k(\cE_t^{(m)})_{i h_1} , \mat_k(\cE_s^{(m)})_{l_2 j} \mat_k(\cE_s^{(m)})_{i h_2} \Big\} \Big| \leq c .
    \end{align*}
    }
\end{enumerate}
\end{assumption}



\begin{assumption}[Dependence and moments across threads]\label{ass: dependence_moment}
For all $m\in[M]$: 
\begin{enumerate}[itemsep=0pt, label = (\alph*), left = 0pt]
    \item Any element in $\cX_t^{(m)}$ is uncorrelated to those in $\cE_t^{(n)}$ for all $t\in[T]$, $n\neq m$.
    \item For all $n\neq m$, there exists some constant $c<\infty$ such that
    \begin{align*}
    & \Bigg\| \var\Bigg\{ \frac{1}{\sqrt{T}} \sum_{t=1}^T \vec(\cX_t^{(m)}) \otimes \vec(\cE_t^{(n)}) \Bigg\} \Bigg\|_{\max} \leq c , \\
    & \Bigg\| \var\Bigg\{ \frac{1}{\sqrt{T}} \sum_{t=1}^T \vec(\cG_t^{(m)}) \otimes \vec(\cF_t^{(n)}) \Bigg\} \Bigg\|_{\max} \leq c , \\
    & \Bigg\| \var\Bigg\{ \frac{1}{\sqrt{T}} \sum_{t=1}^T \vec(\cF_t^{(m)}) \otimes \vec(\cF_t^{(n)}) \Bigg\} \Bigg\|_{\max} \leq c , \\
    & \Bigg\| \var\Bigg\{ \frac{1}{\sqrt{T}} \sum_{t=1}^T \vec(\cG_t^{(m)}) \otimes \vec(\cG_t^{(n)}) \Bigg\} \Bigg\|_{\max} \leq c .
    \end{align*}
\end{enumerate}
\end{assumption}


\begin{assumption}[Dependence and moments within threads]\label{ass: dependence_thread}
For all $m\in[M]$, $k\in[K_m]$, we assume that for any deterministic vectors $\bv\in\R^{p_{m,k}}$, $\bw\in\R^{p_{m,\text{-}k}}$, $\bg\in\R^{r_{m,k}}$, and $\bh\in\R^{r_{m,\text{-}k}}$ with $\|\bv\|, \|\bw\|, \|\bg\|, \|\bh\| =1$, there exists some constant $c$ such that: 
\begin{enumerate}[itemsep=0pt, label = (\alph*), left = 0pt]
    \item $\E\Big\| T^{-1/2} \sum_{t=1}^T \mat_k(\cF_t^{(m)}) \big\{ \bv^\trans \mat_k(\cE_t^{(m)}) \bw \big\} \Big\|_F^2 \leq c$.
    \item $\E\Big\| T^{-1/2} \sum_{t=1}^T \mat_k(\cF_t^{(m)}) \big\{ \bg^\trans \mat_k(\cG_t^{(m)}) \bh \big\} \Big\|_F^2 \leq c$.
    \item $\E\Big\| T^{-1} \sum_{t=1}^T \mat_k(\cG_t^{(m)}) \big\{ \bv^\trans \mat_k(\cE_t^{(m)}) \bw \big\} \Big\|_F^2 \leq c$.
\end{enumerate}

\iffalse % use the above neater presentation
\begin{enumerate}[itemsep=0pt, label = (\alph*), left = 0pt]
    \item 
    \[
    \E\Bigg\| \frac{1}{\sqrt{T}} \sum_{t=1}^T \mat_k(\cF_t^{(m)}) \Big\{ \bv^\trans \mat_k(\cE_t^{(m)}) \bw \Big\}\Bigg\|_F^2 \leq c .
    \]
    \item 
    \[
    \E\Bigg\| \frac{1}{\sqrt{T}} \sum_{t=1}^T \mat_k(\cF_t^{(m)}) \Big\{ \bg^\trans \mat_k(\cG_t^{(m)}) \bh \Big\}\Bigg\|_F^2 \leq c .
    \]
    \item 
    \[
    \E\Bigg\| \frac{1}{T} \sum_{t=1}^T \mat_k(\cG_t^{(m)}) \Big\{ \bv^\trans \mat_k(\cE_t^{(m)}) \bw \Big\}\Bigg\|_F^2 \leq c .
    \]
\end{enumerate}
\fi
\end{assumption}




\Cref{ass: loadings} is standard to describing the factor loading matrices. Part~(b) can be seen as an identification condition. It also indicates that the magnitude of each column of the loading matrices grows proportionally to the dimension $p_{m,k}$---hence factors being pervasive---at the rate $p_k^{-1/2}$, which is also seen in Assumption~2(ii) in \cite{barigozzi2024moving}. The asymptotic orthogonality between $\bA_{m,k}$ and $\bB_{m,k}$ is arguably very mild, compared to the exact orthogonality between global and local factors in Assumption~A1 in \cite{Huetal2025}. Here, our rate of convergence together with the asymptotic orthogonality actually allow the local loading estimator to have faster rate of convergence than the global loading estimator; see \Cref{thm: local_loading}. It is also worth to remark that weak factors as in \cite{Wangetal2019} (\text{cf.} Condition~4 within) are possible, and we decide to consider strong factors only for improving the reading experience.


Assumption~\ref{ass: core_factor}(a) constrains the moments of core factors, while part~(b) describes the local factors and is very common in the factor model literature. Part~(c) complements \Cref{ass: identification} by pinpointing the behaviour of global factors, which is not particularly strong; see the discussion in \Cref{subsec: assump_fac_satisfied} on how part~(c) can be satisfied.
We also note that Assumption~\ref{ass: core_factor}(c) is not the weakest possible in terms of dependency in the global factors, since it requires $\cG_t^{(m)}$ to be correlated to all $\cG_t^{(n)}$ for $n\neq m$. It is possible to allow $\cG_t^{(m)}$ to be independent of global factors in some threads, as long as there is at least one $\cG_t^{(n)}$ with $n\neq m$ that it is correlated to. However, This relaxation could come at the cost of very restricted dimensionality assumptions, and hence is not pursued here.
% Basically need $p_1=\dots=p_M$, which is fine when all tensor orders are the same, but becomes very strong otherwise.

\Cref{ass: noise} includes a set of regularity conditions for the idiosyncratic noise. Parts~(b) and (c) are essentially the same as Assumption~3 in \cite{barigozzi2022statistical} and Assumption~D in \cite{yu2022projected} for matrix factor models. We point out that requiring elements to have bounded fourth moment in Assumption~\ref{ass: noise} is weaker than many literature restricting eighth moment to be bounded \citep[e.g.][]{Choietal2018, ChenFan2023, Huetal2025}.


Lastly, Assumptions~\ref{ass: dependence_moment} and \ref{ass: dependence_thread} depict general moment conditions between and within threads, respectively. \Cref{ass: dependence_moment}(a) strengthens \Cref{ass: identification} such that the correlation among different threads is all pooled inside global factors, whereas \Cref{ass: dependence_moment}(b) is comparable to, for example, Assumption~4 in \cite{barigozzi2022statistical}. \Cref{ass: dependence_thread} significantly relax the dependence structure in similar model framework such as multilevel factor models. Overall, it implies the global component, local component, and noise can be pairwise correlated. A more salient feature according to part~(c) is that the correlation between global factor and noise can be quite strong, while still ensuring satisfactory theoretical results.









\subsection{Theoretical results}\label{subsec: theorem}

To appreciate the performance of parameter estimators, we first show the rate of convergence for the global loading matrix estimator as follows.


\begin{theorem}[Asymptotic consistency of global loading estimators]\label{thm: global_loading}
Under Assumptions~\ref{ass: identification}--\ref{ass: dependence_moment}, it holds for any $m\in[M]$, $k\in[K_m]$ that
\begin{align*}
    \frac{1}{p_{m,k}} \Big\| \wh{\bA}_{m,k} -\bA_{m,k} \Big\|_F^2 = \cO_P\Big( \frac{1}{T} \Big) .
\end{align*}
\end{theorem}

\Cref{thm: global_loading} is an important step for the results hereafter. The approach to obtain $\wh{\bA}_{m,k}$ leverages the condition that the global factor is uncorrelated to local factors and noise in different threads, which is akin to the approach in \cite{Lametal2011} and \cite{Wangetal2019} assuming the noise is white. Hence it is by no accident that our rate of $1/T$ coincides with \text{e.g.} Theorem~1 in \cite{Wangetal2019} under strong factors and the same normalization.

Note that the rate of convergence in \Cref{thm: global_loading} does not involve any cross-sectional dimension. First, in a vector factor model ($K_m=1$), such a $1/T$ rate is at least as good as the rate of convergence for the PCA-type estimator in \cite{BaiNg2002}, or even more superior when $p_{m,1}$ grows in a slower rate than $T$. For higher-order threads, i.e., $K_m>1$, it can be intuited by the difficulty in relating tensor times series with different orders, and the global factor is only identified through the correlation among threads.

In multilevel matrix factor models, \cite{Zhangetal2025} estimate the global loading matrix in a similar way and essentially the same result as our \Cref{thm: global_loading}. However, the same approach is also used in estimating the local loading matrix, resulting in the $1/T$ rate again. This overlooks the opportunity to improve the result by involving information from other modes (rows and columns in a matrix-valued data). Our PCA-type loading estimator based on \eqref{eqn: def_hat_Sigma_B_mk} resolves this. Some results are now in order, providing a theoretical guarantee for the remaining parameter estimators.




\begin{theorem}[Asymptotic consistency of local loading estimators]\label{thm: local_loading}
Let Assumptions~\ref{ass: identification}--\ref{ass: dependence_thread} hold. The following claims are true.
\begin{enumerate}[itemsep=0pt, label = (\alph*), left = 0pt]
    \item For any $m\in[M]$, $k\in[K_m]$ with $K_m> 1$, let $\wh{\bD}_{m,k}$ be the diagonal matrix consisting of the first $u_{m,k}$ eigenvalues of $\wh\bSigma_{B,m,k}$ defined in \eqref{eqn: def_hat_Sigma_B_mk}. Define
    \begin{align*}
        \wh{\bH}_{m,k} &= \frac{1}{Tp_m} \sum_{t=1}^T \mat_k(\cF_t^{(m)}) \bB_{m,\text{-}k}^\trans \wh{\bA}_{m,\text{-}k,\perp} \\
        &\;\quad
        \cdot \wh{\bA}_{m,\text{-}k,\perp}^\trans \bB_{m,\text{-}k} \mat_k(\cF_t^{(m)})^\trans \bB_{m,k}^\trans \wh{\bB}_{m,k} \wh{\bD}_{B,m,k}^{-1} .
    \end{align*}
    Then $\wh{\bH}_{m,k}$ is asymptotically invertible such that $\wh{\bH}_{m,k} \wh{\bH}_{m,k}^\trans = \bI + o_P(1)$. Moreover, the PCA-type estimator $\wh{\bB}_{m,k}$ satisfies
    \[
    \frac{1}{p_{m,k}} \Big\| \wh{\bB}_{m,k} - \bB_{m,k} \wh{\bH}_{m,k} \Big\|_F^2 = \cO_P\Big( \frac{1}{p_{m,k}^2} + \frac{1}{T p_{m,k}} + \frac{1}{T p_{m,\text{-}k}} \Big) .
    \]
    \item For any $m\in[M]$ with $K_m= 1$, let $\wt{\bD}_{m,1}$ be the diagonal matrix consisting of the first $u_{m,1}$ eigenvalues of $\wt\bSigma_{B,m,1}$ defined in \eqref{eqn: def_tilde_Sigma_B_m1}. Define
    \[
    \wt{\bH}_{m,1} = \frac{1}{Tp_{m,1}} \sum_{t=1}^T \mat_1(\cF_t^{(m)}) \mat_1(\cF_t^{(m)})^\trans \bB_{m,1}^\trans \wh{\bA}_{m,1,\perp} \wh{\bA}_{m,1,\perp}^\trans \bB_{m,1}^\trans \wh{\bB}_{m,1} \wt{\bD}_{B,m,1}^{-1} .
    \]
    Then $\wt{\bH}_{m,1}$ is asymptotically invertible, and the estimator $\wh{\bB}_{m,1}$ based on \eqref{eqn: def_tilde_Sigma_B_m1} satisfies
    \[
    \frac{1}{p_{m,1}} \Big\| \wh{\bB}_{m,1} - \bB_{m,1} \wt{\bH}_{m,1} \Big\|_F^2 = \cO_P\Big( \frac{1}{p_{m,1}^2} + \frac{1}{T} \Big) .
    \]
\end{enumerate}
\end{theorem}



\begin{theorem}[Asymptotic consistency of core factor estimators and component estimators]\label{thm: core_factor}
Let assumptions in \Cref{thm: local_loading} hold. Then for any $m\in[M]$, $t\in[T]$:
\begin{enumerate}[itemsep=0pt, label = (\alph*), left = 0pt]
    \item The local factor estimator defined in \eqref{eqn: def_hat_Ft} satisfies that
    \begin{align*}
        & \text{for } K_m> 1, \,
        \Big\| \vec(\wh{\cF}_t^{(m)}) - (\otimes_{k=1}^K \wh{\bH}_{m,k}^\trans) \vec(\cF_t^{(m)}) \Big\|^2 = \cO_P\Big( \sum_{k=1}^{K_m} \frac{1}{p_{m,k}^2} + \frac{1}{T} + \frac{1}{p_m} \Big) ; \\
        & \text{for } K_m =1, \,
        \Big\| \vec(\wh{\cF}_t^{(m)}) - \wt{\bH}_{m,1}^\trans \vec(\cF_t^{(m)}) \Big\|^2 = \cO_P\Big( \frac{1}{T} + \frac{1}{p_{m,1}} \Big) .
    \end{align*}
    \item The global factor estimator defined in \eqref{eqn: def_hat_Gt} is consistent such that 
    \begin{align*}
        \big\| \vec(\wh{\cG}_t^{(m)}) - \vec(\cG_t^{(m)}) \big\|^2 = \cO_P\Big( \sum_{k=1}^{K_m} \frac{1}{p_{m,k}^2} + \frac{1}{T} + \frac{1}{p_m} \Big) .
    \end{align*}
    \item Both global and local component estimators are consistent such that
    \begin{align*}
        \frac{1}{p_m} \big\| \vec(\wh\cX_{G,t}^{(m)}) -\vec(\cX_{G,t}^{(m)}) \big\|^2 &= \cO_P\Big( \sum_{k=1}^{K_m} \frac{1}{p_{m,k}^2} + \frac{1}{T} + \frac{1}{p_m} \Big) , \\
        \frac{1}{p_m} \big\| \vec(\wh\cX_{F,t}^{(m)}) -\vec(\cX_{F,t}^{(m)}) \big\|^2 &= \cO_P\Big( \sum_{k=1}^{K_m} \frac{1}{p_{m,k}^2} + \frac{1}{T} + \frac{1}{p_m} \Big) .
    \end{align*}
\end{enumerate}
\end{theorem}


\Cref{thm: local_loading}(a) spells out the rate of convergence for the local loading estimator under $K_m>1$, and part~(b) under $K_m=1$. As mentioned previously, this result is generally better than our \Cref{thm: global_loading}. In fact, in a balanced-dimension scenario with $T\asymp p_{m,1}\asymp \dots \asymp p_{m,K_m}$, the rate in (a) becomes $1/T^2$ which is better than the result in Theorem~3 in \cite{Zhangetal2025}. Until $p_{m,k}$ is too small such that $p_{m,k}^2 \asymp T$, both rates in \Cref{thm: local_loading}(a) and (b) boil down to $1/T$.

Compared with the classical rate of $1/p_{m,k}^2 + 1/Tp_{m,\text{-}k}$ for PCA-type estimators \citep[e.g.][]{barigozzi2022statistical}, \Cref{thm: local_loading}(a) contains an additional $1/Tp_{m,k}$. This is inherited from the error of the global loading estimator, and would have no impact to the result when $p_{m,1}\asymp \dots \asymp p_{m,K_m}$ and $T$ is of order larger than or equal to $p_{m,1}$, which is often the case in practice.


\Cref{thm: core_factor} spells out the consistency of local/global factor/component estimators, and indicate the same rates of convergence. Although we have a superior result in \Cref{thm: local_loading}, it cannot be translated to \Cref{thm: core_factor} which suffers from the $1/T$ rate in \Cref{thm: global_loading}. This once again reflects the difficulty in parameter estimation under our model. To compare with multilevel factor models, for the vector case such that $K_m=2$ for all $m\in[M]$, our result is the same as Proposition~2 in \cite{Choietal2018}. Under the matrix case when $K_m=2$ and $p_{m,1}\asymp p_{m,2}$ for all $m\in[M]$, our derived asymptotic rates would be the same as Theorems~4 and 5 in \cite{Zhangetal2025}. Thus, our presented results can be also viewed as pointing out the behaviour of parameter estimators in a multilevel factor model for tensor time series with general orders.




\section{Further Discussion}
\label{sec: further_discussion}



\subsection{Explicit form of Global Factors}
\label{subsec: explicit_global_fac}


In \Cref{subsec: model}, JXC following a multi-order tensor factor model means each $\cX_t^{(m)}$ has the representation \eqref{eqn: model}. Although the unspecified form of the global factor is more general, it can be of interest to assume explicit connections among $\big\{ \cG_t^{(1)}, \dots, \cG_t^{(M)} \big\}$. If the tensor orders are the same over $m\in[M]$, a direct choice is $\cG_t^{(1)} =\dots =\cG_t^{(M)}$. Yet with the potentially different orders, the setup is not so straightforward.

With the notation of tensor map in \Cref{sec: intro_map}, we consider a natural choice that $\cG_t^{(m)} = \map(\cG_t, V_m)$ for some order-$K_M$ tensor $\cG_t$ with dimension $r_1\times \dots\times r_{K_M}$, and some pre-specified channel $V_m$.
%such that
%\[ V_m=\Big\{ \{\theta_{m,0}+1, \dots, \theta_{m,1}\}, \{\theta_{m,1}+1, \dots, \theta_{m,2}\}, \dots, \{\theta_{m,K_m-1}+1, \dots, \theta_{m,K_m}\} \Big\} ,\]
%where $\theta_{m,0}=0$, $\theta_{m,K_m}=K_m$, and $\theta_{m,0} <\dots< \theta_{m,K_m}$. 
Essentially, $\cG_t^{(m)} = \map(\cG_t, V_m)$ is some re-organisation of $\cG_t$.
This can be useful if the meaning of modes are related across threads with different orders. For example, one time series contains the fMRI scans of brains over time, while the other reports vital signs of patients such as heart rate, blood pressure, and respiratory rate. Then the sign series is potentially driven by the vectorized core factor for the fMRI series.

Without loss of generality, let $V_M=\big\{\{1\}, \dots, \{K_M\} \big\}$. We may read \eqref{eqn: model} as
\begin{equation}
\label{eqn: model_map_Gt}
\cX_t^{(m)} = \map(\cG_t, V_m) \times_{k=1}^{K_m} \bA_{m,k} + \cF_t^{(m)} \times_{\ell=1}^{K_m} \bB_{m,\ell} + \cE_t^{(m)} .
\end{equation}
By the definition of map operator, $\map(\cG_t, V_m)$ is an order-$K_m$ tensor with each mode-$j$ dimension $r_{m,j} = \prod_{\theta_{m,j-1}+1 <i \leq \theta_{m,j}} r_i$. As a simple example, consider $M=2$ and $(K_1, K_2)=(1,2)$, i.e., the observed JXC consists of a vector time series and a matrix one. Note that $V_1 =\big\{ \{1,2\} \big\}$ and $\map(\cG_t, V_1)=\reshape(\cG_t, \{1,2\}) =\vec(\cG_t)$. Then for each $t\in[T]$, we have
\begin{align*}
    & \cX_t^{(1)} = \bA_{1,1} \vec(\cG_t) + \bB_{1,1} \cF_t^{(1)} + \cE_t^{(1)}
    \quad \in \R^{p_{1,1}} ,\\
    & \cX_t^{(2)} = \bA_{2,1} \cG_t \bA_{2,2}^\trans + \bB_{2,1} \cF_t^{(2)} \bB_{2,2}^\trans + \cE_t^{(2)} 
    \quad \in \R^{p_{2,1} \times p_{2,2}} .
\end{align*}

One advantage of the form \eqref{eqn: model_map_Gt} is that \Cref{ass: core_factor}(c) can be easily satisfied; see Remark~\ref{remark: explicit_fac_satisfy_assump}.
Furthermore, even though all methods and results for \eqref{eqn: model} hold for \eqref{eqn: model_map_Gt}, \Cref{thm: core_factor} implies the global factor estimator has the rate $\sum_{k\in[K_m]} 1/p_{m,k}^2 + 1/T + 1/p_m$. Thus, we can estimate the global factor on $m$th thread such that $\sum_{k\in[K_m]} 1/p_{m,k}^2 + 1/p_m$ is minimised, and invert the map operator to estimate $\cG_t$.

% stack and estimate one single $\cG_t$




\subsection{Estimation of number of factors}
\label{subsec: est_num_fac}

In model~\eqref{eqn: model}, the numbers of global and local factors are unknown, which need to be either pre-specified based on some prior information or appropriately estimated. For the latter, the representation in \eqref{eqn: Xt_rewrite_tfm} enlightens a feasible approach. In particular, $\{\cX_t^{(m)}\}$ follows a classical tensor factor model with each mode-$k$ loading matrix $(\bA_{m,k}, \bB_{m,k})$ with column rank $r_{m,k} +u_{m,k}$ (see also \Cref{ass: loadings}(ii) in \Cref{subsec: assumption}). Hence any consistent factor number estimators can be employed, such as the correlation-thresholding estimator by \cite{Lam2021}, the information-criteria approach by \cite{Hanetal2022}, etc. Denote such estimated factor numbers by $\wh{s}_{m,k}$ which effectively estimates $r_{m,k} +u_{m,k}$. 

Inspired by the eigenvalue-ratio estimator in \cite{Zhangetal2025}, a plausible estimator on the global factor number can be obtained by minimising $\lambda_{i+1}(\wh\bSigma_{m,k}) /\lambda_i (\wh\bSigma_{m,k})$ over $i\in[p_{m,k}]$. However, no theoretical guarantee is made, which is the similar case for the eigenvalue-ratio estimator briefly mentioned in \cite{Wangetal2019}. In fact, Example~1 in \cite{Zhangetal2024} demonstrate that using the direct ratio between consecutive eigenvalues could lead to inconsistent estimator. To circumvent this, suppose $r_{\max,m,k}$ is some upper bound of the factor numbers, we define
\begin{equation}
\label{eqn: est_num_fac_define}
\wh{r}_{m,k} := \argmin_{1\leq i\leq r_{\max,m,k}} \frac{\lambda_{i+1}(\wh\bSigma_{m,k}) + \xi_{m,k}}{ \lambda_i (\wh\bSigma_{m,k}) +\xi_{m,k}} , \quad \text{where} \quad
\xi_{m,k} \asymp \frac{p_m}{\sqrt{T}} \Bigg( \sum_{\substack{n=1\\ n\neq m}}^M p_n \Bigg) .
\end{equation}
In practice, we can take $r_{\max,m,k}=\lfloor p_{m,k}/2 \rfloor$ and $\xi_{m,k} = p_m(\sum_{n\in[M] \setminus \{m\}} p_n) /(5\sqrt{T})$. In contrast to the practical estimator in \cite{Zhangetal2024}, we show $\wh{u}_{m,k}$ is consistent as follows.


\begin{theorem}\label{thm: eigenratio}
Let all assumptions in \Cref{thm: global_loading} hold. Then for any $m\in[M]$, $k\in[K_m]$, the estimator defined in \eqref{eqn: est_num_fac_define} is consistent such that
\[
\P(\wh{r}_{m,k} = r_{m,k} ) \xrightarrow{P} 1.
\]
\end{theorem}

The form in \eqref{eqn: est_num_fac_define} is akin to the perturbed eigenvalue-ratio estimator in Theorem~6 in \cite{Pelger2019}. Intuitively, $\xi_{m,k}$ helps to stabilise the estimator. If one takes $\xi_{m,k}=0$, then we can only show $\P(\wh{r}_{m,k} \geq r_{m,k} ) \xrightarrow{P} 1$. With both $\wh{s}_{m,k}$ and $\wh{r}_{m,k}$ computed, the local factor number can thus be estimated by $\wh{u}_{m,k} := \wh{s}_{m,k} -\wh{r}_{m,k}$, whose consistency is straightforward provided the consistent estimation discussed previously.

As briefly discussed in \Cref{remark: absent_component}, there might be no local components for some threads. Presuming the existence of global components, $\wh{u}_{m,k}=0$ would hint on the absence of local component. Following this, estimators for the local loading matrix and local factor can be trivially set as zero in Sections~\ref{subsec: est_loading} and \ref{subsec: est_fac}, and other procedures therein remain valid. We leave the refinement of this to future research.








\subsection{Fast implementation}
\label{subsec: bootstrap_implement}

When threads in JXC have large dimensions and higher orders, obtaining \eqref{eqn: def_hat_Sigma_mk} can be computationally heavy. To circumvent this, recall from \Cref{thm: global_loading} that only $T$ is involved in the rate of convergence for the global loading matrix estimator. This hints on reformulating \eqref{eqn: def_hat_Sigma_mk} by using a random subset of multi-indices. Indeed, this can be formalized as follows. For each $m\in[M]$, $k\in[K_m]$, define
\begin{equation}
\label{eqn: def_hat_Sigma_mk_bootstrap}
\begin{split}
    \wh\bSigma_{m,k}^S &:= \sum_{\substack{n=1\\ n\neq m}}^M \sum_{\bi \in \cS_n} \wh\bOmega_{k,n,\bi}^{(m)} \wh\bOmega_{k,n,\bi}^{(m) \trans} ,
\end{split}
\end{equation}
where $\cS_n \subseteq [p_{n,1}] \times \cdots \times[p_{n,K_n}]$ is a random set drawn independently of JXC, and $\wh\bOmega_{k,n,\bi}^{(m)}$ is the same as in \eqref{eqn: def_hat_Sigma_mk}. Then the global loading matrix can be estimated by $\wh{\bA}_{m,k}^S$, defined as $\sqrt{p_{m,k}}$ times the $r_{m,k}$ leading eigenvectors of $\wh\bSigma_{m,k}^S$. More importantly, the performance of such estimator can be made mathematically rigorous, as shown below.


\begin{theorem}
\label{prop: global_loading_bootstrap}
Let all assumptions in \Cref{thm: global_loading} hold. Fix $m$, and for any $n\in[M] \setminus\{m\}$, let $\cS_n^\dagger \subseteq [p_n]$ denote the set of row indices of the matrix $\otimes_{k=1}^{K_n} \bA_{n,k}$ that corresponds to the rows of $\{\bA_{n,1}, \dots, \bA_{n,K_n}\}$ with indices in $\cS_n$ from \eqref{eqn: def_hat_Sigma_mk_bootstrap}. For any $j\in[r_n]$, assume that
\[
\sigma_{j}\big\{ (\otimes_{k=1}^{K_n} \bA_{n,k} )_{\cS_n^\dagger} \big\} \asymp |\cS_n^\dagger |^{1/2} ,
\]
where $(\otimes_{k=1}^{K_n} \bA_{n,k} )_{\cS_n^\dagger}$ represents the sub-matrix of $\otimes_{k=1}^{K_n} \bA_{n,k}$ with rows restricted on $\cS_n^\dagger$. Then it holds for any $k\in[K_m]$ that
\begin{align*}
    \frac{1}{p_{m,k}} \Big\| \wh{\bA}_{m,k}^S -\bA_{m,k} \Big\|_F^2 = \cO_P\Big( \frac{1}{T} \Big) .
\end{align*}
\end{theorem}


The additional assumption in \Cref{prop: global_loading_bootstrap} can be regarded as a rank condition to ensure $\wh\bSigma_{m,k}^S$ is well behaved. By \Cref{ass: loadings}(b), it is automatically fulfilled if $\cS_n^\dagger =[p_n]$, i.e., $\wh\bSigma_{m,k}^S =\wh\bSigma_{m,k}$. Through our numerical studies in \Cref{subsec: simulation}, this additional assumption can be satisfied in general. It is also worth pointing out that $\wh\bSigma_{m,k}^S$ can replace $\wh\bSigma_{m,k}$ to compute the factor number estimator in \eqref{eqn: est_num_fac_define}, except that $\xi_{m,k}$ therein needs to fulfil
\[
\xi_{m,k} \asymp \frac{p_m}{\sqrt{T}} \Bigg( \sum_{\substack{n=1\\ n\neq m}}^M |\cS_n^\dagger | \Bigg) .
\]
Then the resulted estimator is still consistent.





\subsection{*How Assumption~\ref{ass: core_factor}(c) can be implied}
\label{subsec: assump_fac_satisfied}

In what follows, we showcase how \Cref{ass: core_factor}(c) can be satisfied for $K_1=\dots=K_M$, i.e., when our framework boils down to a multilevel tensor factor model. As a natural setup, the global factors are the same over threads, so we assume $\cG_t^{(1)} = \dots =\cG_t^{(M)}$ throughout this subsection.

First of all, note that the matrix in the expectation within \Cref{ass: core_factor}(c) has dimension $r_{m,k} \times (p_n r_{m,\text{-}k})$, so the assumption can be treated as a full rank condition on
{\small
\[
\E\Big[ \mat_k(\cG_t^{(m)}) \otimes \Big\{ \vec(\cG_t^{(n)})^\trans \big( \otimes_{h=1}^{K_n} \bA_{h} \big)^\trans \Big\} \Big] \E\Big[ \mat_k(\cG_t^{(m)}) \otimes \Big\{ \vec(\cG_t^{(n)})^\trans \big( \otimes_{h=1}^{K_n} \bA_{h} \big)^\trans \Big\} \Big]^\trans .
\] }
Inspired by this observation, consider the vector case ($K_1=1$). Then we may simplify
\begin{align*}
    \E\Big[ \mat_k(\cG_t^{(m)}) \otimes \Big\{ \vec(\cG_t^{(n)})^\trans \big( \otimes_{h=1}^{K_n} \bA_{h} \big)^\trans \Big\} \Big] = \E\Big\{ \cG_t^{(m)} (\cG_t^{(m)})^\trans \Big\} \bA_{1}^\trans ,
\end{align*}
which implies that \Cref{ass: core_factor}(c) is equivalent to $\E\big\{ \cG_t^{(m)} (\cG_t^{(m)})^\trans \big\}$ being positive definite with eigenvalues bounded away from zero and infinity. In general, \Cref{ass: core_factor}(c) can be replaced by simpler assumptions for any tensor orders. We present the result in the following claim, where the extra assumption on the global factors is not necessary (e.g., as previously shown for the vector case) and can be implied by Assumption~(F1) in \cite{CenLam2025} for instance. The proof of the claim is relegated to the online supplement.

\begin{claim}
\label{claim: assump_fac_simpler}
Under $K_1=\dots=K_M$ and $\cG_t^{(1)} = \dots =\cG_t^{(M)}$, if the global factors have uncorrelated elements with zero mean and bounded variance, then \Cref{ass: core_factor}(c) is satisfied.
\end{claim}


\begin{remark}
\label{remark: explicit_fac_satisfy_assump}
In \Cref{subsec: explicit_global_fac}, different global factors are re-organisation of each other. Hence in \Cref{ass: core_factor}(c), $\vec(\cG_t^{(m)})$ is the same as $\vec(\cG_t^{(n)})$ up to some permutation which can be absorbed into $\otimes_{h=1}^{K_n} \bA_{h}$. Since the Kronecker product structure is preserved by the reshape operator \citep{CenLam2025_KronProd}, all arguments in the proof of \Cref{claim: assump_fac_simpler} follow. This indicates that if $\cG_t$ in \Cref{subsec: explicit_global_fac} has uncorrelated elements with zero mean and bounded variance, then \Cref{ass: core_factor}(c) holds.
\end{remark}





\section{Numerical Analysis}
\label{sec: numerical}

\subsection{Simulation}
\label{subsec: simulation}

In the following, we demonstrate the numerical performance of the estimators described in Sections~\ref{subsec: est_loading} and \ref{subsec: est_fac} using Monte Carlo experiments. As this article includes multilevel factor model as a special case, we will also make comparison with \cite{Zhangetal2025} when JXC only contains matrix time series.

To start with, each thread is generated according to \eqref{eqn: model}, except that we consider the setup in \eqref{eqn: model_map_Gt} such that the global factor for each thread is mapped from $\cG_t^{(M)}$. In detail, each entry of the global and local factors is independent autoregressive process AR(1), with coefficient $0.5$ and innovation following \text{i.i.d.} $\cN(0,1)$. The entries of $\cE_t^{(m)}$ is constructed similarly, except that the innovation is either $\cN(0,1)$ or $t_6$, which will be specified later. Each global loading matrix is generated as $\bA_{m,k} = \bU_{m,k}$ for pervasive factors and $\bA_{m,k} = p_{m,k}^{-0.1} \cdot \bU_{m,k}$ for weak factors, where $\bU_{m,k}\in \R^{p_{m,k}\times r_{m,k}}$ consists of \text{i.i.d.} $\cN(0,1)$ elements. Let $\bA_{m,k,\perp}$ be the half-orthogonal matrix representing the orthogonal complement of $\bA_{m,k}$. Then the local loading matrix is constructed as $p_{m,k}^{0.5}$ (or $p_{m,k}^{0.4}$ for weak factors) times a $p_{m,k}\times u_{m,k}$ sub-matrix of $\bA_{m,k,\perp}$. For simplicity, we only consider one factor for all modes and threads.


As the loading matrices are only identified up to an arbitrary rotation, we measure the distance between the column space spanned by the true loading matrix and the estimator. Formally, given any pair of matrices $\bX$ and $\bY$, the column space distance is computed as
\[
\cD(\bX, \bY) := \Big\| \bX(\bX^\trans \bX)^{-1} \bX^\trans -\bY(\bY^\trans \bY)^{-1} \bY^\trans \Big\| .
\]
Due to the rotational indeterminacy in the loading matrices, performance of the core factor estimator cannot be directly compared, but reflected in the estimated global and local components. We use the (relative) mean squared error to measure the performance of the component estimators:
\begin{align*}
    & \text{MSE}\big(\wh\cX_{G,t}^{(m)} \big) := \frac{ \sum_{t=1}^{T} \big\| \wh\cX_{G,t}^{(m)} - \cX_{G,t}^{(m)} \big\|_F^2}{ \sum_{t=1}^{T} \big\| \cX_{G,t}^{(m)} \big\|_F^2} , \quad
    \text{MSE}\big(\wh\cX_{F,t}^{(m)} \big) := \frac{\sum_{t=1}^{T} \big\| \wh\cX_{F,t}^{(m)} - \cX_{F,t}^{(m)} \big\|_F^2}{ \sum_{t=1}^{T} \big\| \cX_{F,t}^{(m)} \big\|_F^2} .
\end{align*}


We experiment a variety of settings, detailed as follows. Each setting is repeated 400 times, unless otherwise stated.

\vspace{12pt}
\noindent\textbf{Setting~A:}
\begin{enumerate}[itemsep=-1pt, parsep=0pt, topsep=0pt, label = (A.\arabic*), left=1pt]
    \item Number of threads $M=3$ with $(K_1,K_2,K_3)=(1,1,2)$. Fix the cross-sectional dimensions $(p_{1,k}, p_{2,h}, p_{3,l}) =(30, 30, 10)$, and the sample size is experimented through $T\in\{100, 400\}$. Factors are strong and the noise innovation is Gaussian.
    \item Same as (A.1), but $(K_1,K_2,K_3)=(1,2,3)$ and $(p_{1,k}, p_{2,h}, p_{3,l}) =(30, 15, 10)$.
    \item Same as (A.1), except that factors are weak.
    \item Same as (A.1), except that the noise innovation is $t$ distributed.
\end{enumerate}
\vspace{12pt}

\noindent\textbf{Setting~B:}
\begin{enumerate}[itemsep=-1pt, parsep=0pt, topsep=0pt, label = (B.\arabic*), left=1pt]
    \item Number of threads $M=2$ with $(K_1,K_2)=(1,2)$. Fix dimension $(p_{1,k}, p_{2,h}) =(30, 10)$, and the sample size is experimented through $T\in\{100, 400\}$. Factors are strong and the noise innovation is Gaussian.
    \item Same as (B.1), except that $(p_{1,k}, p_{2,h}) =(100, 20)$.
    \item Same as (B.1), except that $(K_1,K_2)=(1,3)$ and $(p_{1,k}, p_{2,h}) =(30, 10)$.
    \item Same as (B.1), except that $(K_1,K_2)=(2,3)$ and $(p_{1,k}, p_{2,h}) =(10, 10)$.
\end{enumerate}
\vspace{12pt}


Settings~A and B consider $M=3$ and $M=2$, respectively. In summary, (A.1) showcases the scenario when some threads can have the same order but different from another thread, while the JXC in (A.2) has threads with different orders. (A.3) and (A.4) respectively illustrate the sensitivity of (A.1) under weak factors and heavy-tailed noise. In Setting~B, we demonstrate scenarios with different combination of thread orders. (B.2) also investigates the effect of increasing cross-sectional dimension to the estimator performance.
For computational concern, we estimate the global loading matrix using $\wh{\bA}_{m,k}^S$ discussed in \Cref{subsec: bootstrap_implement}, with $\sum_{n\in[M] \setminus\{m\}}|\cS_n^\dagger| =50$ in Setting~A, and $\sum_{n\in[M] \setminus\{m\}}|\cS_n^\dagger| =30$ in Setting~B which is the maximum possible in some settings.





%\usepackage{booktabs}
%\newcommand{\ra}[1]{\renewcommand{\arraystretch}{#1}}
\begin{table}%[htp!]\centering
\caption{Results of Settings~(A.1)--(A.4). Each cell shows the mean of measure over 400 runs for the corresponding setting.}
\label{tab: simulation_setting_A}
%\small
%\ra{1.2}
\begin{tabular}{@{}crrrrrrrrrrrrrrrrrrrrrrrrr@{}}
\hline %\toprule
Setting A & \multicolumn{2}{c}{(A.1)} && \multicolumn{2}{c}{(A.2)} && \multicolumn{2}{c}{(A.3)} && \multicolumn{2}{c}{(A.4)}  \\
%\cmidrule{2-2}
$T$
&  $100$ & $400$ && $100$ & $400$ && $100$ & $400$
&& $100$ & $400$ \\
\hline %\midrule
%\cmidrule{2-3} \cmidrule{5-6} \cmidrule{8-9} \cmidrule{11-12}
$\cD(\bA_{1,1}, \wh{\bA}_{1,1}^S)$ &  0.189 & 0.088 && 0.196 & 0.095  && 0.237 & 0.111 && 0.216 & 0.101  \\
$\cD(\bA_{2,1}, \wh{\bA}_{2,1}^S)$ &  0.194 & 0.092 && 0.039 & 0.019  && 0.243 & 0.116 && 0.217 & 0.105 \\
$\cD(\bA_{2,2}, \wh{\bA}_{2,2}^S)$ &  - & - && 0.045 & 0.020  && - & - && - & -  \\
$\cD(\bA_{3,1}, \wh{\bA}_{3,1}^S)$ &  0.045 & 0.022 && 0.018 & 0.008  && .073 & .035 && .055 & .027 \\
$\cD(\bA_{3,2}, \wh{\bA}_{3,2}^S)$ &  0.046 & 0.022 && 0.016 & 0.008  && 0.075 & 0.035 && 0.055 & 0.027 \\
$\cD(\bA_{3,3}, \wh{\bA}_{3,3}^S)$ &  - & - && 0.018 & 0.008  && - & - && - & -  \\
\\
$\cD(\bB_{1,1}, \wh{\bB}_{1,1})$ &  0.186 & 0.087 && 0.191 & 0.093  && 0.232 & 0.109 && 0.211 & 0.100  \\
$\cD(\bB_{2,1}, \wh{\bB}_{2,1})$ &  0.192 & 0.089 && 0.038 & 0.016  && 0.240 & 0.112 && 0.216 & 0.102 \\
$\cD(\bB_{2,2}, \wh{\bB}_{2,2})$ &  - & - && 0.033 & 0.016  && - & - && - & -  \\
$\cD(\bB_{3,1}, \wh{\bB}_{3,1})$ &  0.040 & 0.020 && 0.013 & 0.006  && 0.068 & 0.034 && 0.049 & 0.024 \\
$\cD(\bB_{3,2}, \wh{\bB}_{3,2})$ &  0.040 & 0.020 && 0.013 & 0.006  && 0.068 & 0.033 && 0.049 & 0.024 \\
$\cD(\bB_{3,3}, \wh{\bB}_{3,3})$ &  - & - && 0.013 & 0.006  && - & - && - & -   \\
\\
$\text{MSE}\big(\wh\cX_{F,t}^{(1)} \big)$ &  0.095 & 0.046 && 0.106 & 0.051  && 0.152 & 0.084 && 0.125 & 0.066  \\
$\text{MSE}\big(\wh\cX_{F,t}^{(2)} \big)$ &  0.098 & 0.047 && 0.149 & 0.005  && 0.159 & 0.085 && 0.130 & 0.067 \\
$\text{MSE}\big(\wh\cX_{F,t}^{(3)} \big)$ &  0.014 & 0.011 && 0.002 & 0.001  && 0.036 & 0.028 && 0.021 & 0.016 \\
$\text{MSE}\big(\wh\cX_{G,t}^{(1)} \big)$ &  0.111 & 0.049 && 0.122 & 0.054  && 0.173 & 0.088 && 0.143 & 0.069 \\
$\text{MSE}\big(\wh\cX_{G,t}^{(2)} \big)$ &  0.114 & 0.052 && 0.229 & 0.008  && 0.181 & 0.094 && 0.149 & 0.075 \\
$\text{MSE}\big(\wh\cX_{G,t}^{(3)} \big)$ &  0.020 & 0.016 && 0.006 & 0.002  && 0.052 & 0.040 && 0.030 & 0.024  \\
\hline %\bottomrule
\end{tabular}
\end{table}


%\usepackage{booktabs}
%\newcommand{\ra}[1]{\renewcommand{\arraystretch}{#1}}
\begin{table}%[htp!]\centering
\caption{Results of Settings~(B.1)--(B.4). Each cell shows the mean of measure over 400 runs for the corresponding setting.}
\label{tab: simulation_setting_B}
%\small
%\ra{1.2}
\begin{tabular}{@{}crrrrrrrrrrrrrrrrrrrrrrrrr@{}}
\hline %\toprule
Setting B & \multicolumn{2}{c}{(B.1)} && \multicolumn{2}{c}{(B.2)} && \multicolumn{2}{c}{(B.3)} && \multicolumn{2}{c}{(B.4)}  \\
%\cmidrule{2-2}
$T$
&  $100$ & $400$ && $100$ & $400$ && $100$ & $400$
&& $100$ & $400$ \\
\hline %\midrule
%\cmidrule{2-3} \cmidrule{5-6} \cmidrule{8-9} \cmidrule{11-12}
$\cD(\bA_{1,1}, \wh{\bA}_{1,1}^S)$ &  0.203 & 0.097 && 0.202 & 0.094  && 0.212 & 0.103 && 0.057 & 0.027  \\
$\cD(\bA_{1,2}, \wh{\bA}_{1,2}^S)$ &  - & - && - & -  && - & - && 0.052 & 0.026  \\
$\cD(\bA_{2,1}, \wh{\bA}_{2,1}^S)$ &  0.046 & 0.023 && 0.032 & 0.015  && 0.022 & 0.008 && .017 & .008 \\
$\cD(\bA_{2,2}, \wh{\bA}_{2,2}^S)$ &  0.049 & 0.023 && 0.032 & 0.015  && 0.022 & 0.008 && 0.017 & 0.009  \\
$\cD(\bA_{2,3}, \wh{\bA}_{2,3}^S)$ &  - & - && - & -  && 0.021 & 0.008 && 0.024 & 0.009  \\
\\
$\cD(\bB_{1,1}, \wh{\bB}_{1,1})$ &  0.196 & 0.093 && 0.193 & 0.091  && 0.201 & 0.097 && 0.042 & 0.020  \\
$\cD(\bB_{1,2}, \wh{\bB}_{1,2})$ &  - & - && - & -  && - & - && 0.040 & 0.019  \\
$\cD(\bB_{2,1}, \wh{\bB}_{2,1})$ &  0.040 & 0.020 && 0.028 & 0.014  && 0.017 & 0.006 && 0.015 & 0.006 \\
$\cD(\bB_{2,2}, \wh{\bB}_{2,2})$ &  0.039 & 0.020 && 0.028 & 0.014  && 0.017 & 0.006 && 0.014 & 0.006  \\
$\cD(\bB_{2,3}, \wh{\bB}_{2,3})$ &  - & - && - & -  && 0.017 & 0.006 && 0.013 & 0.006  \\
\\
$\text{MSE}\big(\wh\cX_{F,t}^{(1)} \big)$ &  0.106 & 0.052 && 0.084 & 0.024  && 0.120 & 0.055 && 0.016 & 0.011  \\
$\text{MSE}\big(\wh\cX_{F,t}^{(2)} \big)$ &  0.014 & 0.011 && 0.004 & 0.003  && 0.007 & 0.001 && 0.004 & 0.001 \\
$\text{MSE}\big(\wh\cX_{G,t}^{(1)} \big)$ &  0.127 & 0.059 && 0.099 & 0.026  && 0.140 & 0.061 && 0.028 & 0.018 \\
$\text{MSE}\big(\wh\cX_{G,t}^{(2)} \big)$ &  0.021 & 0.018 && 0.006 & 0.003  && 0.125 & 0.002 && 0.009 & 0.002 \\
\hline %\bottomrule
\end{tabular}
\end{table}






Results for Settings~A and B are reported in Tables~\ref{tab: simulation_setting_A} and \ref{tab: simulation_setting_B}, respectively. First of all, from $T=100$ to $T=400$, the improvement of global loading matrix estimators in all settings corroborates with the rate of convergence presented in Theorem~\ref{thm: global_loading} and \Cref{prop: global_loading_bootstrap}. In Settings~(A.3) and (A.4), all estimation errors are slight inflated, compared to those in Setting~(A.1), but still remains satisfactory. This shows our parameter estimators are quite robust to different factor strengths and heavy-tailed distributed noise. Secondly, note that errors of loading matrices are similar when they are in threads with the same order and dimensionality, \text{e.g.} $\cD(\bA_{1,1}, \wh{\bA}_{1,1}^S)$ and $\cD(\bA_{2,1}, \wh{\bA}_{2,1}^S)$ in (A.1) compared to $\cD(\bA_{1,1}, \wh{\bA}_{1,1}^S)$ in (A.2). It is consistent with the fact that in \Cref{subsec: theorem}, theoretical results for parameters in one thread involve no dimensionality from other threads. The interpretation in \Cref{tab: simulation_setting_B} is analogous to this, and performance is generally better off with larger dimensionality from the results in Setting~(B.2).


As mentioned previously, our proposed framework includes the multilevel matrix factor model by \cite{Zhangetal2025} as a special case. From \Cref{remark: hat_Sigma_rewrite}, our global loading estimator described in \Cref{subsec: est_loading} is the same as theirs and hence would give the same result. In contrast, our local loading estimator is generally more superior than theirs, according to \Cref{thm: local_loading}. We verify this by numerical results shown below. For all settings on this, we fix $M=2$ with $K_1=K_2=2$, and set $T=40$. All factors are pervasive and the noise has Gaussian innovation. We consider two settings representing different noise dynamics, specified below.
%Estimation of the global loading matrix is similar as for Setting~A.

\vspace{12pt}
\noindent\textbf{Setting~C:}
\begin{enumerate}[itemsep=-1pt, parsep=0pt, topsep=0pt, label = (C.\arabic*), left=1pt]
    \item Cross-sectional dimensions are the same such that $p_{1,1}=p_{1,2}=p_{2,1}=p_{2,2}$, and the noise is white such that its AR coefficient is zero. We experiment through $p_{1,1}\in \{10, 40\}$.
    \item Same as (C.1), except that the AR coefficient of noise is 0.5.
\end{enumerate}
\vspace{12pt}



%\usepackage{booktabs}
%\newcommand{\ra}[1]{\renewcommand{\arraystretch}{#1}}
\begin{table}%[htp!]\centering
\caption {Results of local loading estimators under Settings~(C.1)--(C.2), using methods proposed in this paper and in \cite{Zhangetal2025}. Each cell shows the mean of measure over 400 runs for the corresponding setting.}
\label{tab: simulation_setting_C}
%\small
%\ra{1.2}
\begin{tabular}{@{}crrrrrrrrrrrrrrrrrrrrrrrrr@{}}
\hline %\toprule
& \multicolumn{5}{c}{Our method} && \multicolumn{5}{c}{Zhang et al.~(2025)} \\
\hline %\midrule
Setting C & \multicolumn{2}{c}{(C.1)} && \multicolumn{2}{c}{(C.2)} && \multicolumn{2}{c}{(C.1)} && \multicolumn{2}{c}{(C.2)}  \\
%\cmidrule{2-2}
$p_{1,1}$
&  $10$ & $40$ && $10$ & $40$ && $10$ & $40$
&& $10$ & $40$ \\
\hline %\midrule
%\cmidrule{2-3} \cmidrule{5-6} \cmidrule{8-9} \cmidrule{11-12}
$\cD(\bB_{1,1}, \wh{\bB}_{1,1})$ &  0.076 & 0.022 && 0.100 & 0.037  && 0.108 & 0.042 && 0.159 & 0.071  \\
$\cD(\bB_{1,2}, \wh{\bB}_{1,2})$ &  0.083 & 0.025 && 0.106 & 0.042  && 0.110 & 0.043 && 0.160 & 0.070  \\
$\cD(\bB_{2,1}, \wh{\bB}_{2,1})$ &  0.075 & 0.027 && 0.140 & 0.047  && 0.111 & 0.045 && 0.178 & 0.070  \\
$\cD(\bB_{2,2}, \wh{\bB}_{2,2})$ &  0.086 & 0.029 && 0.113 & 0.034  && 0.113 & 0.044 && 0.178 & 0.071  \\
\\
Run Time (s) &  0.073 & 0.091 && 0.108 & 0.089  && 22.347 & 466.275 && 33.267 & 461.613  \\
\hline %\bottomrule
\end{tabular}
\end{table}





\Cref{tab: simulation_setting_C} presents the results for Setting~C, which clearly shows the better performance of our method both in estimation accuracy and computational cost. Note that although both methods have lower estimation errors when $p_{1,1}$ increases from 10 to 40, our method enjoys faster decay on the errors, benefitting from the explicit cross-sectional dimensions involved in the convergence rates spelt out in \Cref{thm: local_loading}. Moreover, it is not surprising that the estimator in \cite{Zhangetal2025} suffers more in Setting~(C.2), since their method relied on the idiosyncratic noise being white.





Lastly, we demonstrate the performance of our global factor number estimator defined in \eqref{eqn: est_num_fac_define}. We adopt the choice of $r_{\max,m,k}$ and $\xi_{m,k}$ below \eqref{eqn: est_num_fac_define}, and consider settings described as follows.

\vspace{12pt}
\noindent\textbf{Setting~D:}
\begin{itemize}
    \item [(D.1--2)] Same as (B.1)--(B.2), except that the each mode in the second thread has two global factors. Besides the standard setup with strong factors and Gaussian noise, we also experiment weak factors and $t_6$ noise.
\end{itemize}
\vspace{12pt}



%\usepackage{booktabs}
%\newcommand{\ra}[1]{\renewcommand{\arraystretch}{#1}}
\begin{table}%[htp!]\centering
\caption{Performance of global factor number estimator in Settings~(D.1)--(D.2). Each cell is formatted as $a \,(b \mid c)$, where $a$, $b$, and $c$ respectively represent the correct, underestimation, and overestimation proportion (in $100\%$) over 400 runs for the corresponding setting. We use ``Strong'', ``Weak'', ``$\cN$'', and ``$t_6$'' to respectively denote settings with strong factors, weak factors, Gaussian noise, and heavy-tailed noise. In each setting, the first row displays results with constant $1/5$ used in $\xi_{m,k}$ and the second displays those using $\xi_{m,k}=0$.}
\label{tab: simulation_setting_D}
%\small
%\ra{1.2}
\begin{tabular}{@{}crrrrrrrrrrrrrrrrrrrrrrrrr@{}}
\hline %\toprule
Setting D & \multicolumn{2}{c}{(D.1)} && \multicolumn{2}{c}{(D.2)}  \\
%\cmidrule{2-2}
$T$
&  $100$ & $400$ && $100$ & $400$ \\
\hline %\midrule
%\cmidrule{2-3} \cmidrule{5-6} \cmidrule{8-9} \cmidrule{11-12}
\multicolumn{1}{c}{Strong, $\cN$}  \\
$r_{1,1}$ &  60.8 (32.5 $\mid$ 06.7)  & 87.5 (08.5 $\mid$ 04.0) && 77.3 (15.7 $\mid$ 07.0) & 93.8 (02.3 $\mid$ 04.0)  \\
          &  56.0 (11.0 $\mid$ 33.0)  & 77.8 (01.2 $\mid$ 21.0) && 00.0 (00.0 $\mid$ 100.0) & 00.0 (00.0 $\mid$ 100.0) \\
$r_{2,1}$ &  73.5 (04.0 $\mid$ 22.5) & 93.0 (01.0 $\mid$ 06.0) && 63.0 (00.5 $\mid$ 36.5) & 90.8 (00.0 $\mid$ 09.2) \\
          &  71.5 (04.0 $\mid$ 24.5) & 92.5 (00.8 $\mid$ 06.7) && 60.3 (00.2 $\mid$ 39.5) & 90.5 (00.0 $\mid$ 09.5) \\
$r_{2,2}$ &  70.8 (04.8 $\mid$ 24.5) & 96.5 (00.0 $\mid$ 03.5) && 63.3 (00.2 $\mid$ 36.5) & 91.3 (00.0 $\mid$ 08.7) \\
          &  69.5 (04.8 $\mid$ 25.7)  & 96.0 (00.0 $\mid$ 04.0) && 59.5 (00.8 $\mid$ 39.7) & 90.8 (00.0 $\mid$ 09.2) \\
\\
\multicolumn{1}{c}{Strong, $t_6$} \\
$r_{1,1}$ &  51.3 (39.8 $\mid$ 09.0) & 85.3 (09.5 $\mid$ 05.2) && 71.5 (21.0 $\mid$ 07.5) & 92.8 (02.5 $\mid$ 04.7) \\
          &  48.8 (20.0 $\mid$ 31.2)  & 83.0 (02.3 $\mid$ 14.7) && 00.0 (00.0 $\mid$ 100.0) & 00.0 (00.0 $\mid$ 100.0) \\
$r_{2,1}$ &  66.5 (10.0 $\mid$ 23.5) & 96.3 (01.0 $\mid$ 02.7) && 70.8 (01.2 $\mid$ 28.0) & 95.3 (00.0 $\mid$ 04.7) \\
          &  65.3 (09.8 $\mid$ 25.0) & 95.8 (01.0 $\mid$ 03.2) && 67.3 (01.2 $\mid$ 31.5) & 94.5 (00.3 $\mid$ 05.2) \\
$r_{2,2}$ &  70.0 (08.5 $\mid$ 21.5) & 98.0 (00.8 $\mid$ 01.2) && 69.5 (00.8 $\mid$ 29.7) & 94.8 (00.0 $\mid$ 05.2)  \\
          &  68.8 (08.0 $\mid$ 23.2) & 97.5 (00.8 $\mid$ 01.7) && 68.8 (00.2 $\mid$ 31.0) & 96.3 (00.0 $\mid$ 03.7) \\
\\
\multicolumn{1}{c}{Weak, $\cN$} \\
$r_{1,1}$ &  29.3 (68.5 $\mid$ 02.2) & 69.8 (29.2 $\mid$ 01.0) && 34.8 (64.2 $\mid$ 01.0) & 75.5 (22.0 $\mid$ 02.5) \\
          &  43.5 (34.8 $\mid$ 21.7) & 78.8 (05.5 $\mid$ 15.7) && 00.0 (00.0 $\mid$ 100.0) & 00.0 (00.0 $\mid$ 100.0) \\
$r_{2,1}$ &  72.5 (16.3 $\mid$ 11.2) & 97.5 (02.5 $\mid$ 00.0) && 80.8 (04.7 $\mid$ 14.5) & 99.0 (00.0 $\mid$ 01.0) \\
          &  71.8 (15.5 $\mid$ 12.7) & 97.0 (02.5 $\mid$ 00.5) && 78.0 (04.2 $\mid$ 17.8) & 98.8 (00.0 $\mid$ 01.2) \\
$r_{2,2}$ &  74.8 (13.0 $\mid$ 12.2) & 98.3 (01.5 $\mid$ 00.2) && 80.8 (03.5 $\mid$ 15.7) & 99.5 (00.0 $\mid$ 00.5)  \\
          &  73.8 (12.0 $\mid$ 14.2) & 98.3 (01.5 $\mid$ 00.2) && 82.0 (03.0 $\mid$ 15.0) & 98.8 (00.0 $\mid$ 01.2) \\
\hline %\bottomrule
\end{tabular}
\end{table}


Results of our estimator and those using $\xi_{m,k}=0$ are all included in \Cref{tab: simulation_setting_D}. First, the performance of our estimator remains satisfactory under $t_6$ noise, comparing with the results under Gaussian noise. It is also clear that our proposed estimator has higher correct proportion than the naive eigenvalue-ratio estimator (using $\xi_{m,k}=0$) in most setting. Exceptions include estimating $r_{1,1}$ for Setting~(D.1) with weak factors, where our estimator suffers more. This can be understood since the number of factors are more underestimated when factors are weak, while the naive estimator tends to overestimate the factor numbers and hence remedy this when cross-sectional dimensions are small. Notwithstanding, the naive estimator significantly suffers and always overestimates $r_{1,1}$ in Setting~(D.2). In comparison, our proposed estimator has improved accuracy when dimensionality increases, numerically reflecting the consistency result in \Cref{thm: eigenratio}.







\subsection{Real data analysis}
\label{subsec: real_data}

In this subsection, we use our proposed multi-order tensor factor model to analyse New York traffic data. We study a JXC constructed from a dataset which includes all individual taxi rides operated by Yellow Taxi within Manhattan Island of New York City. The data is available at:

{\small \texttt{https://www1.nyc.gov/site/tlc/about/tlc-trip-record-data.page}} .

To preclude the influence of Covid-19, we focus on trip records during the pre-Covid period from January 1, 2013 to December 31, 2019. For an overview of the dataset, it contains information on the pick-up and drop-off time and locations which are coded according to 69 predefined taxi zones on Manhattan Island. In particular, each day is divided into 24 hourly periods to represent the pick-up and drop-off time, with the first hourly period from 0 \text{a.m.} to 1 a.m, so that on day $t$ we have $\cY_t\in \R^{69\times 69\times 24}$, where $y_{i_1,i_2,i_3,t}$ is the number of trips from zone $i_1$ to zone $i_2$, with pick-up time within the $i_3$th hourly period.
We consider the business-day series which is 1,763 days long.




The taxi data examples in both \cite{Chenetal2022} and \cite{CenLam2025} suggest that Times Square zone is one of the centre zones leading the traffic dynamic in its local areas. To investigate this in detail, for each $t\in[1763]$, we form two threads: (1) $\cX_t^{(1)} \in\R^{24}$ denoting the hourly trip data within Times Square zone; and (2) $\cX_t^{(2)} \in \R^{18\times 18\times 24}$ as a sub-tensor of $\cY_t$ by selecting 18 neighbouring zones around Times Square. The map of selected taxi zones on Manhattan Island is included in the supplement.


The number of factors are estimated according to \Cref{subsec: est_num_fac}, where we adopt the iTIPUP estimator by \cite{Chenetal2022} to obtain $\wh{s}_{1,1} = 2$, $\wh{s}_{2,1} = 2$, $\wh{s}_{2,2} = 3$, and $\wh{s}_{2,3} = 2$.
\iffalse
\[
\wh{r_{1,1} +u_{1,1}} = 2, \; \wh{r_{2,1} +u_{2,1}} = 2, \; \wh{r_{2,2} +u_{2,2}} = 3, \; \wh{r_{2,3} +u_{2,3}} = 2 .
\]
\fi
The resulted global and hence local factor numbers are estimated as
\[
(\wh{r}_{1,1}, \wh{u}_{1,1}) =(1,1), \;
(\wh{r}_{2,1}, \wh{u}_{2,1}) =(1,1), \;
(\wh{r}_{2,2}, \wh{u}_{2,2}) =(1,2), \;
(\wh{r}_{2,3}, \wh{u}_{2,3}) =(1,1).
\]




\begin{table}%[htp!]
\caption{Estimated global and local factor loading matrices on the hour factor. Values are scaled by 30 and in bold if their magnitudes are larger than 10.}
\label{tab: taxi_biz_hour_loading}
\footnotesize
%\ra{1.5}
\setlength{\tabcolsep}{3.6pt}
\begin{center}
\begin{tabular}{l||ccccccccccccccccccccccccccc}
\hline $0_\text{am}$ & & 2 & & 4 & & 6 & & 8 & & 10 & & $12_\text{pm}$ & & 2 & & 4 & & 6 & & 8 & & 10 & & $12_\text{am}$ \\
\hline $\wh{\bA}_{1,1}$ & 5 & 3 & 2 & 2 & 1 & 1 & 2 & 5 & 6 & 6 & 6 & 6 & 7 & 7 & 7 & 6 & 6 & 6 & \textbf{10} & \textbf{12} & 7 & 6 & 6 & 6 \\
$\wh{\bB}_{1,1}$ & -6 & -5 & -3 & -2 & -1 & -1 & 3 & 5 & 2 & 0 & -1 & 1 & 3 & 3 & 1 & 1 & 0 & 5 & \textbf{11} & 3 & \textbf{-25} & -2 & 1 & -2 \\
\\
$\wh{\bA}_{2,3}$ & 3 & 2 & 1 & 1 & 1 & 1 & 3 & 6 & 8 & 9 & 8 & 8 & 8 & 8 & 8 & 7 & 6 & 6 & 7 & 8 & 7 & 6 & 6 & 5 \\
$\wh{\bB}_{2,3}$ & 2 & 1 & 1 & 1 & 0 & 1 & 2 & 5 & 8 & 8 & 8 & 4 & 8 & 8 & 8 & 9 & 8 & 8 & 9 & 9 & 7 & 5 & 4 & 3 \\
\hline
\end{tabular}
\end{center}
\end{table}



For the first thread, we use \Cref{tab: taxi_biz_hour_loading} to illustrate the estimated loading matrices on the hour factor. One key feature is the sheer difference in the global and local loading estimators for $\{\cX_t^{(1)}\}$: although they resonate at the 7--8 \text{p.m.} period, the local loading estimator contributes to the taxi traffic in Times Squares zone very mildly or even passively. This suggests that the flow in Times Squares can be largely explained by traffic in and out of Times Squares, which is reasonable since Times Squares can be seen as a core for tourism or offices, \text{cf.} Table~4 in \cite{Chenetal2022}. Such a centric position of Times Squares is also reflected by the peaking magnitudes in the table.




\begin{figure}%[htp!]
\begin{center}
\includegraphics[width=0.94\columnwidth]{pickup_loading.pdf}
\caption{Estimated global (left) and local (right) loading matrices on the pickup factor. Values are scaled by 10. Times Squares zone is in white.}
\label{Fig: taxi_biz_pickup_loading}
\end{center}
\begin{center}
\centerline{\includegraphics[width=0.96\columnwidth]{dropoff_loading.pdf}}
\caption{Estimated global (left) and local (middle and right) loading matrices on the drop-off factor. Values are scaled by 10. Times Squares zone is in white.}
\label{Fig: taxi_biz_dropoff_loading}
\end{center}
\end{figure}



Next, consider the second thread which represents the local areas around Times Squares zone. Similarly, we plot the estimated loading matrices on the pickup and drop-off factors in Figures~\ref{Fig: taxi_biz_pickup_loading} and \ref{Fig: taxi_biz_dropoff_loading}, respectively. The first feature is the different patterns between the pickup and drop-off factors. In detail, the pickup factors are generally stronger and centring at the Midtown areas and Union Square, whereas the drop-off factors have more intricate patterns, which is consistent with more drop-off factors, in our selected zones, found in Table~4 in \cite{Chenetal2022}.

Another finding lies at the right panel of \Cref{Fig: taxi_biz_dropoff_loading}, denoting the second local loading matrix estimated on the drop-off factor. Besides pinpointing again the importance of Union Square which is a transportation hub with many shops and restaurants, the active Kips Bay zone can be clearly identified, indicating the busy traffic within. It is also interesting that Kips Bay zone is almost silent in the estimated global loading on the pickup factor. Given that Kips Bay contains many medical facilities and schools, it is natural that there are very few people commuting from Times Squares to Kips Bay. Nevertheless, it is worth pointing out that this zone has been omitted in both \cite{Chenetal2022} and \cite{CenLam2025}, while our discovery takes advantage of the proposed model taking into account the dependency both globally and locally.












%%%%%%%%%%%%%%%%%%%%%%%%%%%%%%%%%%%%%%%%%%%%%%
%% Example with single Appendix:            %%
%%%%%%%%%%%%%%%%%%%%%%%%%%%%%%%%%%%%%%%%%%%%%%
%\begin{appendix}
%\section*{Title}\label{appn} %% if no title is needed, leave empty \section*{}.
%Appendices should be provided in \verb|{appendix}| environment, before Acknowledgements.

%If there is only one appendix, then please refer to it in text as \ldots\ in the \hyperref[appn]{Appendix}.
%\end{appendix}
%%%%%%%%%%%%%%%%%%%%%%%%%%%%%%%%%%%%%%%%%%%%%%
%% Example with multiple Appendixes:        %%
%%%%%%%%%%%%%%%%%%%%%%%%%%%%%%%%%%%%%%%%%%%%%%
\begin{appendix}

\section{Introduction of tensor map operator}
\label{sec: intro_map}

We extend in this subsection the reshape operator introduced in \cite{CenLam2025_KronProd} to a \textit{map} operator. Recall first that given any order-$K$ tensor $\cX\in \R^{I_1\times \dots \times I_K}$ and a set of ordered, strictly increasing elements $v=\{a_1,\dots, a_\ell\} \subseteq[K]$, the reshaped tensor $\reshape(\cX,v)$ is the order-$(K-\ell+1)$ tensor resulting from combining all modes indexed by $v$ into the last mode. Let $V=\{v_1,\dots, v_L\}$ be a set of ordered vectors where each vector $v_l$ is the parameter representing mode indices in a reshape operation and all $v_l$'s, $l\in[L]$, form a partition of $[K]$. Then we denote the mapped tensor $\map(\cX, V)$ as the order-$L$ tensor such that the $l$th mode has dimension $\prod_{j\in v_l} I_j$ and is constructed by combining entries of $\cX$ with modes in $v_l$. This $\map(\cdot,\cdot)$ operation is more general than permuting the array elements in $\cX$ as it also allows for stacking elements from different modes into one, and more importantly, preserves the factor structure in the data, by iteratively applying Theorem~1 in \cite{CenLam2025_KronProd}.

Hereafter, we refer to $V$ as \textit{channel}. Without loss of generality, let the data tensor $\cX$ be pre-processed such that, for $a<b$, elements in $v_a$ are larger than those in $v_b$. Then the procedure of mapping a tensor can be summarized in \Cref{alg: map}. Finally, by \cite{CenLam2025_KronProd}, the reshape operator can be inverted, meaning that one can recover $\cX$ from $\reshape(\cX, v)$ given the original dimensions of $\cX$, we may also recover $\cX$ from $\map(\cX, V)$.


\begin{algorithm}
\caption{Tensor map}\label{alg: map}
\begin{algorithmic}[1]
%\setstretch{1.25}
\State \textbf{Input:} Tensor $\cX$, channel $V=\{v_1,\dots, v_L\}$
\State \textbf{Initialize:} Set $\map(\cX, V) \gets \cX$
\For{$l \in [L]$}
\State Set $\map(\cX, V) \gets \reshape\big\{ \map(\cX, V), v_l \big\}$
\EndFor
\State \textbf{Output:} The mapped tensor $\map(\cX, V)$
\end{algorithmic}
\end{algorithm}



%\section{Title of the first appendix}\label{appA}
%If there are more than one appendix, then please refer to it as \ldots\ in Appendix \ref{appA}, Appendix \ref{appB}, etc.

%\section{Title of the second appendix}\label{appB}
%\subsection{First subsection of Appendix \protect\ref{appB}}

%Use the standard \LaTeX\ commands for headings in \verb|{appendix}|.
%Headings and other objects will be numbered automatically.
%\begin{equation}
%\mathcal{P}=(j_{k,1},j_{k,2},\dots,j_{k,m(k)}). \label{path}
%\end{equation}

%Sample of cross-reference to the formula (\ref{path}) in Appendix \ref{appB}.
\end{appendix}

%%%%%%%%%%%%%%%%%%%%%%%%%%%%%%%%%%%%%%%%%%%%%%
%% Support information, if any,             %%
%% should be provided in the                %%
%% Acknowledgements section.                %%
%%%%%%%%%%%%%%%%%%%%%%%%%%%%%%%%%%%%%%%%%%%%%%
%\begin{acks}[Acknowledgments]
%The authors would like to thank the anonymous referees, an Associate Editor and the Editor for their constructive comments that improved the quality of this paper.
%\end{acks}

%%%%%%%%%%%%%%%%%%%%%%%%%%%%%%%%%%%%%%%%%%%%%%
%% Funding information, if any,             %%
%% should be provided in the                %%
%% funding section.                         %%
%%%%%%%%%%%%%%%%%%%%%%%%%%%%%%%%%%%%%%%%%%%%%%
%\begin{funding}
%The first author was supported by NSF Grant DMS-??-??????.
%\end{funding}

%%%%%%%%%%%%%%%%%%%%%%%%%%%%%%%%%%%%%%%%%%%%%%
%% Supplementary Material, including data   %%
%% sets and code, should be provided in     %%
%% {supplement} environment with title      %%
%% and short description. It cannot be      %%
%% available exclusively as external link.  %%
%% All Supplementary Material must be       %%
%% available to the reader on Project       %%
%% Euclid with the published article.       %%
%%%%%%%%%%%%%%%%%%%%%%%%%%%%%%%%%%%%%%%%%%%%%%
\begin{supplement}
\stitle{Supplement to ``Identification and estimation of multi-order tensor factor models''}
\sdescription{The online supplement includes additional details on the traffic dataset in \Cref{subsec: real_data}, in addition to the proof of all theoretical results and auxiliary results.}
\end{supplement}
%\begin{supplement}
%\stitle{Title of Supplement B}
%\sdescription{Short description of Supplement B.}
%\end{supplement}

%%%%%%%%%%%%%%%%%%%%%%%%%%%%%%%%%%%%%%%%%%%%%%%%%%%%%%%%%%%%%
%%                  The Bibliography                       %%
%%                                                         %%
%%  imsart-???.bst  will be used to                        %%
%%  create a .BBL file for submission.                     %%
%%                                                         %%
%%  Note that the displayed Bibliography will not          %%
%%  necessarily be rendered by Latex exactly as specified  %%
%%  in the online Instructions for Authors.                %%
%%                                                         %%
%%  MR numbers will be added by VTeX.                      %%
%%                                                         %%
%%  Use \cite{...} to cite references in text.             %%
%%                                                         %%
%%%%%%%%%%%%%%%%%%%%%%%%%%%%%%%%%%%%%%%%%%%%%%%%%%%%%%%%%%%%%

%% if your bibliography is in bibtex format, uncomment commands:
\bibliographystyle{imsart-nameyear} % Style BST file (imsart-number.bst or imsart-nameyear.bst)
\bibliography{reference}       % Bibliography file (usually '*.bib')

%% or include bibliography directly:
%\begin{thebibliography}{4}
%%
%\bibitem{r1}
%\textsc{Billingsley, P.} (1999). \textit{Convergence of
%Probability Measures}, 2nd ed.
%Wiley, New York.

%\bibitem{r2}
%\textsc{Bourbaki, N.}  (1966). \textit{General Topology}  \textbf{1}.
%Addison--Wesley, Reading, MA.

%\bibitem{r3}
%\textsc{Ethier, S. N.} and \textsc{Kurtz, T. G.} (1985).
%\textit{Markov Processes: Characterization and Convergence}.
%Wiley, New York.

%\bibitem{r4}
%\textsc{Prokhorov, Yu.} (1956).
%Convergence of random processes and limit theorems in probability
%theory. \textit{Theory  Probab.  Appl.}
%\textbf{1} 157--214.
%\end{thebibliography}




\clearpage
\setcounter{section}{0}
\setcounter{equation}{0}
%\appendix
\renewcommand{\thesection}{S.\arabic{section}}
\renewcommand{\theHsection}{S.\arabic{section}} 
\renewcommand{\thesubsection}{S\arabic{section}.\arabic{subsection}}
\renewcommand{\theHsubsection}{S\arabic{section}.\arabic{subsection}} 

\def\theequation{S.\arabic{equation}}
\def\thelemma{S.\arabic{lemma}}
\def\theremark{S.\arabic{remark}}


\setcounter{page}{1}

\title{Supplement to ``Identification and estimation of multi-order tensor factor models''}
\vspace{12pt}

%\tableofcontents \newpage

In this supplementary material, we provide additional details on the traffic data in the main text.
%additional mathematical details for Section 2.2 in the main text,
%extra simulation results,
%and additional explanations for the real data analysis.
Proofs of the main results stated in the main text are also included, together with auxiliary results and their proofs.

\section{Additional details}







\subsection{New York traffic data: visualisation}

\Cref{Fig: manhattan} displays the map of 69 taxi zones on Manhattan Island with the selected zones coloured.

\begin{figure}[htp!]
\begin{center}
\centerline{\includegraphics[width=.95\columnwidth]{manhattan_zone.pdf}}
\caption{Taxi zones on Manhattan Island. Times Square zone is in orange, 18 neighbouring zones are in different colours, and all remaining zones are shaded.}
\label{Fig: manhattan}
\end{center}
\end{figure}



\section{Proof of theorems and auxiliary results}



\subsection{Lemmas with their proofs}


\begin{lemma}\label{lemma: rate_Omega}
Let Assumptions~\ref{ass: identification}, \ref{ass: loadings}, \ref{ass: noise}, and \ref{ass: dependence_moment} hold. For any $m\in[M]$, $k\in[K_m]$, $n\neq m$, we have
\begin{align*}
    &\;\quad
    \sum_{\bi \in [p_{n,1}] \times \cdots \times[p_{n,K_n}]} \Big\| \wh\bOmega_{k,n,\bi}^{(m)} - \bOmega_{k,n,\bi}^{(m)} \Big\|_F^2 
    = \cO_P\Big( \frac{p_m p_n}{T} \Big) .
\end{align*}
\end{lemma}



\begin{proof}[Proof of Lemma~\ref{lemma: rate_Omega}]
From \eqref{eqn: unfold_Xt} and \eqref{eqn: def_hat_Sigma_mk}, we may further decompose
\begin{align}
    &\;\quad
    \wh\bOmega_{k,n,\bi}^{(m)} = \frac{1}{T} \sum_{t=1}^T \mat_k(\cX_t^{(m)}) \cdot (\cX_t^{(n)})_{\bi} \notag \\
    &=
    \frac{1}{T} \sum_{t=1}^T \Big\{ \bA_{m,k} \mat_k(\cG_t^{(m)}) \bA_{m,\text{-}k}^\trans + \bB_{m,k} \mat_k(\cF_t^{(m)}) \bB_{m,\text{-}k}^\trans + \mat_k(\cE_t^{(m)}) \Big\} \cdot (\cX_t^{(n)})_{\bi} \notag \\
    &=
    \frac{1}{T} \sum_{t=1}^T \bA_{m,k} \mat_k(\cG_t^{(m)}) \Big( \cG_t^{(n)} \times_{j=1}^{K_n} \bA_{n,j} \Big)_{\bi} \bA_{m,\text{-}k}^\trans \notag \\
    &\;\quad
    + \frac{1}{T} \sum_{t=1}^T \bA_{m,k} \mat_k(\cG_t^{(m)}) \Big( \cE_t^{(n)} \Big)_{\bi} \bA_{m,\text{-}k}^\trans \notag \\
    &\;\quad
    + \frac{1}{T} \sum_{t=1}^T \bA_{m,k} \mat_k(\cG_t^{(m)}) \Big( \cF_t^{(n)} \times_{\ell=1}^{K_n} \bB_{n,\ell} \Big)_{\bi} \bA_{m,\text{-}k}^\trans \notag \\
    &\;\quad
    + \frac{1}{T} \sum_{t=1}^T \bB_{m,k} \mat_k(\cF_t^{(m)}) \Big( \cE_t^{(n)} \Big)_{\bi} \bB_{m,\text{-}k}^\trans \notag \\
    &\;\quad
    + \frac{1}{T} \sum_{t=1}^T \bB_{m,k} \mat_k(\cF_t^{(m)}) \Big( \cG_t^{(n)} \times_{j=1}^{K_n} \bA_{n,j} \Big)_{\bi} \bB_{m,\text{-}k}^\trans \notag \\
    &\;\quad
    + \frac{1}{T} \sum_{t=1}^T \mat_k(\cE_t^{(m)}) \Big( \cG_t^{(n)} \times_{j=1}^{K_n} \bA_{n,j} \Big)_{\bi} \notag \\
    &\;\quad
    + \frac{1}{T} \sum_{t=1}^T \bB_{m,k} \mat_k(\cF_t^{(m)}) \Big( \cF_t^{(n)} \times_{\ell=1}^{K_n} \bB_{n,\ell} \Big)_{\bi} \bB_{m,\text{-}k}^\trans \notag \\
    &\;\quad
    + \frac{1}{T} \sum_{t=1}^T \mat_k(\cE_t^{(m)}) \Big( \cF_t^{(n)} \times_{\ell=1}^{K_n} \bB_{n,\ell} \Big)_{\bi} \notag \\
    &\;\quad
    + \frac{1}{T} \sum_{t=1}^T \mat_k(\cE_t^{(m)}) \Big( \cE_t^{(n)} \Big)_{\bi} \notag \\
    &
    =: \cI_1 + \cI_2 + \cI_3 + \cI_4 + \cI_5 + \cI_6 + \cI_7 + \cI_8 + \cI_9 .
    \label{eqn: decomp_Omega}
\end{align}


Note that by Assumptions~\ref{ass: identification}(b), \ref{ass: noise}, and \ref{ass: dependence_moment}(a), $\cI_2$--$\cI_9$ all have zero mean. For $\cI_2$, consider $\sum_{\bi \in [p_{n,1}] \times \cdots \times[p_{n,K_n}]} \big\| T^{-1} \sum_{t=1}^T \mat_k(\cG_t^{(m)}) \big( \cE_t^{(n)} \big)_{\bi} \big\|_F^2$ first. Then we have
\begin{align*}
    \sum_{\bi \in \substack{[p_{n,1}] \times \cdots\\ \times[p_{n,K_n}] }} \Bigg\| \frac{1}{T} \sum_{t=1}^T \mat_k(\cG_t^{(m)}) \Big( \cE_t^{(n)} \Big)_{\bi} \Bigg\|_F^2 
    &= \Bigg\| \frac{1}{T} \sum_{t=1}^T \vec(\cG_t^{(m)}) \otimes \vec(\cE_t^{(n)}) \Bigg\|_F^2 \\
    &= \cO_P\Big( \frac{p_n}{T} \Big) ,
\end{align*}
where the last equality used \Cref{ass: dependence_moment}(b) and zero mean of $\cI_2$. Hence, by \Cref{ass: loadings}(b),
\begin{equation}
\label{eqn: cI2_rate}
\sum_{\bi \in [p_{n,1}] \times \cdots \times[p_{n,K_n}]} \|\cI_2\|_F^2 \leq \cO_P\Big( \frac{p_n}{T} \Big)\cdot \|\bA_{m,k}\|_F^2 \cdot \|\bA_{m,\text{-}k}\|_F^2 = \cO_P\Big( \frac{p_m p_n}{T} \Big) .
\end{equation}

For $\cI_3$--$\cI_9$, note that by \Cref{ass: loadings}(a), $\Big( \cG_t^{(n)} \times_{j=1}^{K_n} \bA_{n,j} \Big)_{\bi}$ and $\Big( \cF_t^{(n)} \times_{\ell=1}^{K_n} \bB_{n,\ell} \Big)_{\bi}$ are respectively finite linear combinations of elements in $\cG_t^{(n)}$ and $\cF_t^{(n)}$. Then by the similar arguments as the above steps for $\cI_2$, for $j=3, \dots, 9$, we have
\begin{equation}
\label{eqn: cI3_to_cI9_rate}
\sum_{\bi \in [p_{n,1}] \times \cdots \times[p_{n,K_n}]} \|\cI_j\|_F^2 = \cO_P\Big( \frac{p_m p_n}{T} \Big) .
\end{equation}


Finally, consider $\sum_{\bi \in [p_{n,1}] \times \cdots \times[p_{n,K_n}]} \|\cI_1 - \E(\cI_1)\|_F^2$. By \Cref{ass: loadings}(b), we have
\begin{align}
    &\;\quad
    \sum_{\bi \in [p_{n,1}] \times \cdots \times[p_{n,K_n}]} \|\cI_1 - \E(\cI_1)\|_F^2 \notag \\
    &\leq
    \|\bA_{m,k}\|_F^2 \cdot \|\bA_{m,\text{-}k}\|_F^2
    \sum_{\bi \in \substack{[p_{n,1}] \times \cdots\\ \times[p_{n,K_n}] }} \Bigg\| \frac{1}{T} \sum_{t=1}^T \Bigg[ \mat_k(\cG_t^{(m)}) \Big( \cG_t^{(n)} \times_{j=1}^{K_n} \bA_{n,j} \Big)_{\bi} \notag \\
    &\;\quad
    - \E\Big\{ \mat_k(\cG_t^{(m)}) \Big( \cG_t^{(n)} \times_{j=1}^{K_n} \bA_{n,j} \Big)_{\bi} \Big\} \Bigg] \Bigg\|_F^2 \notag \\
    &\lesssim
    p_m \cdot \sum_{\bi \in \substack{[p_{n,1}] \times \cdots\\ \times[p_{n,K_n}] }} \Bigg\| \frac{1}{T} \sum_{t=1}^T \Bigg[ \vec(\cG_t^{(m)}) \Big( \cG_t^{(n)} \times_{j=1}^{K_n} \bA_{n,j} \Big)_{\bi} \notag \\
    &\;\quad
    - \E\Big\{ \vec(\cG_t^{(m)}) \Big( \cG_t^{(n)} \times_{j=1}^{K_n} \bA_{n,j} \Big)_{\bi} \Big\} \Bigg] \Bigg\|_F^2 \notag \\
    &\lesssim
    p_m \cdot \sum_{\bi \in \substack{[p_{n,1}] \times \cdots\\ \times[p_{n,K_n}] }} \Bigg\| \frac{1}{T} \sum_{t=1}^T \Bigg[ \vec(\cG_t^{(m)}) \otimes \vec(\cG_t^{(n)}) - \E\Big\{ \vec(\cG_t^{(m)}) \otimes \vec(\cG_t^{(n)}) \Big\} \Bigg] \Bigg\|_F^2 \notag \\
    &=
    \cO_P\Big( \frac{p_m p_n}{T} \Big) ,
    \label{eqn: cI1_var}
\end{align}
where the last line used Assumptions~\ref{ass: loadings}(a) and \ref{ass: dependence_moment}(b).

Combining \eqref{eqn: decomp_Omega}, \eqref{eqn: cI2_rate}, \eqref{eqn: cI3_to_cI9_rate}, and \eqref{eqn: cI1_var}, we conclude that
\begin{align*}
    &\;\quad
    \sum_{\bi \in [p_{n,1}] \times \cdots \times[p_{n,K_n}]} \Big\| \wh\bOmega_{k,n,\bi}^{(m)} - \bOmega_{k,n,\bi}^{(m)} \Big\|_F^2 \\
    &\leq
    \sum_{\bi \in [p_{n,1}] \times \cdots \times[p_{n,K_n}]} \|\cI_1 - \E(\cI_1)\|_F^2 + \sum_{j=2}^9 \sum_{\bi \in [p_{n,1}] \times \cdots \times[p_{n,K_n}]} \|\cI_j\|_F^2 = \cO_P\Big( \frac{p_m p_n}{T} \Big) ,
\end{align*}
as desired. This completes the proof of the lemma.
\end{proof}






\begin{lemma}\label{lemma: rate_Sigma}
Let Assumptions~\ref{ass: identification}, \ref{ass: loadings}, \ref{ass: noise}, and \ref{ass: dependence_moment} hold. For any $m\in[M]$, $k\in[K_m]$, $n\neq m$, we have
\begin{align*}
    \Big\| \wh\bSigma_{m,k} -\bSigma_{m,k} \Big\|_F
    = \cO_P\Bigg\{ \frac{p_m}{\sqrt{T}} \Bigg( \sum_{\substack{n=1\\ n\neq m}}^M p_n \Bigg) \Bigg\} .
\end{align*}
\end{lemma}


\begin{proof}[Proof of Lemma~\ref{lemma: rate_Sigma}]
From \eqref{eqn: def_hat_Sigma_mk}, using the triangle inequality and Cauchy--Schwarz inequality,
\begin{align}
    &\;\quad
    \Big\| \wh\bSigma_{m,k} -\bSigma_{m,k} \Big\|_F = \Bigg\| \sum_{\substack{n=1\\ n\neq m}}^M \sum_{\bi \in [p_{n,1}] \times \cdots \times[p_{n,K_n}]} \Big( \wh\bOmega_{k,n,\bi}^{(m)} \wh\bOmega_{k,n,\bi}^{(m) \trans} - \bOmega_{k,n,\bi}^{(m)} \bOmega_{k,n,\bi}^{(m) \trans} \Big) \Bigg\|_F \notag \\
    &=
    \Bigg\| \sum_{\substack{n=1\\ n\neq m}}^M \sum_{\bi \in [p_{n,1}] \times \cdots \times[p_{n,K_n}]} \Bigg\{ \Big( \wh\bOmega_{k,n,\bi}^{(m)} - \bOmega_{k,n,\bi}^{(m)} \Big) \wh\bOmega_{k,n,\bi}^{(m) \trans} + \bOmega_{k,n,\bi}^{(m)} \Big( \wh\bOmega_{k,n,\bi}^{(m)} - \bOmega_{k,n,\bi}^{(m)} \Big)^\trans \Bigg\} \Bigg\|_F \notag \\
    &\leq
    \Bigg\| \sum_{\substack{n=1\\ n\neq m}}^M \sum_{\bi \in [p_{n,1}] \times \cdots \times[p_{n,K_n}]} \Big( \wh\bOmega_{k,n,\bi}^{(m)} - \bOmega_{k,n,\bi}^{(m)} \Big) \Big( \wh\bOmega_{k,n,\bi}^{(m)} - \bOmega_{k,n,\bi}^{(m)} \Big)^\trans \Bigg\|_F \notag \\
    &\;\quad
    + 2\Bigg\| \sum_{\substack{n=1\\ n\neq m}}^M \sum_{\bi \in [p_{n,1}] \times \cdots \times[p_{n,K_n}]} \bOmega_{k,n,\bi}^{(m)} \Big( \wh\bOmega_{k,n,\bi}^{(m)} - \bOmega_{k,n,\bi}^{(m)} \Big)^\trans \Bigg\|_F \notag \\
    &\lesssim
    \sum_{\substack{n=1\\ n\neq m}}^M \sum_{\bi \in \substack{[p_{n,1}] \times \cdots\\ \times[p_{n,K_n}] }} \Big\| \wh\bOmega_{k,n,\bi}^{(m)} - \bOmega_{k,n,\bi}^{(m)} \Big\|_F^2 \notag \\
    &\;\quad
    + \sum_{\substack{n=1\\ n\neq m}}^M \Bigg( \sum_{\bi \in \substack{[p_{n,1}] \times \cdots\\ \times[p_{n,K_n}] }} \Big\| \wh\bOmega_{k,n,\bi}^{(m)} - \bOmega_{k,n,\bi}^{(m)} \Big\|_F^2 \Bigg)^{1/2} \Bigg( \sum_{\bi \in \substack{[p_{n,1}] \times \cdots\\ \times[p_{n,K_n}] }} \Big\| \bOmega_{k,n,\bi}^{(m)} \Big\|_F^2 \Bigg)^{1/2} .
    \label{eqn: whSigma_minus_Sigma_rate}
\end{align}

By \Cref{ass: loadings}(b), we have
\begin{align}
    &\;\quad
    \sum_{\bi \in \substack{[p_{n,1}] \times \cdots\\ \times[p_{n,K_n}] }} \big\| \bOmega_{k,n,\bi}^{(m)} \big\|_F^2 \notag \\
    &=
    \sum_{\bi \in \substack{[p_{n,1}] \times \cdots\\ \times[p_{n,K_n}] }} \Big\| \bA_{m,k} \E\Big\{ \mat_k(\cG_t^{(m)}) \Big( \cG_t^{(n)} \times_{h=1}^{K_n} \bA_{n,h} \Big)_{\bi} \Big\} \bA_{m,\text{-}k}^\trans \Big\|_F^2 \notag \\
    &\leq
    \|\bA_{m,k}\|_F^2 \cdot \|\bA_{m,\text{-}k}\|_F^2 \cdot \Big\| \E\Big\{ \vec(\cG_t^{(m)}) \vec(\cG_t^{(n)})^\trans (\otimes_{h=1}^{K_n} \bA_{n,h})^\trans \Big\} \Big\|_F^2 \notag \\
    &\lesssim
    p_m \cdot \Big\| \E\Big\{ \vec(\cG_t^{(m)}) \otimes \vec(\cG_t^{(n)}) \Big\} \Big\|_F^2 \cdot \Big\| \otimes_{h=1}^{K_n} \bA_{n,h} \Big\|_F^2
    = \cO(p_m p_n) ,
    \label{eqn: sum_i_Omega_rate}
\end{align}
where the last equality used \Cref{ass: core_factor}(c). Combining \eqref{eqn: whSigma_minus_Sigma_rate}, \eqref{eqn: sum_i_Omega_rate}, and \Cref{lemma: rate_Omega}, we concludes the proof of this lemma.
\end{proof}




\begin{lemma}\label{lemma: eigenvalue_Sigma}
Let Assumptions~\ref{ass: identification}, \ref{ass: loadings}, \ref{ass: core_factor}, \ref{ass: noise}, and \ref{ass: dependence_moment} hold. For any $m\in[M]$, $k\in[K_m]$,
%there exists a positive constant $c$ such that
\begin{align*}
    %\lambda_{r_{m,k}}\big( \bSigma_{m,k} \big) \geq c\cdot p_m \Bigg( \sum_{\substack{n=1\\ n\neq m}}^M p_n \Bigg) .
    \lambda_{r_1}\big( \bSigma_{m,k} \big) \asymp \dots \asymp \lambda_{r_{m,k}}\big( \bSigma_{m,k} \big) \asymp p_m \Bigg( \sum_{\substack{n=1\\ n\neq m}}^M p_n \Bigg) .
\end{align*}
\end{lemma}


\begin{proof}[Proof of Lemma~\ref{lemma: eigenvalue_Sigma}]
Using Assumption~\ref{ass: loadings}(b), we have
\begin{align}
    &\;\quad
    \lambda_{r_{m,k}}\big( \bSigma_{m,k} \big) = \lambda_{r_{m,k}}\Bigg( \sum_{\substack{n=1\\ n\neq m}}^M \sum_{\bi \in [p_{n,1}] \times \cdots \times[p_{n,K_n}]} \bOmega_{k,n,\bi}^{(m)} \bOmega_{k,n,\bi}^{(m) \trans} \Bigg) \notag \\
    &=
    \lambda_{r_{m,k}}\Bigg( \bA_{m,k} \sum_{\substack{n=1\\ n\neq m}}^M \sum_{\bi \in \substack{[p_{n,1}] \times \cdots\\ \times[p_{n,K_n}] }}
    \E\Big\{ \mat_k(\cG_t^{(m)}) \Big( \cG_t^{(n)} \times_{h=1}^{K_n} \bA_{n,h} \Big)_{\bi} \Big\} \bA_{m,\text{-}k}^\trans \bA_{m,\text{-}k} \notag \\
    &\;\quad
    \cdot \E\Big\{ \mat_k(\cG_t^{(m)}) \Big( \cG_t^{(n)} \times_{h=1}^{K_n} \bA_{n,h} \Big)_{\bi} \Big\}^\trans
    \bA_{m,k}^\trans \Bigg) \notag \\
    &\geq
    \lambda_{r_{m,k}}\big( \bA_{m,k} \bA_{m,k}^\trans \big) \cdot \lambda_{r_{m,k}}\Bigg( \sum_{\substack{n=1\\ n\neq m}}^M \sum_{\bi \in \substack{[p_{n,1}] \times \cdots\\ \times[p_{n,K_n}] }} \E\Big\{ \mat_k(\cG_t^{(m)}) \Big( \cG_t^{(n)} \times_{h=1}^{K_n} \bA_{n,h} \Big)_{\bi} \Big\} \notag \\
    &\;\quad
    \cdot \bA_{m,\text{-}k}^\trans \bA_{m,\text{-}k} \E\Big\{ \mat_k(\cG_t^{(m)}) \Big( \cG_t^{(n)} \times_{h=1}^{K_n} \bA_{n,h} \Big)_{\bi} \Big\}^\trans
    \Bigg) \notag \\
    &\geq
    \Big\{ \lambda_{r_{m,k}}\big( p_{m,k} \bI \big) - \lambda_1\big( p_{m,k} \bI - \bA_{m,k} \bA_{m,k}^\trans \big) \Big\} \notag \\
    &\;\quad
    \cdot \Bigg\{ \lambda_{r_{m,k}}\Bigg( \sum_{\substack{n=1\\ n\neq m}}^M \sum_{\bi \in \substack{[p_{n,1}] \times \cdots\\ \times[p_{n,K_n}] }} \E\Big\{ \mat_k(\cG_t^{(m)}) \Big( \cG_t^{(n)} \times_{h=1}^{K_n} \bA_{n,h} \Big)_{\bi} \Big\} \notag \\
    &\;\quad
    \cdot p_{m,\text{-}k} \bI \cdot \E\Big\{ \mat_k(\cG_t^{(m)}) \Big( \cG_t^{(n)} \times_{h=1}^{K_n} \bA_{n,h} \Big)_{\bi} \Big\}^\trans \Bigg) \notag \\
    &\;\quad
    - \lambda_{1}\Bigg( \sum_{\substack{n=1\\ n\neq m}}^M \sum_{\bi \in \substack{[p_{n,1}] \times \cdots\\ \times[p_{n,K_n}] }} \E\Big\{ \mat_k(\cG_t^{(m)}) \Big( \cG_t^{(n)} \times_{h=1}^{K_n} \bA_{n,h} \Big)_{\bi} \Big\} \notag \\
    &\;\quad
    \cdot \big( p_{m,\text{-}k} \bI - \bA_{m,\text{-}k}^\trans \bA_{m,\text{-}k} \big) \cdot \E\Big\{ \mat_k(\cG_t^{(m)}) \Big( \cG_t^{(n)} \times_{h=1}^{K_n} \bA_{n,h} \Big)_{\bi} \Big\}^\trans \Bigg) \Bigg\} \notag \\
    &\asymp
    p_{m,k} \cdot p_{m,\text{-}k} \cdot \lambda_{r_{m,k}}\Bigg( \sum_{\substack{n=1\\ n\neq m}}^M \sum_{\bi \in \substack{[p_{n,1}] \times \cdots\\ \times[p_{n,K_n}] }} \E\Big\{ \mat_k(\cG_t^{(m)}) \Big( \cG_t^{(n)} \times_{h=1}^{K_n} \bA_{n,h} \Big)_{\bi} \Big\} \notag \\
    &\;\quad
    \cdot \E\Big\{ \mat_k(\cG_t^{(m)}) \Big( \cG_t^{(n)} \times_{h=1}^{K_n} \bA_{n,h} \Big)_{\bi} \Big\}^\trans \Bigg) \notag \\
    &\geq
    p_{m,k} \cdot p_{m,\text{-}k} \cdot \sum_{\substack{n=1\\ n\neq m}}^M \lambda_{r_{m,k}}\Bigg( \sum_{\bi \in \substack{[p_{n,1}] \times \cdots\\ \times[p_{n,K_n}] }} \E\Big\{ \mat_k(\cG_t^{(m)}) \Big( \cG_t^{(n)} \times_{h=1}^{K_n} \bA_{n,h} \Big)_{\bi} \Big\} \notag \\
    &\;\quad
    \cdot \E\Big\{ \mat_k(\cG_t^{(m)}) \Big( \cG_t^{(n)} \times_{h=1}^{K_n} \bA_{n,h} \Big)_{\bi} \Big\}^\trans \Bigg) . \label{eqn: eigenvalue_Sigma_step}
\end{align}
where the first and second inequalities used Weyl's inequality and the fact that for any compatible matrices $\bX,\bY$ with $\bX\bY$ being positive semi-definite ($\bSigma_{m,k}$ is $\bX\bY$ here), then $\lambda_i(\bX\bY)=\lambda_i(\bY\bX)$; and the asymptotic equality (in probability) used \Cref{ass: loadings}(b) and \Cref{lemma: inequality_sandwich}(a); and the last inequality used Weyl's inequality again.


Next, consider the following for any $n\neq m$.
First, for any $\bi\in [p_{n,1}] \times \cdots \times[p_{n,K_n}]$, there exists some $i\in[p_n]$ such that we can write
\begin{align*}
    &\;\quad
    \E\Big\{ \mat_k(\cG_t^{(m)}) \Big( \cG_t^{(n)} \times_{h=1}^{K_n} \bA_{n,h} \Big)_{\bi} \Big\} \\
    &=
    \E\Big[ \mat_k(\cG_t^{(m)}) \otimes \Big\{ \be_{p_n,i}^\trans \big( \otimes_{h=1}^{K_n} \bA_{n,h} \big) \vec(\cG_t^{(n)}) \Big\} \Big] \\
    &=
    \E\Big[ \mat_k(\cG_t^{(m)}) \otimes \Big\{ \vec(\cG_t^{(n)})^\trans \big( \otimes_{h=1}^{K_n} \bA_{n,h} \big)^\trans \be_{p_n,i} \Big\} \Big] ,
\end{align*}
where $\be_{p_n,i} \in\R^{p_n}$ contains $1$ in the $i$th entry and zero otherwise. Then
\begin{align}
    &\;\quad
    \lambda_{r_{m,k}}\Bigg( \sum_{\bi \in \substack{[p_{n,1}] \times \cdots\\ \times[p_{n,K_n}] }} \E\Big\{ \mat_k(\cG_t^{(m)}) \Big( \cG_t^{(n)} \times_{h=1}^{K_n} \bA_{n,h} \Big)_{\bi} \Big\} \notag \\
    &\;\quad
    \cdot \E\Big\{ \mat_k(\cG_t^{(m)}) \Big( \cG_t^{(n)} \times_{h=1}^{K_n} \bA_{n,h} \Big)_{\bi} \Big\}^\trans \Bigg) \notag \\
    &=
    \lambda_{r_{m,k}}\Bigg( \sum_{i \in [p_n]} \E\Big[ \mat_k(\cG_t^{(m)}) \otimes \Big\{ \vec(\cG_t^{(n)})^\trans \big( \otimes_{h=1}^{K_n} \bA_{n,h} \big)^\trans \be_{p_n,i} \Big\} \Big] \notag \\
    &\;\quad
    \cdot \E\Big[ \mat_k(\cG_t^{(m)}) \otimes \Big\{ \vec(\cG_t^{(n)})^\trans \big( \otimes_{h=1}^{K_n} \bA_{n,h} \big)^\trans \be_{p_n,i} \Big\} \Big]^\trans \Bigg) \notag \\
    &=
    \lambda_{r_{m,k}}\Bigg( \E\Big[ \mat_k(\cG_t^{(m)}) \otimes \Big\{ \vec(\cG_t^{(n)})^\trans \big( \otimes_{h=1}^{K_n} \bA_{n,h} \big)^\trans \Big\} \Big] \notag \\
    &\;\quad
    \cdot \E\Big[ \mat_k(\cG_t^{(m)})^\trans \otimes \Big\{ \big( \otimes_{h=1}^{K_n} \bA_{n,h} \big) \vec(\cG_t^{(n)}) \Big\} \Big] \Bigg) 
    %& \text{\ZC{it seems to work here, directly assume the above}}
    \asymp
    p_n \cdot c ,
    \label{eqn: eigenvalue_Sigma_step_2}
\end{align}
where $c$ is some positive constant and the last line used Assumptions~\ref{ass: loadings}(b) and \ref{ass: core_factor}(c). Plugging the above back in \eqref{eqn: eigenvalue_Sigma_step}, we conclude there exists some constant $c$ such that
\[
\lambda_{r_{m,k}}\big( \bSigma_{m,k} \big) \geq c \cdot p_m \sum_{n\in[M] \setminus\{m\}} p_n .
\]

It remains to show the same rate as an upper bound of $\lambda_{1}\big( \bSigma_{m,k} \big)$, but this is analogous to all the above steps with signs of inequalities reversed. The proof of this lemma is thus complete.
\end{proof}





\begin{lemma}[Useful inequalities to bound sandwiched terms]\label{lemma: inequality_sandwich}
We have the following.
\begin{enumerate}[itemsep=0pt, label = (\alph*), left = 0pt]
    \item For any square matrix $\bM$ and sequence of matrices $\bA_t$ with compatible dimensions for $t\in[T]$, we have
    \[
    \left\| \sum_{t=1}^T \bA_t \bM \bA_t^\trans \right\|_F^2 \leq \|\bM\|^2 \cdot \left\| \sum_{t=1}^T \bA_t \bA_t^\trans \right\|_F^2 .
    \]
    Note that this inequality is tight, which can be seen by taking $\bM$ as the identity matrix.
    \item (Lemma~B.16 in \cite{Barigozzietal2025_robust}.) For any sequence of matrices $\bA_t=[\bA_{t,ij}]$, $\bB_t$ for $t\in[T]$, and any matrix $\bM$ with compatible dimensions, we have
    \[
    \left\| \sum_{t=1}^T \bA_t \bM \bB_t^\trans \right\|_F^2 \leq \|\bM\|_F^2 \cdot \sum_{i,j} \left\| \sum_{t=1}^T \bA_{t,ij} \bB_t \right\|_F^2 .
    \]
\end{enumerate}
\end{lemma}

\begin{proof}[Proof of Lemma~\ref{lemma: inequality_sandwich}]
Part~(b) is a direct quotation of Lemma~B.16 in \cite{Barigozzietal2025_robust}, so it suffices to show part~(a). Expanding the squared Frobenius norm, we have
\begin{equation}
\label{eqn: sum_AMA_decomp}
\begin{split}
    \left\| \sum_{t=1}^T \bA_t \bM \bA_t^\trans \right\|_F^2 &= \tr\left\{ \left( \sum_{t=1}^T \bA_t \bM \bA_t^\trans \right) \left( \sum_{s=1}^T \bA_s \bM^\trans \bA_s^\trans \right) \right\} \\
    &=
    \sum_{t=1}^T \sum_{s=1}^T \tr(\bA_t \bM \bA_t^\trans \bA_s \bM^\trans \bA_s^\trans ) \\
    &=
    \sum_{t=1}^T \sum_{s=1}^T \tr(\bA_s^\trans \bA_t \bM \bA_t^\trans \bA_s \bM^\trans )
    =: \sum_{t=1}^T \sum_{s=1}^T \tr(\bA_{st} \bM \bA_{st}^\trans \bM^\trans ) ,
\end{split}
\end{equation}
where the second line used the cyclic property of trace, and the definition $\bA_{st}= \bA_s^\trans \bA_t$. For any $t,s\in[T]$, we may read $\tr(\bA_{st} \bM \bA_{st}^\trans \bM^\trans)$ as a Frobenius inner product and apply the Cauchy--Schwarz inequality such that
\begin{equation}
\label{eqn: trace_AMAM}
\tr(\bA_{st} \bM \bA_{st}^\trans \bM^\trans) = \langle \bA_{st} \bM, \bM \bA_{st} \rangle_F \leq \|\bA_{st} \bM\|_F \cdot \|\bM \bA_{st}\|_F \leq \|\bM\|^2 \|\bA_{st}\|_F^2 ,
\end{equation}
where the last inequality used the fact that for any matrix products $\bA\bB$, we have $\|\bA\bB\|_F=\sqrt{\sum_{i=1}^n \|\bA_{i\cdot}^\top \bB\|^2} \leq \sqrt{\sum_{i=1}^n \|\bA_{i\cdot}^\top\|^2 \|\bB\|^2} = \|\bA\|_F \cdot \|\bB\|$.

Combining \eqref{eqn: sum_AMA_decomp} and \eqref{eqn: trace_AMAM}, we have
\begin{align*}
    &\quad
    \left\| \sum_{t=1}^T \bA_t \bM \bA_t^\trans \right\|_F^2 = \sum_{t=1}^T \sum_{s=1}^T \tr(\bA_{st} \bM \bA_{st}^\trans \bM^\trans) \\
    &\leq
    \|\bM\|^2 \cdot \sum_{t=1}^T \sum_{s=1}^T \|\bA_{st}\|_F^2
    = \|\bM\|^2 \sum_{t=1}^T \sum_{s=1}^T \tr(\bA_s^\trans \bA_t \bA_t^\trans \bA_s) \\
    &=
    \|\bM\|^2 \cdot \tr\left\{ \left( \sum_{t=1}^T \bA_t \bA_t^\trans \right) \left( \sum_{s=1}^T \bA_s \bA_s^\trans \right) \right\}
    = \|\bM\|^2 \left\| \sum_{t=1}^T \bA_t \bA_t^\trans \right\|_F^2 ,
\end{align*}
as desired. This completes the proof of this lemma.
\end{proof}



\begin{lemma}\label{lemma: A_perp_rate}
Let all assumptions in \Cref{thm: global_loading} hold. For any $m\in[M]$, $k\in[K_m]$, denote the orthogonal complement of $\bA_{m,k}$ as $\bA_{m,k,\perp}$ which is defined as an orthogonal basis spanning the complement of $\textnormal{col}(\bA_{m,k})$. Then
\begin{align*}
    \Big\| \wh{\bA}_{m,k,\perp} -\bA_{m,k,\perp} \Big\|_F = \cO_P\Big( \frac{1}{\sqrt{T}} \Big) &, \\
    \Big\| \wh{\bA}_{m,k,\perp} \wh{\bA}_{m,k,\perp}^\trans - \bA_{m,k,\perp} \bA_{m,k,\perp}^\trans \Big\|_F = \cO_P\Big( \frac{1}{\sqrt{T}} \Big) &.
\end{align*}
\end{lemma}

\begin{proof}[Proof of Lemma~\ref{lemma: A_perp_rate}]
By the definition of $\wh{\bA}_{m,k,\perp}$ and $\wh{\bA}_{m,k}$, we have
\begin{align*}
    &\;\quad
    p_{m,k} \bA_{m,k,\perp} \bA_{m,k,\perp}^\trans + \bA_{m,k} \bA_{m,k}^\trans = \Big( \sqrt{p_{m,k}} \bA_{m,k,\perp}, \bA_{m,k} \Big) \Big( \sqrt{p_{m,k}} \bA_{m,k,\perp}, \bA_{m,k} \Big)^\trans \\
    &=
    p_{m,k} \cdot \bI = \Big( \sqrt{p_{m,k}} \wh{\bA}_{m,k,\perp}, \wh{\bA}_{m,k} \Big) \Big( \sqrt{p_{m,k}} \wh{\bA}_{m,k,\perp}, \wh{\bA}_{m,k} \Big)^\trans \\
    &=
    p_{m,k} \wh{\bA}_{m,k,\perp} \wh{\bA}_{m,k,\perp}^\trans + \wh{\bA}_{m,k} \wh{\bA}_{m,k}^\trans ,
\end{align*}
which implies
\begin{equation}
\label{eqn: A_perp_minus_hat_A_perp}
\wh{\bA}_{m,k,\perp} \wh{\bA}_{m,k,\perp}^\trans = \bA_{m,k,\perp} \bA_{m,k,\perp}^\trans + \frac{1}{p_{m,k}} \big( \bA_{m,k} \bA_{m,k}^\trans - \wh{\bA}_{m,k} \wh{\bA}_{m,k}^\trans \big) .
\end{equation}

Hence applying Lemma~3 in \cite{Lametal2011} on \eqref{eqn: A_perp_minus_hat_A_perp}, we have
\begin{align*}
    &\;\quad
    \Big\| \wh{\bA}_{m,k,\perp} -\bA_{m,k,\perp} \Big\|_F \\
    &\leq
    \frac{8}{\lambda_{p_{m,k} -r_{m,k}}\big( \bA_{m,k,\perp} \bA_{m,k,\perp}^\trans \big)} \Big\| \frac{1}{p_{m,k}} \big( \bA_{m,k} \bA_{m,k}^\trans - \wh{\bA}_{m,k} \wh{\bA}_{m,k}^\trans \big) \Big\|_F \\
    &\lesssim
    \frac{1}{p_{m,k}} \Big\| \big( \bA_{m,k} - \wh{\bA}_{m,k} \big) \bA_{m,k}^\trans + \wh{\bA}_{m,k} \big( \bA_{m,k} - \wh{\bA}_{m,k} \big)^\trans \Big\|_F
    = \cO_P\Big( \frac{1}{\sqrt{T}} \Big) ,
\end{align*}
where the last equality used the triangle inequality and \Cref{thm: global_loading}. The second statement in the lemma is straightforward from the above, and this completes the proof of this lemma.
\end{proof}




\begin{lemma}\label{lemma: eigenvalue_Sigma_B}
Let Assumptions~\ref{ass: identification}, \ref{ass: loadings}, \ref{ass: core_factor}, \ref{ass: noise}, \ref{ass: dependence_moment}, and \ref{ass: dependence_thread} hold.
%For any $m\in[M]$, $k\in[K_m]$ with $K_m> 1$, let $\wh{\bD}_{m,k}$ be the diagonal matrix consisting of the first $u_{m,k}$ eigenvalues of $\wh\bSigma_{B,m,k}$ defined in \eqref{eqn: def_hat_Sigma_B_mk}.
Given any $m\in[M]$, $k\in[K_m]$ with $K_m> 1$, recall $\wh{\bD}_{m,k}$ from the statement of Theorem~\ref{thm: local_loading}.
Then the eigenvalues of $\wh{\bD}_{m,k}$ satisfy $\lambda_i(\wh{\bD}_{m,k}) = \lambda_i(\bSigma_{F,m,k}) + o_P(1)$ for any $i\leq u_{m,k}$.
\end{lemma}


\begin{proof}[Proof of Lemma~\ref{lemma: eigenvalue_Sigma_B}]
By the definition in \eqref{eqn: def_hat_Sigma_B_mk}, we can decompose
\begin{align}
    &\;\quad
    \wh\bSigma_{B,m,k} = \frac{1}{Tp_m} \sum_{t=1}^T \mat_k(\cX_t^{(m)}) \wh{\bA}_{m,\text{-}k,\perp} \wh{\bA}_{m,\text{-}k,\perp}^\trans \mat_k(\cX_t^{(m)})^\trans \notag \\
    &=
    \frac{1}{Tp_m} \sum_{t=1}^T \Big( \bB_{m,k} \mat_k(\cF_t^{(m)}) \bB_{m,\text{-}k}^\trans \wh{\bA}_{m,\text{-}k,\perp} + \mat_k(\cE_t^{(m)}) \wh{\bA}_{m,\text{-}k,\perp} \notag \\
    &\;\quad
    + \bA_{m,k} \mat_k(\cG_t^{(m)}) \bA_{m,\text{-}k}^\trans \wh{\bA}_{m,\text{-}k,\perp} \Big) 
    \cdot \Big( \bB_{m,k} \mat_k(\cF_t^{(m)}) \bB_{m,\text{-}k}^\trans \wh{\bA}_{m,\text{-}k,\perp} \notag \\
    &\;\quad
    + \mat_k(\cE_t^{(m)}) \wh{\bA}_{m,\text{-}k,\perp} + \bA_{m,k} \mat_k(\cG_t^{(m)}) \bA_{m,\text{-}k}^\trans \wh{\bA}_{m,\text{-}k,\perp} \Big)^\trans \notag \\
    &=
    \frac{1}{Tp_m} \sum_{t=1}^T \bB_{m,k} \mat_k(\cF_t^{(m)}) \bB_{m,\text{-}k}^\trans \wh{\bA}_{m,\text{-}k,\perp} \wh{\bA}_{m,\text{-}k,\perp}^\trans \bB_{m,\text{-}k} \mat_k(\cF_t^{(m)})^\trans \bB_{m,k}^\trans \notag \\
    &\;\quad
    + \frac{1}{Tp_m} \sum_{t=1}^T \bB_{m,k} \mat_k(\cF_t^{(m)}) \bB_{m,\text{-}k}^\trans \wh{\bA}_{m,\text{-}k,\perp} \wh{\bA}_{m,\text{-}k,\perp}^\trans \mat_k(\cE_t^{(m)})^\trans \notag \\
    &\;\quad
    + \frac{1}{Tp_m} \sum_{t=1}^T \bB_{m,k} \mat_k(\cF_t^{(m)}) \bB_{m,\text{-}k}^\trans \wh{\bA}_{m,\text{-}k,\perp} \wh{\bA}_{m,\text{-}k,\perp}^\trans \bA_{m,\text{-}k} \mat_k(\cG_t^{(m)})^\trans \bA_{m,k}^\trans \notag \\
    &\;\quad
    + \frac{1}{Tp_m} \sum_{t=1}^T \mat_k(\cE_t^{(m)}) \wh{\bA}_{m,\text{-}k,\perp} \wh{\bA}_{m,\text{-}k,\perp}^\trans \bB_{m,\text{-}k} \mat_k(\cF_t^{(m)})^\trans \bB_{m,k}^\trans \notag \\
    &\;\quad
    + \frac{1}{Tp_m} \sum_{t=1}^T \mat_k(\cE_t^{(m)}) \wh{\bA}_{m,\text{-}k,\perp} \wh{\bA}_{m,\text{-}k,\perp}^\trans \mat_k(\cE_t^{(m)})^\trans \notag \\
    &\;\quad
    + \frac{1}{Tp_m} \sum_{t=1}^T \mat_k(\cE_t^{(m)}) \wh{\bA}_{m,\text{-}k,\perp} \wh{\bA}_{m,\text{-}k,\perp}^\trans \bA_{m,\text{-}k} \mat_k(\cG_t^{(m)})^\trans \bA_{m,k}^\trans \notag \\
    &\;\quad
    + \frac{1}{Tp_m} \sum_{t=1}^T \bA_{m,k} \mat_k(\cG_t^{(m)}) \bA_{m,\text{-}k}^\trans \wh{\bA}_{m,\text{-}k,\perp} \wh{\bA}_{m,\text{-}k,\perp}^\trans \bB_{m,\text{-}k} \mat_k(\cF_t^{(m)})^\trans \bB_{m,k}^\trans \notag \\
    &\;\quad
    + \frac{1}{Tp_m} \sum_{t=1}^T \bA_{m,k} \mat_k(\cG_t^{(m)}) \bA_{m,\text{-}k}^\trans \wh{\bA}_{m,\text{-}k,\perp} \wh{\bA}_{m,\text{-}k,\perp}^\trans \mat_k(\cE_t^{(m)})^\trans \notag \\
    &\;\quad
    + \frac{1}{Tp_m} \sum_{t=1}^T \bA_{m,k} \mat_k(\cG_t^{(m)}) \bA_{m,\text{-}k}^\trans \wh{\bA}_{m,\text{-}k,\perp} \wh{\bA}_{m,\text{-}k,\perp}^\trans \bA_{m,\text{-}k} \mat_k(\cG_t^{(m)})^\trans \bA_{m,k}^\trans \notag \\
    &=:
    \cII_1 + \cII_2 + \cII_3 + \cII_4 + \cII_5 + \cII_6 + \cII_7 + \cII_8 + \cII_9 .
    \label{eqn: decomp_Sigma_B}
\end{align}


Consider $\cII_1$ first. By Lemmas~\ref{lemma: inequality_sandwich} and \ref{lemma: A_perp_rate}, we can write
\begin{align*}
    &\;\quad
    \cII_1 = \frac{1}{Tp_m} \sum_{t=1}^T \bB_{m,k} \mat_k(\cF_t^{(m)}) \bB_{m,\text{-}k}^\trans \wh{\bA}_{m,\text{-}k,\perp} \wh{\bA}_{m,\text{-}k,\perp}^\trans \bB_{m,\text{-}k} \mat_k(\cF_t^{(m)})^\trans \bB_{m,k}^\trans \\
    &=
    \frac{1}{Tp_m} \sum_{t=1}^T \bB_{m,k} \mat_k(\cF_t^{(m)}) \bB_{m,\text{-}k}^\trans \bA_{m,\text{-}k,\perp} \bA_{m,\text{-}k,\perp}^\trans \bB_{m,\text{-}k} \mat_k(\cF_t^{(m)})^\trans \bB_{m,k}^\trans \\
    &\;\quad
    + o_P\Big\{ \frac{1}{Tp_m} \sum_{t=1}^T \bB_{m,k} \mat_k(\cF_t^{(m)}) \bB_{m,\text{-}k}^\trans \bB_{m,\text{-}k} \mat_k(\cF_t^{(m)})^\trans \bB_{m,k}^\trans \Big\} \\
    &=: \cII_{\Sigma,1} + o_P(\cII_{\Sigma,2}) .
\end{align*}
By \Cref{ass: loadings}(b), for any $j\in[K_m]$, there exists a half-orthogonal matrix $\bC_{m,j}$ with dimension $p_{m,j}\times (p_{m,j} -r_{m,j} -u_{m,j})$ such that $\bB_{m,j}^\trans \bC_{m,j} =\zero$ and, without loss of generality, $\bA_{m,j,\perp} = \big( p_{m,j}^{-1/2} \bB_{m,j}, \bC_{m,j} \big) + o(1)$. Hence,
\begin{align}
    &\;\quad
    \bB_{m,j}^\trans \bA_{m,j,\perp} \bA_{m,j,\perp}^\trans \bB_{m,j} = \bB_{m,j}^\trans \big( p_{m,j}^{-1/2} \bB_{m,j}, \bC_{m,j} \big) \big( p_{m,j}^{-1/2} \bB_{m,j}, \bC_{m,j} \big)^\trans \bB_{m,j} \notag \\
    &=
    p_{m,j}\cdot \big\{ \big( \bI, \zero\big) \big( \bI, \zero\big)^\trans + o(1) \big\} = p_{m,j} \cdot \{\bI + o(1)\} ,
    \label{eqn: B_A_perp_A_perp_B}
\end{align}
which also implies that
\begin{equation}
\label{eqn: BAAB_converge}
\bB_{m,\text{-}k}^\trans \bA_{m,\text{-}k,\perp} \bA_{m,\text{-}k,\perp}^\trans \bB_{m,\text{-}k}
    = \otimes_{j\in[K_m] \setminus \{k\}} \bB_{m,j}^\trans \bA_{m,j,\perp} \bA_{m,j,\perp}^\trans \bB_{m,j}
    = p_{m,\text{-}k} \cdot \{ \bI+o(1)\} .
\end{equation}
Note that in $\cII_{\Sigma,2}$, we also have $\bB_{m,\text{-}k}^\trans \bB_{m,\text{-}k} = p_{m,\text{-}k} \cdot \{ \bI+o(1)\}$ by \Cref{ass: loadings}(b). Hence $\cII_{\Sigma}$ dominates $o_P(\cII_{\Sigma,2})$ and it remains to consider $\cII_{\Sigma,1}$ for $\cII_1$. To this end, using \Cref{lemma: inequality_sandwich}(a), the $o(1)$ term in \eqref{eqn: BAAB_converge} is dominated when plugged in $\cII_{\Sigma,1}$. We thus only consider
\[
\cII_1 \asymp \frac{1}{Tp_{m,k}} \sum_{t=1}^T \bB_{m,k} \mat_k(\cF_t^{(m)}) \mat_k(\cF_t^{(m)})^\trans \bB_{m,k}^\trans \xrightarrow{P} \frac{1}{p_{m,k}} \bB_{m,k} \bSigma_{F,m,k} \bB_{m,k}^\trans ,
\]
according to \Cref{ass: core_factor}(b). Hence with \Cref{ass: loadings}(b) again, we conclude $\lambda_i(\cII_1) = \lambda_i(\bSigma_{F,m,k}) + o_P(1)$ for all $i\leq u_{m,k}$, while $\lambda_i(\cII_1) =0$ for $i> u_{m,k}$ by $\rank(\cII_1)\leq u_{m,k}$.


Next, consider $\cII_2$. By definition, we have
\begin{align}
    &\;\quad
    \|\cII_2\|_F^2 = \Bigg\| \frac{1}{Tp_m} \sum_{t=1}^T \bB_{m,k} \mat_k(\cF_t^{(m)}) \bB_{m,\text{-}k}^\trans \wh{\bA}_{m,\text{-}k,\perp} \wh{\bA}_{m,\text{-}k,\perp}^\trans \mat_k(\cE_t^{(m)})^\trans \Bigg\|_F^2 \notag \\
    &\asymp
    \sum_{i=1}^{p_{m,k}} \Bigg\| \frac{1}{Tp_m} \sum_{j=1}^{r_{m,\text{-}k}} \sum_{t=1}^T \bB_{m,k} \mat_k(\cF_t^{(m)})_{\cdot j} (\bB_{m,\text{-}k}^\trans \bA_{m,\text{-}k,\perp} \bA_{m,\text{-}k,\perp}^\trans)_{j\cdot}^\trans \mat_k(\cE_t^{(m)})_{i\cdot} \Bigg\|_F^2 \notag \\
    &\leq
    \frac{1}{T p_m^2} \cdot \big\| \bB_{m,k} \big\|_F^2 \cdot p_{m,\text{-}k} \cdot \max_{j\in[r_{m,\text{-}k}]} \sum_{i=1}^{p_{m,k}} \Bigg\| \frac{1}{\sqrt{T}} \sum_{t=1}^T \mat_k(\cF_t^{(m)})_{\cdot j} \notag \\
    &\;\quad
    \cdot \Bigg\{ \be_i^\trans \mat_k(\cE_t^{(m)}) \Bigg( \frac{\bB_{m,\text{-}k}^\trans \bA_{m,\text{-}k,\perp} \bA_{m,\text{-}k,\perp}^\trans}{\sqrt{p_{m,\text{-}k}}} \Bigg)_{j\cdot} \Bigg\} \Bigg\|_F^2 \notag \\
    &\leq
    \frac{1}{T p_m^2} \cdot \big\| \bB_{m,k} \big\|_F^2 \cdot p_{m,\text{-}k} \cdot \cO_P(p_{m,k})
    = \cO_P\Big( \frac{1}{T p_{m,\text{-}k}} \Big) ,
    \label{eqn: cII_2}
\end{align}
where $\be_i \in\R^{p_{m,k}}$ contains $1$ in its $i$th entry and $0$ elsewhere, and the last line used Assumptions~\ref{ass: loadings}(b) and \ref{ass: dependence_thread}(a). The same arguments hold for $\cII_4$, i.e.,
\begin{align}
    &\;\quad
    \|\cII_4\|_F^2 = \cO_P\Big( \frac{1}{T p_{m,\text{-}k}} \Big) .
    \label{eqn: cII_4}
\end{align}

For $\cII_3$ (and hence $\cII_7$), we have the following by similar steps in \eqref{eqn: cII_2}, Assumptions~\ref{ass: loadings}(b) and \ref{ass: dependence_thread}(b):
\begin{align}
    &\;\quad
    \|\cII_3\|_F^2, \; \|\cII_7\|_F^2 = \cO_P\Big( \frac{1}{T p_{m,\text{-}k}} \Big) .
    \label{eqn: cII_3_7}
\end{align}


Consider $\cII_6$ now. Using \Cref{lemma: A_perp_rate}, we have
\begin{equation}
\label{eqn: hat_A_perp_minus_A_kron}
    \Big\| \wh{\bA}_{m,\text{-}k,\perp} -\bA_{m,\text{-}k,\perp} \Big\|_F = \cO_P\Big( \frac{1}{\sqrt{T}} \Big)
    = \Big\| \wh{\bA}_{m,\text{-}k,\perp} \wh{\bA}_{m,\text{-}k,\perp}^\trans - \bA_{m,\text{-}k,\perp} \bA_{m,\text{-}k,\perp}^\trans \Big\|_F^2 ,
\end{equation}
by a simple induction argument (see \text{e.g.} the induction argument in the proof of Lemma~6 in \cite{CenLam2025_KronProd}). Then, using \Cref{lemma: inequality_sandwich}(b),
\begin{align}
    &\;\quad
    \|\cII_6\|_F^2 = \Bigg\| \frac{1}{Tp_m} \sum_{t=1}^T \mat_k(\cE_t^{(m)}) \wh{\bA}_{m,\text{-}k,\perp} \wh{\bA}_{m,\text{-}k,\perp}^\trans \bA_{m,\text{-}k} \mat_k(\cG_t^{(m)})^\trans \bA_{m,k}^\trans \Bigg\|_F^2 \notag \\
    &\lesssim
    \Bigg\| \frac{1}{Tp_m} \sum_{t=1}^T \mat_k(\cE_t^{(m)}) \notag \\
    &\;\quad
    \cdot \big( \wh{\bA}_{m,\text{-}k,\perp} \wh{\bA}_{m,\text{-}k,\perp}^\trans - \bA_{m,\text{-}k,\perp} \bA_{m,\text{-}k,\perp}^\trans \big) \bA_{m,\text{-}k} \mat_k(\cG_t^{(m)})^\trans \bA_{m,k}^\trans \Bigg\|_F^2 \notag \\
    &\leq
    \frac{1}{p_m^2} \cdot \|\bA_{m,k}\|_F^2 \cdot \Big\| \wh{\bA}_{m,\text{-}k,\perp} \wh{\bA}_{m,\text{-}k,\perp}^\trans - \bA_{m,\text{-}k,\perp} \bA_{m,\text{-}k,\perp}^\trans \Big\|_F^2 \cdot p_{m,\text{-}k} \notag \\
    &\;\quad
    \cdot \sum_{i=1}^{p_{m,k}} \sum_{j=1}^{p_{m,\text{-}k}} \Bigg\| \frac{1}{T} \sum_{t=1}^T \mat_k(\cE_t^{(m)})_{i,j} \frac{\bA_{m,\text{-}k}}{\sqrt{p_{m,\text{-}k}}} \mat_k(\cG_t^{(m)})^\trans \Bigg\|_F^2 \notag \\
    &=
    \frac{1}{p_m^2} \cdot \|\bA_{m,k}\|_F^2 \cdot \frac{p_{m,\text{-}k}}{T} \cdot \cO_P(p_{m})
    = \cO_P\Big( \frac{1}{T} \Big) ,
    \label{eqn: cII_6}
\end{align}
where the second last equality used \eqref{eqn: hat_A_perp_minus_A_kron}, \Cref{ass: dependence_thread}(c) and similar arguments in \eqref{eqn: cII_2}, and the last used \Cref{ass: loadings}(b). Similarly, we also have
\begin{align}
    &\;\quad
    \|\cII_8\|_F^2 = \cO_P\Big( \frac{1}{T p_{m,\text{-}k}} \Big) ,
    \label{eqn: cII_8}
\end{align}


For $\cII_9$, by Lemmas~\ref{lemma: inequality_sandwich}(a) and \ref{lemma: A_perp_rate}, we have
\begin{align}
    &\;\quad
    \|\cII_9\|_F^2 \notag \\
    &=
    \Bigg\| \frac{1}{Tp_m} \sum_{t=1}^T \bA_{m,k} \mat_k(\cG_t^{(m)}) \bA_{m,\text{-}k}^\trans \wh{\bA}_{m,\text{-}k,\perp} \wh{\bA}_{m,\text{-}k,\perp}^\trans \bA_{m,\text{-}k} \mat_k(\cG_t^{(m)})^\trans \bA_{m,k}^\trans \Bigg\|_F^2 \notag \\
    &\leq
    \Big\| \wh{\bA}_{m,\text{-}k,\perp} \wh{\bA}_{m,\text{-}k,\perp}^\trans - \bA_{m,\text{-}k,\perp} \bA_{m,\text{-}k,\perp}^\trans \Big\|_F^2 \notag \\
    &\;\quad
    \cdot \Bigg\| \frac{1}{T p_{m,k}} \sum_{t=1}^T \bA_{m,k} \mat_k(\cG_t^{(m)}) \bA_{m,\text{-}k}^\trans \bA_{m,\text{-}k} \mat_k(\cG_t^{(m)})^\trans \bA_{m,k}^\trans \Bigg\|_F^2
    = \cO_P\Big( \frac{1}{T} \Big) ,
    \label{eqn: cII_9}
\end{align}
where the last equality used \Cref{ass: core_factor} and similar arguments on $\cII_1$. Similarly for $\cII_5$, using also Lemma~\ref{lemma: inequality_sandwich}(a) and noting that $\big\| \wh{\bA}_{m,\text{-}k,\perp} \wh{\bA}_{m,\text{-}k,\perp}^\trans -I \big\|_F =\cO(1)$,
\begin{align}
    &\;\quad
    \|\cII_5\|_F^2 = \Bigg\| \frac{1}{Tp_m} \sum_{t=1}^T \mat_k(\cE_t^{(m)}) \wh{\bA}_{m,\text{-}k,\perp} \wh{\bA}_{m,\text{-}k,\perp}^\trans \mat_k(\cE_t^{(m)})^\trans \Bigg\|_F^2 \notag \\
    &\lesssim
    \Bigg\| \frac{1}{Tp_m} \sum_{t=1}^T \mat_k(\cE_t^{(m)}) \mat_k(\cE_t^{(m)})^\trans \Bigg\|_F^2
    = \cO_P\Big( \frac{1}{T p_{m,\text{-}k}} + \frac{1}{p_{m,k}} \Big) ,
    \label{eqn: cII_5}
\end{align}
where the last equality used \Cref{ass: noise} and the proof of Lemma~3 in \cite{barigozzi2022statistical}.

Finally, combining all \eqref{eqn: cII_2}--\eqref{eqn: cII_5}, then the statement in the lemma is concluded by \eqref{eqn: decomp_Sigma_B} and applying Weyl's inequality iteratively. This ends the proof of the lemma.
\end{proof}




\begin{lemma}\label{lemma: eigenvalue_tilde_Sigma_B}
Let Assumptions~\ref{ass: identification}, \ref{ass: loadings}, \ref{ass: core_factor}, \ref{ass: noise}, \ref{ass: dependence_moment}, and \ref{ass: dependence_thread} hold.
%For any $m\in[M]$ with $K_m> 1$, let $\wt{\bD}_{m,1}$ be the diagonal matrix consisting of the first $u_{m,1}$ eigenvalues of $\wt\bSigma_{B,m,1}$ defined in \eqref{eqn: def_tilde_Sigma_B_m1}.
Given any $m\in[M]$ with $K_m= 1$, recall $\wt{\bD}_{m,1}$ from the statement of Theorem~\ref{thm: local_loading}.
Then the eigenvalues of $\wt{\bD}_{m,1}$ satisfy $\lambda_i(\wt{\bD}_{m,1}) = \lambda_i(\bSigma_{F,m,1}) + o_P(1)$ for any $i\leq u_{m,1}$.
\end{lemma}


\begin{proof}[Proof of Lemma~\ref{lemma: eigenvalue_tilde_Sigma_B}]
By the definition in \eqref{eqn: def_tilde_Sigma_B_m1}, we decompose
\begin{align}
    &\;\quad
    \wt\bSigma_{B,m,1} = \frac{1}{Tp_{m,1}} \sum_{t=1}^T \wh{\bA}_{m,1,\perp} \wh{\bA}_{m,1,\perp}^\trans \mat_1(\cX_t^{(m)}) \mat_1(\cX_t^{(m)})^\trans \wh{\bA}_{m,1,\perp} \wh{\bA}_{m,1,\perp}^\trans \notag \\
    &=
    \frac{1}{Tp_{m,1}} \sum_{t=1}^T \wh{\bA}_{m,1,\perp} \wh{\bA}_{m,1,\perp}^\trans \Big\{ \bA_{m,1} \mat_1(\cG_t^{(m)}) + \bB_{m,1} \mat_1(\cF_t^{(m)}) + \mat_k(\cE_t^{(m)}) \Big\} \notag \\
    &\;\quad
    \cdot \Big\{ \bA_{m,1} \mat_1(\cG_t^{(m)}) + \bB_{m,1} \mat_1(\cF_t^{(m)}) + \mat_k(\cE_t^{(m)}) \Big\}^\trans \wh{\bA}_{m,1,\perp} \wh{\bA}_{m,1,\perp}^\trans \notag \\
    &=
    \frac{1}{Tp_{m,1}} \sum_{t=1}^T \wh{\bA}_{m,1,\perp} \wh{\bA}_{m,1,\perp}^\trans \bB_{m,1} \mat_1(\cF_t^{(m)}) \mat_1(\cF_t^{(m)})^\trans \bB_{m,1}^\trans \wh{\bA}_{m,1,\perp} \wh{\bA}_{m,1,\perp}^\trans \notag \\
    &\;\quad
    + \frac{1}{Tp_{m,1}} \sum_{t=1}^T \wh{\bA}_{m,1,\perp} \wh{\bA}_{m,1,\perp}^\trans \bB_{m,1} \mat_1(\cF_t^{(m)}) \mat_1(\cG_t^{(m)})^\trans \bA_{m,1}^\trans \wh{\bA}_{m,1,\perp} \wh{\bA}_{m,1,\perp}^\trans \notag \\
    &\;\quad
    + \frac{1}{Tp_{m,1}} \sum_{t=1}^T \wh{\bA}_{m,1,\perp} \wh{\bA}_{m,1,\perp}^\trans \bB_{m,1} \mat_1(\cF_t^{(m)}) \mat_1(\cE_t^{(m)})^\trans \wh{\bA}_{m,1,\perp} \wh{\bA}_{m,1,\perp}^\trans \notag \\
    &\;\quad
    + \frac{1}{Tp_{m,1}} \sum_{t=1}^T \wh{\bA}_{m,1,\perp} \wh{\bA}_{m,1,\perp}^\trans \bA_{m,1} \mat_1(\cG_t^{(m)}) \mat_1(\cG_t^{(m)})^\trans \bA_{m,1}^\trans \wh{\bA}_{m,1,\perp} \wh{\bA}_{m,1,\perp}^\trans \notag \\
    &\;\quad
    + \frac{1}{Tp_{m,1}} \sum_{t=1}^T \wh{\bA}_{m,1,\perp} \wh{\bA}_{m,1,\perp}^\trans \bA_{m,1} \mat_1(\cG_t^{(m)}) \mat_1(\cF_t^{(m)})^\trans \bB_{m,1}^\trans \wh{\bA}_{m,1,\perp} \wh{\bA}_{m,1,\perp}^\trans \notag \\
    &\;\quad
    + \frac{1}{Tp_{m,1}} \sum_{t=1}^T \wh{\bA}_{m,1,\perp} \wh{\bA}_{m,1,\perp}^\trans \bA_{m,1} \mat_1(\cG_t^{(m)}) \mat_1(\cE_t^{(m)})^\trans \wh{\bA}_{m,1,\perp} \wh{\bA}_{m,1,\perp}^\trans \notag \\
    &\;\quad
    + \frac{1}{Tp_{m,1}} \sum_{t=1}^T \wh{\bA}_{m,1,\perp} \wh{\bA}_{m,1,\perp}^\trans \mat_1(\cE_t^{(m)}) \mat_1(\cG_t^{(m)})^\trans \bA_{m,1}^\trans \wh{\bA}_{m,1,\perp} \wh{\bA}_{m,1,\perp}^\trans \notag \\
    &\;\quad
    + \frac{1}{Tp_{m,1}} \sum_{t=1}^T \wh{\bA}_{m,1,\perp} \wh{\bA}_{m,1,\perp}^\trans \mat_1(\cE_t^{(m)}) \mat_1(\cF_t^{(m)})^\trans \bB_{m,1}^\trans \wh{\bA}_{m,1,\perp} \wh{\bA}_{m,1,\perp}^\trans \notag \\
    &\;\quad
    + \frac{1}{Tp_{m,1}} \sum_{t=1}^T \wh{\bA}_{m,1,\perp} \wh{\bA}_{m,1,\perp}^\trans \mat_1(\cE_t^{(m)}) \mat_1(\cE_t^{(m)})^\trans \wh{\bA}_{m,1,\perp} \wh{\bA}_{m,1,\perp}^\trans \notag \\
    &=: \wt\cII_1 + \wt\cII_2 + \wt\cII_3 + \wt\cII_4 + \wt\cII_5 + \wt\cII_6 + \wt\cII_7 + \wt\cII_8 + \wt\cII_9 .
    \label{eqn: decomp_tilde_Sigma_B}
\end{align}


By definition, the rank of $\wt\cII_1$ is $u_{m,1}$. As the nonzero eigenvalues of some matrix product $\bX\bY$ are the same as those of $\bY\bX$ (with multiplicity), from \Cref{ass: core_factor}(b) we have $\lambda_i(\wt\cII_1) \asymp \lambda_i(\bSigma_{F,m,1})$. Then similar to the proof of \Cref{lemma: eigenvalue_Sigma_B}, it remains to show the (squared) Frobenius norm of $\wt\cII_2$--$\wt\cII_9$ are $o_P(1)$.

To this end, consider $\wt\cII_2$. By the triangle inequality, we have
\iffalse %stronger but unnecessary
\begin{align}
    &\;\quad
    \|\wt\cII_2\|_F^2 = \Bigg\| \frac{1}{T p_{m,1}} \sum_{t=1}^T \wh{\bA}_{m,1,\perp} \wh{\bA}_{m,1,\perp}^\trans \bB_{m,1} \mat_1(\cF_t^{(m)}) \mat_1(\cG_t^{(m)})^\trans \bA_{m,1}^\trans \wh{\bA}_{m,1,\perp} \wh{\bA}_{m,1,\perp}^\trans \Bigg\|_F^2 \notag \\
    &\leq
    \frac{1}{T p_{m,1}^2} \big\|\bB_{m,1} \big\|_F^2 \cdot \Bigg\| \frac{1}{\sqrt{T}} \sum_{t=1}^T \mat_1(\cF_t^{(m)}) \mat_1(\cG_t^{(m)})^\trans \bA_{m,1}^\trans \big( \wh{\bA}_{m,1,\perp} - \bA_{m,1,\perp} \big) \Bigg\|_F^2 \notag \\
    &=
    \cO_P\Big( \frac{1}{T p_{m,1}} \Big) \cdot \sum_{j=1}^{p_{m,1}} \Bigg\| \frac{1}{\sqrt{T}} \sum_{t=1}^T \mat_1(\cF_t^{(m)}) \mat_1(\cG_t^{(m)})^\trans (\bA_{m,1})_{j\cdot} \Bigg\|_F^2 \cdot \big\| \wh{\bA}_{m,1,\perp} - \bA_{m,1,\perp} \big\|_F^2 \notag \\
    &= \cO_P\Big( \frac{1}{T^2} \Big)
    \label{eqn: }
\end{align}
where the second last equality used \Cref{ass: loadings}(b), and the last used \Cref{ass: dependence_thread}(b) and \Cref{lemma: A_perp_rate}.
\fi
\begin{align}
    &\;\quad
    \|\wt\cII_2\|_F^2 \notag \\
    &=
    \Bigg\| \frac{1}{T p_{m,1}} \sum_{t=1}^T \wh{\bA}_{m,1,\perp} \wh{\bA}_{m,1,\perp}^\trans \bB_{m,1} \mat_1(\cF_t^{(m)}) \mat_1(\cG_t^{(m)})^\trans \bA_{m,1}^\trans \wh{\bA}_{m,1,\perp} \wh{\bA}_{m,1,\perp}^\trans \Bigg\|_F^2 \notag \\
    &\leq
    \frac{1}{T p_{m,1}^2} \big\|\bB_{m,1} \big\|_F^2 \cdot \sum_{j=1}^{p_{m,1}} \Bigg\| \frac{1}{\sqrt{T}} \sum_{t=1}^T \mat_1(\cF_t^{(m)}) \mat_1(\cG_t^{(m)})^\trans (\bA_{m,1})_{j\cdot} \Bigg\|_F^2 
    = \cO_P\Big( \frac{1}{T} \Big)
    \label{eqn: tilde_cII_2}
\end{align}
where the last equality used Assumptions~\ref{ass: loadings}(b) and \ref{ass: dependence_thread}(b). Similarly, we have
\begin{align}
    &\;\quad
    \|\wt\cII_3\|_F^2, \;
    \|\wt\cII_5\|_F^2, \;
    \|\wt\cII_8\|_F^2
    = \cO_P\Big( \frac{1}{T} \Big) .
    \label{eqn: tilde_cII_3_5_8}
\end{align}

With \Cref{lemma: A_perp_rate}, Assumptions~\ref{ass: core_factor}(a) and \ref{ass: dependence_thread}(c), we also have 
\begin{align}
    &\;\quad
    \|\wt\cII_4\|_F^2, \;
    \|\wt\cII_6\|_F^2, \;
    \|\wt\cII_7\|_F^2
    = \cO_P\Big( \frac{1}{T} \Big) .
    \label{eqn: tilde_cII_4_6_7}
\end{align}

Lastly, using the same argument in \eqref{eqn: cII_5} (except that $p_{m,\text{-}k}=1$ here) \begin{align}
    &\;\quad
    \|\wt\cII_9\|_F^2 = \cO_P\Big( \frac{1}{T} + \frac{1}{p_{m,1}} \Big) ,
    \label{eqn: tilde_cII_9}
\end{align}
hence concluding the proof of this lemma.
\end{proof}









\subsection{Proof of theorems}


\begin{proof}[Proof of \Cref{prop: identification}]
Fix $m\in[M]$. Then we have $\cX_{G,t}^{(m)} + \cX_{F,t}^{(m)} = \wt\cX_{G,t}^{(m)} + \wt\cX_{F,t}^{(m)}$. Take any $k\in[K]$, for any $n\neq m$ and multi-index $\bi=(i_1,\ldots,i_{K_n}) \in [p_{n,1}] \times \cdots \times[p_{n,K_n}]$, we have
\begin{align*}
    &\;\quad
    \E\Big\{ \mat_k(\cX_{G,t}^{(m)}) \cdot (\cX_{G,t}^{(n)})_{\bi} \Big\}
    = \E\Big\{ \mat_k\Big( \cX_{G,t}^{(m)} +\cX_{F,t}^{(m)} \Big) \cdot \Big( \cX_{G,t}^{(n)} +\cX_{F,t}^{(n)} \Big)_{\bi} \Big\} \\
    &=
    \E\Big\{ \mat_k\Big( \wt\cX_{G,t}^{(m)} +\wt\cX_{F,t}^{(m)} \Big) \cdot \Big( \cX_{G,t}^{(n)} +\cX_{F,t}^{(n)} \Big)_{\bi} \Big\}
    = \E\Big\{ \mat_k(\wt\cX_{G,t}^{(m)}) \cdot (\cX_{G,t}^{(n)})_{\bi} \Big\} \\
    &=
    \E\Big\{ \mat_k\Big( \cX_{G,t}^{(m)} + [\wt\cX_{G,t}^{(m)} - \cX_{G,t}^{(m)}] \Big) \cdot (\cX_{G,t}^{(n)})_{\bi} \Big\} ,
\end{align*}
where the first and third equalities used Assumption~\ref{ass: identification}(b). The above implies
\[
\E\Big\{ \mat_k\Big( \wt\cX_{G,t}^{(m)} - \cX_{G,t}^{(m)} \Big) \cdot (\cX_{G,t}^{(n)})_{\bi} \Big\} = 0 ,
\]
which, together with Assumption~\ref{ass: identification}(a), suggests that any element of $\wt\cX_{G,t}^{(m)} - \cX_{G,t}^{(m)}$ is uncorrelated to any element in $\cX_{G,t}^{(n)}$, which holds true for any $n\in m$. By Assumption~\ref{ass: identification}(c), we conclude that $\wt\cX_{G,t}^{(m)} = \cX_{G,t}^{(m)}$. The global component---hence the local component---is thus identified, which completes the proof of the theorem.
\end{proof}







\begin{proof}[Proof of Theorem~\ref{thm: global_loading}]
By applying Lemma~3 in \cite{Lametal2011} and notice the matrix $P$ within is of fixed rank, we have
\begin{align*}
    \frac{1}{\sqrt{p_{m,k}}} \Big\| \wh{\bA}_{m,k} -\bA_{m,k} \Big\|_F \leq \frac{8}{\lambda_{r_{m,k}}\big( \bSigma_{m,k} \big)} \Big\| \wh\bSigma_{m,k} -\bSigma_{m,k} \Big\|_F
    = \cO_P\Big( \frac{1}{\sqrt{T}} \Big) ,
\end{align*}
where the last equality used Lemmas~\ref{lemma: rate_Sigma} and \ref{lemma: eigenvalue_Sigma}. This completes the proof of the theorem.
\end{proof}




\begin{proof}[Proof of Theorem~\ref{thm: local_loading}]
By the definition of $\wh{\bB}_{m,k}$, we have
\[
\wh{\bB}_{m,k} = \wh\bSigma_{B,m,k} \wh{\bB}_{m,k} \wh{\bD}_{B,m,k}^{-1} .
\]
First consider $K_m> 1$, and note that from the decomposition in \eqref{eqn: decomp_Sigma_B},
%we can define
recall from the statement of \Cref{thm: local_loading} that
{\small
\[
\wh{\bH}_{m,k} = \frac{1}{Tp_m} \sum_{t=1}^T \mat_k(\cF_t^{(m)}) \bB_{m,\text{-}k}^\trans \wh{\bA}_{m,\text{-}k,\perp} \wh{\bA}_{m,\text{-}k,\perp}^\trans \bB_{m,\text{-}k} \mat_k(\cF_t^{(m)})^\trans \bB_{m,k}^\trans \wh{\bB}_{m,k} \wh{\bD}_{B,m,k}^{-1} ,
\]
}
so that with the notation in \eqref{eqn: decomp_Sigma_B}, by Cauchy--Schwarz inequality, we have
\begin{align}
    &\;\quad
    \frac{1}{p_{m,k}} \Big\| \wh{\bB}_{m,k} - \bB_{m,k} \wh{\bH}_{m,k} \Big\|_F^2 = \frac{1}{p_{m,k}} \Big\| \Big( \sum_{i=2}^9 \cII_i \Big) \wh{\bB}_{m,k} \wh{\bD}_{B,m,k}^{-1} \Big\|_F^2 \notag \\
    &\lesssim
    \frac{1}{p_{m,k}} \Big\| \cII_5 \wh{\bB}_{m,k} \wh{\bD}_{B,m,k}^{-1} \Big\|_F^2 + \frac{1}{p_{m,k}} \Big\| \cII_9 \wh{\bB}_{m,k} \wh{\bD}_{B,m,k}^{-1} \Big\|_F^2 \notag \\
    &\;\quad
    + \frac{1}{p_{m,k}}  \cdot \big\| \wh{\bD}_{B,m,k}^{-1} \big\|_F^2 \cdot \sum_{\substack{i=2\\ i\neq 5}}^8 \Big\| \cII_i \wh{\bB}_{m,k} \Big\|_F^2 \notag \\
    &=
    \frac{1}{p_{m,k}} \Big\| \cII_5 \wh{\bB}_{m,k} \Big\|_F^2 + \frac{1}{p_{m,k}} \Big\| \cII_9 \wh{\bB}_{m,k} \Big\|_F^2 + \cO_P\Big( \frac{1}{T p_{m,\text{-}k}} \Big) + \cO_P\Big( \frac{1}{T p_{m,k}} \Big) \notag \\
    &\;\quad
    + o_P\Big( \frac{1}{p_{m,k}} \Big\| \wh{\bB}_{m,k} - \bB_{m,k} \wh{\bH}_{m,k} \Big\|_F^2 \Big) \notag \\
    &\lesssim
    \frac{1}{p_{m,k}} \Big\| \cII_5 \bB_{m,k} \Big\|_F^2 + \frac{1}{p_{m,k}} \Big\| \cII_9 \bB_{m,k} \Big\|_F^2 + \cO_P\Big( \frac{1}{T p_{m,\text{-}k}} \Big) + \cO_P\Big( \frac{1}{T p_{m,k}} \Big) \notag \\
    &\;\quad
    + o_P\Big( \frac{1}{p_{m,k}} \Big\| \wh{\bB}_{m,k} - \bB_{m,k} \wh{\bH}_{m,k} \Big\|_F^2 \Big) ,
    \label{eqn: hat_B_minus_BH}
\end{align}
where: (1) the second last equality used \Cref{lemma: eigenvalue_Sigma_B}, the rates in its proof, and the fact that $\big\|\bA_{m,k}^\trans \bB_{m,k} \big\|_F =\cO_P(\sqrt{p_{m,k}})$ by \Cref{ass: loadings}(b) and hence, for example, for $\cII_6 \wh{\bB}_{m,k}$, similar to \eqref{eqn: cII_6}:
\begin{align*}
    \Big\| \cII_6 \wh{\bB}_{m,k} \Big\|_F^2 &\lesssim
    \frac{1}{p_m^2} \cdot \big\|\bA_{m,k}^\trans \bB_{m,k} \big\|_F^2 \cdot \frac{p_{m,\text{-}k}}{T} \cdot \cO_P(p_{m})
    + o_P\Big( \Big\| \wh{\bB}_{m,k} - \bB_{m,k} \wh{\bH}_{m,k} \Big\|_F^2 \Big) \\
    &=
    \cO_P\Big( \frac{1}{T} \Big) + o_P\Big( \Big\| \wh{\bB}_{m,k} - \bB_{m,k} \wh{\bH}_{m,k} \Big\|_F^2 \Big) ;
\end{align*}
and (2) the last line used \eqref{eqn: cII_9}, \eqref{eqn: cII_5}, and the fact that $\big\| \wh{\bH}_{m,k} \big\|_F^2 =\cO_P(1)$ by the convergence of $\cII_1$ in the proof of \Cref{lemma: eigenvalue_Sigma_B}.




Then by \Cref{lemma: inequality_sandwich}(b) and \Cref{ass: loadings}(a), consider
\begin{align}
    &\;\quad
    \frac{1}{p_{m,k}} \Big\| \cII_5 \bB_{m,k} \Big\|_F^2 \notag \\
    &=
    \frac{1}{p_{m,k}} \Bigg\| \frac{1}{Tp_m} \sum_{t=1}^T \mat_k(\cE_t^{(m)}) \wh{\bA}_{m,\text{-}k,\perp} \wh{\bA}_{m,\text{-}k,\perp}^\trans \mat_k(\cE_t^{(m)})^\trans \bB_{m,k} \Bigg\|_F^2 \notag \\
    &\leq
    \frac{1}{T^2 p_m^2 p_{m,k}} \sum_{i=1}^{u_{m,k}} \sum_{j=1}^{p_{m,\text{-}k}} \Bigg\| \sum_{t=1}^T \big\{ \bB_{m,k}^\trans \mat_k(\cE_t^{(m)}) \big\}_{ij} \mat_k(\cE_t^{(m)})^\trans \Bigg\|_F^2 \notag \\
    &=
    \frac{1}{T^2 p_m^2 p_{m,k}} \max_{i\in[u_{m,k}]} \sum_{j=1}^{p_{m,\text{-}k}} \Bigg\| \sum_{t=1}^T (\bB_{m,k})_{\cdot i}^\trans \mat_k(\cE_t^{(m)})_{\cdot j} \mat_k(\cE_t^{(m)}) \Bigg\|_F^2 \notag \\
    &=
    \frac{1}{T^2 p_m^2 p_{m,k}} \max_{i\in[u_{m,k}]} \sum_{l=1}^{p_{m,k}} \sum_{h,j=1}^{p_{m,\text{-}k}} \Bigg( \sum_{t=1}^T \sum_{w=1}^{p_{m,k}} (\bB_{m,k})_{wi} \mat_k(\cE_t^{(m)})_{wj} \mat_k(\cE_t^{(m)})_{lh} \Bigg)^2 \notag \\
    &\lesssim
    \frac{1}{T^2 p_m^2 p_{m,k}} \max_{i\in[u_{m,k}]} \sum_{l=1}^{p_{m,k}} \sum_{h,j=1}^{p_{m,\text{-}k}} \Bigg\{ \E\Bigg( \sum_{t=1}^T \sum_{w=1}^{p_{m,k}} (\bB_{m,k})_{wi} \mat_k(\cE_t^{(m)})_{wj} \mat_k(\cE_t^{(m)})_{lh} \Bigg) \Bigg\}^2 \notag \\
    &\;\quad
    + \frac{1}{T^2 p_m^2 p_{m,k}} \max_{i\in[u_{m,k}]} \sum_{l=1}^{p_{m,k}} \sum_{h,j=1}^{p_{m,\text{-}k}} \sum_{t,s=1}^T \sum_{w,q=1}^{p_{m,k}} \cov\Big\{ \mat_k(\cE_t^{(m)})_{wj} \mat_k(\cE_t^{(m)})_{lh}, \notag \\
    &\;\quad
    \mat_k(\cE_s^{(m)})_{qj} \mat_k(\cE_s^{(m)})_{lh} \Big\}
    =
    \cO_P\Big( \frac{1}{p_{m,k}^2} + \frac{1}{T p_{m,k}} \Big) ,
    \label{eqn: cII_5_Bmk}
\end{align}
where the last equality used Assumptions~\ref{ass: noise}(b) and (c).

On the other hand, by triangle inequality, we have
\begin{align*}
    &\;\quad
    \frac{1}{p_{m,k}} \Big\| \cII_9 \bB_{m,k} \Big\|_F^2 \\
    &=
    \frac{1}{p_{m,k}} \Bigg\| \frac{1}{Tp_m} \sum_{t=1}^T \bA_{m,k} \mat_k(\cG_t^{(m)}) \\
    &\;\quad
    \cdot \bA_{m,\text{-}k}^\trans \wh{\bA}_{m,\text{-}k,\perp} \wh{\bA}_{m,\text{-}k,\perp}^\trans \bA_{m,\text{-}k} \mat_k(\cG_t^{(m)})^\trans \bA_{m,k}^\trans \bB_{m,k} \Bigg\|_F^2 \\
    &\leq
    \frac{1}{p_m^2 p_{m,k}} \big\|\bA_{m,k}\big\|_F^2 \cdot \Big\| \wh{\bA}_{m,\text{-}k,\perp} \wh{\bA}_{m,\text{-}k,\perp}^\trans - \bA_{m,\text{-}k,\perp} \bA_{m,\text{-}k,\perp}^\trans \Big\|_F^2 \\
    &\;\quad
    \cdot \Bigg\| \frac{1}{T} \sum_{t=1}^T \mat_k(\cG_t^{(m)}) \bA_{m,\text{-}k}^\trans \bA_{m,\text{-}k} \mat_k(\cG_t^{(m)})^\trans \Bigg\|_F^2 \cdot \big\| \bA_{m,k}^\trans \bB_{m,k} \big\|_F^2 \\
    &\lesssim
    \frac{1}{T} \cdot \frac{1}{p_{m,k}^2} \big\| \bA_{m,k}^\trans \bB_{m,k} \big\|_F^2
    = \frac{1}{T} \cdot \Bigg\| \frac{1}{p_{m,k}} \bA_{m,k}^\trans \bB_{m,k} \Bigg\|_F^2
    = \cO_P\Big( \frac{1}{T p_{m,k}} \Big) ,
\end{align*}
where the last line used Assumption~\ref{ass: loadings}(b), \Cref{lemma: A_perp_rate}, and similar arguments in \eqref{eqn: cII_9}.


Finally, noting that the last term in \eqref{eqn: hat_B_minus_BH} can be omitted, we conclude that
\[
\frac{1}{p_{m,k}} \Big\| \wh{\bB}_{m,k} - \bB_{m,k} \wh{\bH}_{m,k} \Big\|_F^2 = \cO_P\Big( \frac{1}{p_{m,k}^2} + \frac{1}{T p_{m,k}} + \frac{1}{T p_{m,\text{-}k}} \Big) .
\]

It remains to show $\wh{\bH}_{m,k} \wh{\bH}_{m,k}^\trans \xrightarrow{P} \bI$, which is obvious by
\begin{align*}
    \bI &= \frac{1}{p_{m,k}} \wh{\bB}_{m,k}^\trans \wh{\bB}_{m,k} = \frac{1}{p_{m,k}} \wh{\bB}_{m,k}^\trans (\wh{\bB}_{m,k} - \bB_{m,k} \wh{\bH}_{m,k}) + \frac{1}{p_{m,k}} \wh{\bB}_{m,k}^\trans \bB_{m,k} \wh{\bH}_{m,k} \\
    &=
    \frac{1}{p_{m,k}} \wh{\bB}_{m,k}^\trans (\wh{\bB}_{m,k} - \bB_{m,k} \wh{\bH}_{m,k}) + \frac{1}{p_{m,k}} (\wh{\bB}_{m,k} - \bB_{m,k} \wh{\bH}_{m,k})^\trans \bB_{m,k} \wh{\bH}_{m,k} \\
    &\;\quad
    + \frac{1}{p_{m,k}} \wh{\bH}_{m,k}^\trans \bB_{m,k}^\trans \bB_{m,k} \wh{\bH}_{m,k}
    =
    \wh{\bH}_{m,k} \wh{\bH}_{m,k}^\trans + o_P(1) ,
\end{align*}
which completes part~(a) of the theorem.


Consider part~(b). Similar to part~(a), we also have $\wh{\bB}_{m,1} = \wt\bSigma_{B,m,1} \wh{\bB}_{m,1} \wt{\bD}_{B,m,1}^{-1}$. For $\wt\cII_1$ in \eqref{eqn: decomp_tilde_Sigma_B}, by \Cref{lemma: A_perp_rate} and \eqref{eqn: B_A_perp_A_perp_B},
\begin{align*}
    \wt\cII_1 &= \frac{1}{Tp_{m,1}} \sum_{t=1}^T \bA_{m,1,\perp} \bA_{m,1,\perp}^\trans \bB_{m,1} \mat_1(\cF_t^{(m)}) \mat_1(\cF_t^{(m)})^\trans \bB_{m,1}^\trans \wh{\bA}_{m,1,\perp} \wh{\bA}_{m,1,\perp}^\trans \\
    &\;\quad
    + o_P(\wt\cII_1) \\
    &\asymp
    \sqrt{p_{m,1}} \cdot \big( p_{m,1}^{-1/2} \bB_{m,1}, \bC_{m,1} \big) \big( \bI, \zero\big)^\trans \\
    &\;\quad
    \cdot \frac{1}{Tp_{m,1}} \sum_{t=1}^T \mat_1(\cF_t^{(m)}) \mat_1(\cF_t^{(m)})^\trans \bB_{m,1}^\trans \wh{\bA}_{m,1,\perp} \wh{\bA}_{m,1,\perp}^\trans \\
    &\asymp
    \frac{1}{Tp_{m,1}} \sum_{t=1}^T \bB_{m,1} \mat_1(\cF_t^{(m)}) \mat_1(\cF_t^{(m)})^\trans \bB_{m,1}^\trans \wh{\bA}_{m,1,\perp} \wh{\bA}_{m,1,\perp}^\trans ,
\end{align*}
so that the remaining of part~(b) can be shown by repeating the steps for part~(a), except that \Cref{lemma: eigenvalue_tilde_Sigma_B} is used in place of \Cref{lemma: eigenvalue_Sigma_B}, and for \eqref{eqn: cII_5_Bmk}, Lemma~4 in \cite{barigozzi2022statistical} can be directly applied. The proof of this theorem is now complete.
\end{proof}








\begin{proof}[Proof of Theorem~\ref{thm: core_factor}]
For the local factor estimator, consider $K_m> 1$ first. From \eqref{eqn: def_hat_Ft}, we have
\begin{align}
    &\;\quad
    \vec(\wh{\cF}_t^{(m)}) = (\wh{\bC}_m^\trans \wh{\bC}_m)^{-1} \wh{\bC}_m^\trans (\otimes_{k=1}^K \wh{\bA}_{m,k,\perp})^\trans \vec(\cX_t^{m}) \notag \\
    &=
    (\wh{\bC}_m^\trans \wh{\bC}_m)^{-1} \wh{\bC}_m^\trans \Big\{ (\otimes_{k=1}^K \wh{\bA}_{m,k,\perp})^\trans (\otimes_{k=1}^K \bB_{m,k}) \vec(\cF_t^{(m)}) \notag \\
    &\;\quad
    + (\otimes_{k=1}^K \wh{\bA}_{m,k,\perp})^\trans (\otimes_{k=1}^K \bA_{m,k}) \vec(\cG_t^{(m)}) + (\otimes_{k=1}^K \wh{\bA}_{m,k,\perp})^\trans \vec(\cE_t^{m}) \Big\} \notag \\
    &=
    (\otimes_{k=1}^K \wh{\bH}_{m,k}^\trans) \vec(\cF_t^{(m)}) \notag \\
    &\;\quad
    + (\wh{\bC}_m^\trans \wh{\bC}_m)^{-1} \wh{\bC}_m^\trans (\otimes_{k=1}^K \wh{\bA}_{m,k,\perp})^\trans \big( \otimes_{k=1}^K \bB_{m,k} - \otimes_{k=1}^K \wh{\bB}_{m,k} \wh{\bH}_{m,k}^\trans \big) \vec(\cF_t^{(m)}) \notag \\
    &\;\quad
    + (\wh{\bC}_m^\trans \wh{\bC}_m)^{-1} \wh{\bC}_m^\trans (\otimes_{k=1}^K \wh{\bA}_{m,k,\perp})^\trans (\otimes_{k=1}^K \bA_{m,k}) \vec(\cG_t^{(m)}) \notag \\
    &\;\quad
    + (\wh{\bC}_m^\trans \wh{\bC}_m)^{-1} \wh{\bC}_m^\trans (\otimes_{k=1}^K \wh{\bA}_{m,k,\perp})^\trans \vec(\cE_t^{m}) \notag \\
    &=: (\otimes_{k=1}^K \wh{\bH}_{m,k}^\trans) \vec(\cF_t^{(m)}) + \cIII_1 + \cIII_2 + \cIII_3 .
    \label{eqn: vec_F_decomp}
\end{align}


Before considering the last line in \eqref{eqn: vec_F_decomp}, first note that
\begin{align}
    &\;\quad
    \Big\| (\wh{\bC}_m^\trans \wh{\bC}_m)^{-1} \wh{\bC}_m^\trans \Big\|_F^2 \notag \\
    &\leq
    \Big\| \big( \otimes_{k=1}^K \wh{\bB}_{m,k}^\trans \wh{\bA}_{m,k,\perp} \wh{\bA}_{m,k,\perp}^\trans \wh{\bB}_{m,k} \big)^{-1} \Big\|_F^2 \cdot \big\| \otimes_{k=1}^K \wh{\bA}_{m,k,\perp} \big\|_F^2 \cdot \big\| \otimes_{k=1}^K \wh{\bB}_{m,k} \big\|_F^2 \notag \\
    &\asymp
    \big\| (p_m \bI)^{-1} \big\|_F^2 \cdot \big\| \otimes_{k=1}^K \wh{\bA}_{m,k,\perp} \big\|_F^2 \cdot \big\| \otimes_{k=1}^K \wh{\bB}_{m,k} \big\|_F^2
    = \cO_P\big( p_m^{-1} \big) ,
    \label{eqn: CCC_rate}
\end{align}
where the last line used \Cref{lemma: A_perp_rate}, \eqref{eqn: B_A_perp_A_perp_B}, and \Cref{ass: loadings}(b).

With \eqref{eqn: CCC_rate}, by Assumption~\ref{ass: core_factor}(a) and \Cref{thm: local_loading}, we immediately have (by a simple induction):
\begin{equation}
\label{eqn: cIII_1}
\begin{split}
    \| \cIII_1\|_F^2 &\lesssim \sum_{k\in[K_m]} \frac{1}{p_{m,k}} \Big\| \bB_{m,k} - \wh{\bB}_{m,k} \wh{\bH}_{m,k}^\trans \Big\|_F^2 \\
    &=
    \cO_P\Bigg\{ \sum_{k=1}^{K_m} \Big( \frac{1}{p_{m,k}^2} + \frac{1}{T p_{m,k}} + \frac{1}{T p_{m,\text{-}k}} \Big) \Bigg\} .
\end{split}
\end{equation}

For $\cIII_2$, using \eqref{eqn: CCC_rate} and \Cref{lemma: A_perp_rate},
\begin{equation}
\label{eqn: cIII_2}
\| \cIII_2\|_F^2 \lesssim \frac{1}{p_m} \big\| \otimes_{k=1}^K \bA_{m,k} \big\|_F^2 \cdot \Big\| \wh{\bA}_{m,k,\perp} -\bA_{m,k,\perp} \Big\|_F^2 
= \cO_P\Big( \frac{1}{T} \Big).
\end{equation}

For $\cIII_3$, first consider $\wh{\bC}_m^\trans (\otimes_{k=1}^K \wh{\bA}_{m,k,\perp})^\trans \vec(\cE_t^{m})$. By \eqref{eqn: B_A_perp_A_perp_B}, we have
\begin{align*}
    &\;\quad
    \Big\| \wh{\bC}_m^\trans (\otimes_{k=1}^K \wh{\bA}_{m,k,\perp})^\trans \vec(\cE_t^{m}) \Big\|_F^2 \\
    &\lesssim \max_i \big\| (\otimes_{k=1}^K \bB_{m,k})_{i\cdot} \big\|_F^2 \sum_{j,l=1}^{p_m} \big| \E\big\{ \vec(\cE_t^{m})_j \vec(\cE_t^{m})_l \big\} \Big|
    = \cO_P(p_m) ,
\end{align*}
where the last equality used Assumptions~\ref{ass: loadings}(a) and \ref{ass: noise}(b). Noting that $\big\| (\wh{\bC}_m^\trans \wh{\bC}_m)^{-1} \big\|_F^2 =\cO_P\big( p_m^{-2} \big)$ from \eqref{eqn: CCC_rate}, it hence holds that
\begin{equation}
\label{eqn: cIII_3}
\| \cIII_3\|_F^2 = \cO_P\Big( \frac{1}{p_m} \Big).
\end{equation}

Therefore, combining \eqref{eqn: vec_F_decomp}, \eqref{eqn: cIII_1}, \eqref{eqn: cIII_2}, and \eqref{eqn: cIII_3}, we conclude
\[
\Big\| \vec(\wh{\cF}_t^{(m)}) - (\otimes_{k=1}^K \wh{\bH}_{m,k}^\trans) \vec(\cF_t^{(m)}) \Big\|_F^2 = \cO_P\Big( \sum_{k\in[K_m]} \frac{1}{p_{m,k}^2} + \frac{1}{T} + \frac{1}{p_m} \Big) ,
\]
which completes the proof for the scenario $K_m> 1$. When $K_m=1$, all the previous arguments follow, except that for \eqref{eqn: cIII_1}, we have
\[
\| \cIII_1\|_F^2 = \cO_P\Big( \frac{1}{p_{m,1}^2} + \frac{1}{T} \Big) ,
\]
which leads to essentially the same expression (note that $p_m=p_{m,1}$ for $K_m=1$, and $\wh{\bH}_{m,k}$ is replaced by $\wt{\bH}_{m,1}$), i.e.,
\[
\Big\| \vec(\wh{\cF}_t^{(m)}) - \wt{\bH}_{m,1}^\trans \vec(\cF_t^{(m)}) \Big\|_F^2 = \cO_P\Big( \frac{1}{T} + \frac{1}{p_{m,1}} \Big) .
\]
This completes the proof for part~(a) of the theorem.



Next, we first show the consistency of the vectorized local component as follows. For $K_m> 1$,
\begin{align}
    &\;\quad
    \frac{1}{p_m} \Big\| (\otimes_{k=1}^K \bB_{m,k}) \vec(\cF_t^{(m)}) - (\otimes_{k=1}^K \wh{\bB}_{m,k}) \vec(\wh\cF_t^{(m)}) \Big\|_F^2 \notag \\
    &=
    \frac{1}{p_m} \Big\| (\otimes_{k=1}^K \bB_{m,k} \wh{\bH}_{m,k}) (\otimes \wh{\bH}_{m,k}^\trans) \vec(\cF_t^{(m)}) - (\otimes_{k=1}^K \wh{\bB}_{m,k}) \vec(\wh\cF_t^{(m)}) \Big\|_F^2 \notag \\
    &\lesssim
    \frac{1}{p_m} \Big\| (\otimes_{k=1}^K \bB_{m,k} \wh{\bH}_{m,k}) - (\otimes_{k=1}^K \wh{\bB}_{m,k}) \Big\|_F^2
    + \Big\| (\otimes \wh{\bH}_{m,k}^\trans) \vec(\cF_t^{(m)}) - \vec(\wh\cF_t^{(m)}) \Big\|_F^2 \notag \\
    &=
    \cO_P\Big( \sum_{k\in[K_m]} \frac{1}{p_{m,k}^2} + \frac{1}{T} + \frac{1}{p_m} \Big) ,
    \label{eqn: rate_local_component_Km_not1}
\end{align}
by Theorems~\ref{thm: local_loading} and \ref{thm: core_factor}(a). Similarly, for $K_m=1$,
\begin{align}
    &\;\quad
    \frac{1}{p_m} \Big\| (\otimes_{k=1}^K \bB_{m,k}) \vec(\cF_t^{(m)}) - (\otimes_{k=1}^K \wh{\bB}_{m,k}) \vec(\wh\cF_t^{(m)}) \Big\|_F^2 \notag \\
    &=
    \cO_P\Big( \frac{1}{T} + \frac{1}{p_{m,1}} \Big) .
    \label{eqn: rate_local_component_Km_1}
\end{align}


Then for part~(b), by the definition in \eqref{eqn: def_hat_Gt}, we may write
\begin{align*}
    &\;\quad
    \vec(\wh{\cG}_t^{(m)}) \\
    &=
    \Big\{ (\otimes_{k=1}^K \wh{\bA}_{m,k})^\trans (\otimes_{k=1}^K \wh{\bA}_{m,k}) \Big\}^{-1} (\otimes_{k=1}^K \wh{\bA}_{m,k})^\trans \Big\{ \vec(\cX_t^{m}) - (\otimes_{k=1}^K \wh{\bB}_{m,k}) \vec(\wh\cF_t^{(m)}) \Big\} \\
    &=
    \Big\{ (\otimes_{k=1}^K \wh{\bA}_{m,k})^\trans (\otimes_{k=1}^K \wh{\bA}_{m,k}) \Big\}^{-1} (\otimes_{k=1}^K \wh{\bA}_{m,k})^\trans \Big\{ \vec(\cX_t^{m}) - (\otimes_{k=1}^K \bB_{m,k}) \vec(\cF_t^{(m)}) \Big\} \\
    &\;\quad
    + \Big\{ (\otimes_{k=1}^K \wh{\bA}_{m,k})^\trans (\otimes_{k=1}^K \wh{\bA}_{m,k}) \Big\}^{-1} (\otimes_{k=1}^K \wh{\bA}_{m,k})^\trans \\
    &\;\quad
    \cdot \Big\{ (\otimes_{k=1}^K \bB_{m,k}) \vec(\cF_t^{(m)}) - (\otimes_{k=1}^K \wh{\bB}_{m,k}) \vec(\wh\cF_t^{(m)}) \Big\} \\
    &=
    \Big\{ (\otimes_{k=1}^K \wh{\bA}_{m,k})^\trans (\otimes_{k=1}^K \wh{\bA}_{m,k}) \Big\}^{-1} (\otimes_{k=1}^K \wh{\bA}_{m,k})^\trans \Big\{ (\otimes_{k=1}^K \bA_{m,k}) \vec(\cG_t^{(m)}) + \vec(\cE_t^{m}) \Big\} \\
    &\;\quad
    + \Big\{ (\otimes_{k=1}^K \wh{\bA}_{m,k})^\trans (\otimes_{k=1}^K \wh{\bA}_{m,k}) \Big\}^{-1} (\otimes_{k=1}^K \wh{\bA}_{m,k})^\trans \\
    &\;\quad
    \cdot \Big\{ (\otimes_{k=1}^K \bB_{m,k}) \vec(\cF_t^{(m)}) - (\otimes_{k=1}^K \wh{\bB}_{m,k}) \vec(\wh\cF_t^{(m)}) \Big\} \\
    &=
    \vec(\cG_t^{(m)}) + \Big\{ (\otimes_{k=1}^K \wh{\bA}_{m,k})^\trans (\otimes_{k=1}^K \wh{\bA}_{m,k}) \Big\}^{-1} \\
    &\;\quad
    \cdot (\otimes_{k=1}^K \wh{\bA}_{m,k})^\trans \Big\{ (\otimes_{k=1}^K \bA_{m,k}) -(\otimes_{k=1}^K \wh{\bA}_{m,k}) \Big\} \vec(\cG_t^{(m)}) \\
    &\;\quad
    + \Big\{ (\otimes_{k=1}^K \wh{\bA}_{m,k})^\trans (\otimes_{k=1}^K \wh{\bA}_{m,k}) \Big\}^{-1} (\otimes_{k=1}^K \wh{\bA}_{m,k})^\trans \\
    &\;\quad
    \cdot \Big\{ (\otimes_{k=1}^K \bB_{m,k}) \vec(\cF_t^{(m)}) - (\otimes_{k=1}^K \wh{\bB}_{m,k}) \vec(\wh\cF_t^{(m)}) \Big\} \\
    &\;\quad
    + \Big\{ (\otimes_{k=1}^K \wh{\bA}_{m,k})^\trans (\otimes_{k=1}^K \wh{\bA}_{m,k}) \Big\}^{-1} (\otimes_{k=1}^K \wh{\bA}_{m,k})^\trans \vec(\cE_t^{m}) \\
    &=: \vec(\cG_t^{(m)}) + \cIV_1 + \cIV_2 + \cIV_3 .
\end{align*}

The terms $\cIV_1$--$\cIV_3$ are comparable to $\cIII_1$--$\cIII_3$ in \eqref{eqn: vec_F_decomp}. Indeed, combining \Cref{thm: global_loading}, \eqref{eqn: rate_local_component_Km_not1}, \eqref{eqn: rate_local_component_Km_1}, and analogous argument as for $\cIII_3$ (see details above \eqref{eqn: cIII_3}), it holds similarly as the previous proof for part~(a) that
\begin{align*}
    \|\cIV_1\|_F^2 &\lesssim \sum_{k=1}^{K_m} \frac{1}{p_{m,k}} \Big\| \wh{\bA}_{m,k} -\bA_{m,k} \Big\|_F^2 = \cO_P\Big( \frac{1}{T} \Big) ; \\
    \|\cIV_2\|_F^2 &\lesssim \frac{1}{p_m} \Big\| (\otimes_{k=1}^K \bB_{m,k}) \vec(\cF_t^{(m)}) - (\otimes_{k=1}^K \wh{\bB}_{m,k}) \vec(\wh\cF_t^{(m)}) \Big\|_F^2 \\
    &= \begin{cases}
        \cO_P\Big( \sum_{k\in[K_m]} \frac{1}{p_{m,k}^2} + \frac{1}{T} + \frac{1}{p_m} \Big) & \text{if $K_m> 1$,} \\
        \cO_P\Big( \frac{1}{T} + \frac{1}{p_{m,1}} \Big) & \text{if $K_m= 1$;}
    \end{cases} \\
    \|\cIV_3\|_F^2 &= \cO_P\Big( \frac{1}{p_m} \Big) .
\end{align*}

Putting all together, we hence conclude
\[
\Big\| \vec(\wh{\cG}_t^{(m)}) - \vec(\cG_t^{(m)}) \Big\|_F^2 =
\begin{cases}
    \cO_P\Big( \sum_{k\in[K_m]} \frac{1}{p_{m,k}^2} + \frac{1}{T} + \frac{1}{p_m} \Big) & \text{if $K_m> 1$,} \\
    \cO_P\Big( \frac{1}{T} + \frac{1}{p_{m,1}} \Big) & \text{if $K_m= 1$.}
\end{cases}
\]

The rate of convergence for the local component estimator has been spelled out in \eqref{eqn: rate_local_component_Km_not1} and \eqref{eqn: rate_local_component_Km_1}. The arguments for the global component estimator follow similarly and hence omitted here. This ends the proof of the theorem.
\end{proof}







\begin{proof}[Proof of \Cref{prop: global_loading_bootstrap}]
The proof is direct by following Lemmas~\ref{lemma: rate_Omega}--\ref{lemma: eigenvalue_Sigma} similarly (with the rate $p_n$ being $|\cS_n^\dagger|$ now) and hence \Cref{thm: global_loading}, except that in \eqref{eqn: eigenvalue_Sigma_step_2} in the proof of \Cref{lemma: eigenvalue_Sigma}:
\begin{align*}
    &\;\quad
    \lambda_{r_{m,k}}\Bigg( \sum_{\bi \in \cS_n} \E\Big\{ \mat_k(\cG_t^{(m)}) \Big( \cG_t^{(n)} \times_{h=1}^{K_n} \bA_{n,h} \Big)_{\bi} \Big\} \\
    &\;\quad
    \cdot \E\Big\{ \mat_k(\cG_t^{(m)}) \Big( \cG_t^{(n)} \times_{h=1}^{K_n} \bA_{n,h} \Big)_{\bi} \Big\}^\trans \Bigg) \notag \\
    &=
    \lambda_{r_{m,k}}\Bigg( \E\Big[ \mat_k(\cG_t^{(m)}) \otimes \Big\{ \vec(\cG_t^{(n)})^\trans \big( \otimes_{h=1}^{K_n} \bA_{n,h} \big)_{\cS_n^\dagger}^\trans \Big\} \Big] \notag \\
    &\;\quad
    \cdot \E\Big[ \mat_k(\cG_t^{(m)})^\trans \otimes \Big\{ \big( \otimes_{h=1}^{K_n} \bA_{n,h} \big)_{\cS_n^\dagger} \vec(\cG_t^{(n)}) \Big\} \Big] \Bigg) 
    %& \text{\ZC{it seems to work here, directly assume the above}}
    \asymp
    \cS_n^\dagger \cdot c ,
\end{align*}
where the last line used \Cref{ass: core_factor}(c) and the additional condition in the statement of \Cref{prop: global_loading_bootstrap}. This concludes the proof of the theorem.
\end{proof}





\begin{proof}[Proof of \Cref{thm: eigenratio}]
Fix any $m\in[M]$ and $k\in[K_m]$. By Weyl's inequality and \Cref{lemma: rate_Sigma}, we have, for any $j\in[p_{m,k}]$,
%for any $j\in[r_{m,k}]$,
\begin{align*}
    \Big| \lambda_j(\wh\bSigma_{m,k}) -\lambda_j(\bSigma_{m,k}) \Big| \leq \Big\| \wh\bSigma_{m,k} -\bSigma_{m,k} \Big\|
    = \cO_P\Bigg\{ \frac{p_m}{\sqrt{T}} \Bigg( \sum_{\substack{n=1\\ n\neq m}}^M p_n \Bigg) \Bigg\}
    = \cO_P(\xi_{m,k}) ,
\end{align*}
where the last equality used the definition of $\xi_{m,k}$.

Consider the scenario of $r_{m,k}>1$. For any $j\in[r_{m,k} -1]$, we hence have
\begin{equation}
\label{eqn: eigenratio_proof_1}
\frac{\lambda_{j+1}(\wh\bSigma_{m,k}) +\xi_{m,k}}{\lambda_j(\wh\bSigma_{m,k}) +\xi_{m,k}} 
\geq \frac{\lambda_{j+1}(\bSigma_{m,k}) +\cO_P(\xi_{m,k})}{\lambda_j(\bSigma_{m,k}) +\cO_P(\xi_{m,k})}
\asymp c ,
\end{equation}
for some positive constant $c$, where the asymptotic (in probability) equality used \Cref{lemma: eigenvalue_Sigma} and the definition of $\xi_{m,k}$. The above is unnecessary to prove this theorem when $r_{m,k}=1$.

Next, note that the rank of $\bSigma_{m,k}$ is upper bounded by $r_{m,k}$ due to the definition of $\bSigma_{m,k}$ where $\bA_{m,k}\in \R^{p_{m,k}\times r_{m,k}}$ within. We then have
\begin{equation}
\label{eqn: eigenratio_proof_2}
\begin{split}
    &\;\quad
    \frac{\lambda_{r_{m,k}+1}(\wh\bSigma_{m,k}) +\xi_{m,k}}{\lambda_{r_{m,k}}(\wh\bSigma_{m,k}) +\xi_{m,k}} \leq \frac{\lambda_{r_{m,k}+1}(\bSigma_{m,k}) +\cO_P(\xi_{m,k})}{\lambda_{r_{m,k}}(\bSigma_{m,k}) +\cO_P(\xi_{m,k})} \\
    &=
    \frac{\cO_P(\xi_{m,k})}{\lambda_{r_{m,k}}(\bSigma_{m,k}) +\cO_P(\xi_{m,k})}
    = \cO_P\Big( \frac{1}{\sqrt{T}} \Big) ,
\end{split}
\end{equation}
where the last equality used again \Cref{lemma: eigenvalue_Sigma} and the definition of $\xi_{m,k}$.

It remains to consider the eigenvalue ratio for $j\in \{r_{m,k}+1, \dots, r_{\max,m,k}\}$. In detail, we have
\begin{equation}
\label{eqn: eigenratio_proof_3}
\frac{\lambda_{j+1}(\wh\bSigma_{m,k}) +\xi_{m,k}}{\lambda_j(\wh\bSigma_{m,k}) +\xi_{m,k}}
\geq \frac{\xi_{m,k}}{\lambda_j(\wh\bSigma_{m,k}) +\xi_{m,k}}
\geq \frac{\xi_{m,k}}{\cO_P(\xi_{m,k}) +\xi_{m,k}}
\geq 1/c ,
\end{equation}
in probability for some positive constant $c$, where we also used the definition of $\xi_{m,k}$. Lastly, combining \eqref{eqn: eigenratio_proof_1}, \eqref{eqn: eigenratio_proof_2}, and \eqref{eqn: eigenratio_proof_3}, we conclude the proof of the theorem.
\end{proof}








\subsection{Proof of claims}

\begin{proof}[Proof of \Cref{claim: assump_fac_simpler}]
Note that elements of the global factors have zero mean according to \Cref{ass: identification}(a). To prove the claim, it suffices to consider the case when global factors have elements with unit variance, since the variance can be absorbed into the global loading matrices otherwise. For simplicity, we use throughout this proof the following notation.
\begin{align*}
    & \cG_t := \cG_t^{(1)} =\dots =\cG_t^{(M)}, \\
    & K:= K_1=\dots =K_M, \\
    % & \bA_{h} := \bA_{n,h} \quad \text{for each $h\in[K]$}, \\
    & r_{h} := r_{1,h}=\dots =r_{M,h} \quad \text{for each $h\in[K]$} , \\
    & r_{\text{-}k} := r_1 \dots r_{k-1} r_{k+1} \dots r_K .
\end{align*}
where we fix $k\in[K]$ and $n\in[K] \setminus\{m\}$.

Further define
\begin{align*}
    \bM_t &:= \mat_k(\cG_t^{(m)}) \otimes \Big\{ \vec(\cG_t^{(n)})^\trans \big( \otimes_{h=1}^{K_n} \bA_{h} \big)^\trans \Big\} \\
    &=
    \mat_k(\cG_t) \otimes \Big\{ \vec(\cG_t)^\trans \big( \otimes_{h=1}^{K} \bA_{h} \big)^\trans \Big\} ,
\end{align*}
which is a matrix of dimension $r_k \times (r_{\text{-}k} p_n)$. We next write the expression in index form.
Let the entries of $\cG_t$ be denoted by $g_{t,i_1, \dots, i_K}$ with $i_h \in [r_h]$.
Denote the entries of mode-$k$ unfolding of $\cG_t$ as
\[
\mat_k(\cG_t)_{i_k, \balpha_k} = g_{t,i_1, \dots, i_K},
\]
where $\balpha_k$ corresponds to the multi-index $(i_1, \dots, i_{k-1}, i_{k+1}, \dots, i_K)$. For the row vector $\vec(\cG_t^{(n)})^\trans \big( \otimes_{h=1}^{K_n} \bA_{n,h} \big)^\trans$ which has length $p_n=p_{n,1}\dots p_{n,K}$, let $\bbeta = (j_1, \dots, j_K)$ denote its multi-index (corresponding to each $\bA_{n,h}$). Then we can write
\[
\Big\{ \vec(\cG_t^{(n)})^\trans \big( \otimes_{h=1}^{K_n} \bA_{n,h} \big)^\trans \Big\}_{\bbeta} = \sum_{i'_1, \dots, i'_K} \left( \prod_{h=1}^K (\bA_h)_{j_h, i'_h} \right) g_{t,i'_1, \dots, i'_K}.
\]

Hence, the $(i_k, (\balpha_k, \bbeta))$ entry of $\bM_t$ is
\[
\bM_t(i_k, (\balpha_k, \bbeta)) = g_{t,i_1, \dots, i_K} \cdot \left\{ \sum_{i'_1, \dots, i'_K} \left( \prod_{h=1}^K (\bA_h)_{j_h, i'_h} \right) g_{t,i'_1, \dots, i'_K} \right\},
\]
where $i_k$ determines the row index and $\balpha_k$ determines the other indices.


Since the entries of $\cG_t$ are uncorrelated with mean 0 and variance 1, we have
\[
\E(g_{t,i_1, \dots, i_K} \cdot g_{t,i'_1, \dots, i'_K}) = 
\begin{cases}
1 & \text{if } (i_1, \dots, i_K) = (i'_1, \dots, i'_K), \\
0 & \text{otherwise}.
\end{cases}
\]
Therefore,
\begin{equation}
\label{eqn: Mt_multi_index_simplify}
\E[\bM_t(i_k, (\balpha_k, \bbeta))] = \prod_{h=1}^K (\bA_h)_{j_h, i_h}
= (\bA_1)_{j_1, i_1} (\bA_2)_{j_2, i_2} \dots (\bA_K)_{j_K, i_K} .
\end{equation}

In what follows, we recognise a Kronecker product structure in $\bM_t$. To this end, rewrite the product in \eqref{eqn: Mt_multi_index_simplify} as
\begin{equation}
\label{eqn: Mt_multi_index_simplify_2}
\E[\bM_t(i_k, (\balpha_k, \bbeta))] = (\bA_k)_{j_k, i_k} \cdot \left\{ \prod_{h \neq k} (\bA_h)_{j_h, i_h} \right\}.
\end{equation}


For each $h \neq k$, denote the $i_h$th column of $\bA_h$ as $\bA_{h,\cdot i_h}$. Then the product over $h \neq k$ in \eqref{eqn: Mt_multi_index_simplify_2} is the $(j_1, \dots, j_{k-1}, j_{k+1}, \dots, j_K)$ entry of
\[
\bA_{1,\cdot i_1} \otimes \bA_{2,\cdot i_2} \otimes \dots \otimes \bA_{k-1,\cdot i_{k-1}} \otimes \bA_{k+1,\cdot i_{k+1}} \otimes \dots \otimes \bA_{K,\cdot i_K} .
\]

Define $\bA_{\text{-}k} := \otimes_{h\in[K] \setminus \{k\}} \bA_h$. Then the element of $\bA_{\text{-}k}$ corresponding to row (index by multi-index) $(j_1, \dots, j_{k-1}, j_{k+1}, \dots, j_K)$ and column $\balpha_k$ is 
\[
\bA_{\text{-}k}(\bbeta, \balpha_k) = \prod_{h \neq k} (\bA_h)_{j_h, i_h} .
\]
Thus by \eqref{eqn: Mt_multi_index_simplify_2}, we have
\begin{align*}
    \E[\bM_t(i_k, (\balpha_k, \bbeta))] &= (\bA_k)_{j_k, i_k} \cdot \bA_{\text{-}k}(\bbeta, \balpha_k)
    = (\bA_k^\trans)_{i_k, j_k} \cdot \bA_{\text{-}k}(\bbeta, \balpha_k) ,
\end{align*}
which implies $\E[\bM_t(:, (\balpha_k, \bbeta))] = (\bA_k^\trans)_{\cdot j_k} \cdot \bA_{\text{-}k}(\bbeta, \balpha_k)$ and hence
\begin{align*}
    \E[\bM_t(:, (\balpha_k, :))] &= \bA_k^\trans \otimes \{ \bA_{\text{-}k}(:, \balpha_k) \}^\trans 
    = \bA_k^\trans \otimes (\bA_{\text{-}k})_{\cdot \balpha_k}^\trans .
\end{align*}

Finally, $\E[\bM_t]$ is the row block vector by concatenating all $\E[\bM_t(:, (\balpha_k, :))]$ in order. Thus, we conclude
\begin{align*}
    \E[\bM_t] = \bA_k^\trans \otimes \vec(\bA_{\text{-}k})^\trans
    = \bA_k^\trans \otimes \vec\big( \otimes_{h\in[K] \setminus \{k\}} \bA_h \big)^\trans .
\end{align*}


Then it is easy to see that for any $j\in[r_k]$,
\begin{align*}
    &\;\quad
    \sigma_{j}^2\big( \E[\bM_t] \big) = \lambda_{j}\big( \E[\bM_t] \E[\bM_t]^\trans \big) \\
    &=
    \lambda_{j}\Big\{ \bA_k^\trans \bA_k \otimes \vec\big( \otimes_{h\in[K] \setminus \{k\}} \bA_h \big)^\trans \vec\big( \otimes_{h\in[K] \setminus \{k\}} \bA_h \big) \Big\}
    \asymp c ,
\end{align*}
for some positive constant $c$, where the last line used the fact that $\bA_k^\trans \bA_k$ has all first $r_k$ eigenvalues bounded away from zero and infinity, and
\begin{align*}
    &\;\quad
    \vec\big( \otimes_{h\in[K] \setminus \{k\}} \bA_h \big)^\trans \vec\big( \otimes_{h\in[K] \setminus \{k\}} \bA_h \big) \Big\} \\
    &=
    \Big\| \vec\big( \otimes_{h\in[K] \setminus \{k\}} \bA_h \big) \Big\|^2
    = \Big\| \otimes_{h\in[K] \setminus \{k\}} \bA_h \Big\|_F^2 \asymp c .
\end{align*}
The proof of the claim is thus complete.
\end{proof}



\end{document}
